\documentclass[11pt]{article}
\usepackage[margin=1.1in]{geometry}
\usepackage{amsmath,amssymb,amsthm,mathtools,booktabs,array,tabularx,longtable}
\usepackage[dvipsnames]{xcolor}
\usepackage[colorlinks=true,linkcolor=MidnightBlue,citecolor=MidnightBlue,urlcolor=MidnightBlue]{hyperref}
\usepackage{fancyhdr}
\pagestyle{fancy}\fancyhf{}\renewcommand{\headrulewidth}{0pt}
\fancyfoot[R]{\thepage}
\fancyfoot[L]{\scriptsize\color{gray}Generated by the Transport Lab, not the version of record. Official version: \href{https://doi.org/10.5281/zenodo.21743863}{doi:10.5281/zenodo.21743863}}
\fancypagestyle{plain}{\fancyhf{}\fancyfoot[R]{\thepage}\fancyfoot[L]{\scriptsize\color{gray}Generated by the Transport Lab, not the version of record. Official version: \href{https://doi.org/10.5281/zenodo.21743863}{doi:10.5281/zenodo.21743863}}\renewcommand{\headrulewidth}{0pt}}
\newtheorem{theorem}{Theorem}[section]
\newtheorem{proposition}[theorem]{Proposition}
\newtheorem{lemma}[theorem]{Lemma}
\newtheorem{corollary}[theorem]{Corollary}
\newtheorem{conjecture}[theorem]{Conjecture}
\newtheorem{refutation}[theorem]{Refutation}
\newtheorem{prediction}[theorem]{Prediction}
\newtheorem{principle}[theorem]{Principle}
\newtheorem{observation}[theorem]{Observation}
\newtheorem{axiom}[theorem]{Axiom}
\theoremstyle{definition}
\newtheorem{definition}[theorem]{Definition}
\newtheorem{fact}[theorem]{Fact}
\newtheorem{assembled}[theorem]{Assembly}
\newtheorem{postulate}[theorem]{Postulate}
\newtheorem{premise}[theorem]{Premise}
\newtheorem{hypothesis}[theorem]{Hypothesis}
\newtheorem{remark}[theorem]{Remark}
\newtheorem{example}[theorem]{Example}
\theoremstyle{plain}
\setlength{\marginparwidth}{0.8in}
\newcommand{\changedmark}[2]{\marginpar{\raggedright\scriptsize\color{#1}#2}}
\providecommand{\Aa}{}\renewcommand{\Aa}{\mathcal{A}}
\providecommand{\Ad}{}\renewcommand{\Ad}{\widetilde{\mathrm{Ad}}}
\providecommand{\Aev}{}\renewcommand{\Aev}{\mathcal{A}^{+}}
\providecommand{\Aod}{}\renewcommand{\Aod}{\mathcal{A}^{-}}
\providecommand{\B}{}\renewcommand{\B}{\mathrm{B}}
\providecommand{\C}{}\renewcommand{\C}{\mathbb{C}}
\providecommand{\CAR}{}\renewcommand{\CAR}{\mathrm{CAR}}
\providecommand{\CS}{}\renewcommand{\CS}{{\mathbb{S}^{+}}}
\providecommand{\CSm}{}\renewcommand{\CSm}{{\mathbb{S}^{-}}}
\providecommand{\Cl}{}\renewcommand{\Cl}{\mathrm{Cl}}
\providecommand{\Con}{}\renewcommand{\Con}{\mathcal{C}}
\providecommand{\Dl}{}\renewcommand{\Dl}{\Delta}
\providecommand{\EN}{}\renewcommand{\EN}{E_{N}}
\providecommand{\End}{}\renewcommand{\End}{\mathrm{End}}
\providecommand{\F}{}\renewcommand{\F}{\mathbb{F}}
\providecommand{\G}{}\renewcommand{\G}{\Gamma}
\providecommand{\HH}{}\renewcommand{\HH}{\mathbb{H}}
\providecommand{\IS}{}\renewcommand{\IS}{I\sigma_3}
\providecommand{\J}{}\renewcommand{\J}{\mathbf{J}}
\providecommand{\LH}{}\renewcommand{\LH}{L_{H}}
\providecommand{\Mat}{}\renewcommand{\Mat}{\mathrm{M}}
\providecommand{\Mo}{}\renewcommand{\Mo}{M^{\circ}}
\providecommand{\Mp}{}\renewcommand{\Mp}{M^{+}}
\providecommand{\Msun}{}\renewcommand{\Msun}{M_{\odot}}
\providecommand{\Mt}{}\renewcommand{\Mt}{\widetilde{M}}
\providecommand{\Ncone}{}\renewcommand{\Ncone}{\mathcal{N}}
\providecommand{\Neff}{}\renewcommand{\Neff}{N_{\mathrm{eff}}}
\providecommand{\PCS}{}\renewcommand{\PCS}{\mathbb{P}{\mathbb{S}^{+}}}
\providecommand{\PT}{}\renewcommand{\PT}{\mathbb{PT}}
\providecommand{\Pin}{}\renewcommand{\Pin}{\mathrm{Pin}}
\providecommand{\Q}{}\renewcommand{\Q}{\mathbb{Q}}
\providecommand{\R}{}\renewcommand{\R}{\mathbb{R}}
\providecommand{\Ric}{}\renewcommand{\Ric}{\mathrm{Ric}}
\providecommand{\SBH}{}\renewcommand{\SBH}{S_{\mathrm{BH}}}
\providecommand{\Sc}{}\renewcommand{\Sc}{\mathbb{S}}
\providecommand{\Sch}{}\renewcommand{\Sch}{\mathrm{P}}
\providecommand{\Sig}{}\renewcommand{\Sig}{\Sigma}
\providecommand{\Som}{}\renewcommand{\Som}{S_\omega}
\providecommand{\Sp}{}\renewcommand{\Sp}{\mathbb{S}}
\providecommand{\Sph}{}\renewcommand{\Sph}{\mathrm{S}}
\providecommand{\Spi}{}\renewcommand{\Spi}{S_\pi}
\providecommand{\Spin}{}\renewcommand{\Spin}{\mathrm{Spin}}
\providecommand{\Sth}{}\renewcommand{\Sth}{S_{\mathrm{th}}}
\providecommand{\T}{}\renewcommand{\T}{\mathbb{T}}
\providecommand{\TH}{}\renewcommand{\TH}{T_{\mathrm{H}}}
\providecommand{\Tr}{}\renewcommand{\Tr}{\operatorname{Tr}}
\providecommand{\UHF}{}\renewcommand{\UHF}{\mathrm{UHF}}
\providecommand{\Z}{}\renewcommand{\Z}{\mathbb{Z}}
\providecommand{\Zc}{}\renewcommand{\Zc}{\mathrm{Z}}
\providecommand{\cmg}{}\renewcommand{\cmg}{\mathrm{cm^{2}\,g^{-1}}}
\providecommand{\dS}{}\renewcommand{\dS}{\mathrm{dS}}
\providecommand{\dsc}{}\renewcommand{\dsc}{\mathrm{dsc}}
\providecommand{\dtPG}{}\renewcommand{\dtPG}{dt}
\providecommand{\dtS}{}\renewcommand{\dtS}{dt_{\mathrm{S}}}
\providecommand{\ellv}{}\renewcommand{\ellv}{\ell}
\providecommand{\er}{}\renewcommand{\er}{e_r}
\providecommand{\et}{}\renewcommand{\et}{e_t}
\providecommand{\ev}{}\renewcommand{\ev}{\mathrm{ev}}
\providecommand{\gl}{}\renewcommand{\gl}{\mathfrak{gl}}
\providecommand{\hatV}{}\renewcommand{\hatV}{\widehat{V}}
\providecommand{\hb}{}\renewcommand{\hb}{\bar{h}}
\providecommand{\hf}{}\renewcommand{\hf}{\hat{f}}
\providecommand{\id}{}\renewcommand{\id}{\mathrm{id}}
\providecommand{\kpc}{}\renewcommand{\kpc}{\mathrm{kpc}}
\providecommand{\lP}{}\renewcommand{\lP}{\ell_{P}}
\providecommand{\mD}{}\renewcommand{\mD}{\mathbf{m}_{D}}
\providecommand{\mb}{}\renewcommand{\mb}{\bar m}
\providecommand{\mnu}{}\renewcommand{\mnu}{\Sigma m_{\nu}}
\providecommand{\nInf}{}\renewcommand{\nInf}{n_\infty}
\providecommand{\nO}{}\renewcommand{\nO}{n_0}
\providecommand{\ninf}{}\renewcommand{\ninf}{n_{\infty}}
\providecommand{\nz}{}\renewcommand{\nz}{n_{0}}
\providecommand{\proj}{}\renewcommand{\proj}{\operatorname{proj}}
\providecommand{\pv}{}\renewcommand{\pv}{\pi}
\providecommand{\rad}{}\renewcommand{\rad}{\mathrm{rad}}
\providecommand{\sgn}{}\renewcommand{\sgn}{\operatorname{sgn}}
\providecommand{\shA}{}\renewcommand{\shA}{\mathcal{A}}
\providecommand{\shC}{}\renewcommand{\shC}{\mathcal{C}}
\providecommand{\shE}{}\renewcommand{\shE}{\mathcal{E}}
\providecommand{\shJ}{}\renewcommand{\shJ}{\mathcal{J}}
\providecommand{\shK}{}\renewcommand{\shK}{\mathcal{K}}
\providecommand{\shL}{}\renewcommand{\shL}{L}
\providecommand{\shW}{}\renewcommand{\shW}{\mathcal{W}}
\providecommand{\sll}{}\renewcommand{\sll}{\mathfrak{sl}}
\providecommand{\so}{}\renewcommand{\so}{\mathfrak{so}}
\providecommand{\su}{}\renewcommand{\su}{\mathfrak{su}}
\providecommand{\supp}{}\renewcommand{\supp}{\operatorname{supp}}
\providecommand{\tr}{}\renewcommand{\tr}{\operatorname{tr}}
\providecommand{\uu}{}\renewcommand{\uu}{\mathfrak{u}}
\providecommand{\br}{}\renewcommand{\br}[1]{\langle #1\rangle}
\providecommand{\bra}{}\renewcommand{\bra}[1]{\langle #1|}
\providecommand{\file}{}\renewcommand{\file}[1]{\texttt{#1}}
\providecommand{\hp}{}\renewcommand{\hp}[2]{( #1 \,|\, #2 )}
\providecommand{\ip}{}\renewcommand{\ip}[2]{\langle #1, #2\rangle}
\providecommand{\ket}{}\renewcommand{\ket}[1]{|#1\rangle}
\providecommand{\rel}{}\renewcommand{\rel}[2]{D\bigl(#1 \,\|\, #2\bigr)}
\providecommand{\rev}{}\renewcommand{\rev}[1]{\widetilde{#1}}
\providecommand{\slashed}{}\renewcommand{\slashed}[1]{{#1}\mkern-10mu/}
\title{\bfseries The Transport Principle\\[4pt]\large\itshape Total Division in Clifford Algebras via Minimal Completion}
\author{Leon Butler\\ \normalsize\href{mailto:lbutler.math@gmail.com}{\texttt{lbutler.math@gmail.com}}\\ \normalsize ORCID: \href{https://orcid.org/0009-0006-4449-5617}{0009-0006-4449-5617}}
\date{Generated by the Transport Lab, 27 October 2026 \textperiodcentered{} Transport Papers, T1\\[3pt]\small Official version: \href{https://doi.org/10.5281/zenodo.21743863}{doi:10.5281/zenodo.21743863}}
\begin{document}
\maketitle
\begin{abstract}
\normalsize
The usual response to a zero divisor is exclusion. This paper, the root of the programme, takes the opposite: read the failure as an instruction. Division fails in a real Clifford algebra of nondegenerate form in exactly three ways: at isotropic vectors, central idempotents, and non-central idempotents. We prove the trichotomy exhaustive and, for each failure, the minimal repair and the theorem that makes it forced. From dimension three, isotropic vectors admit no equivariant repair at the level of elements, in the Clifford algebra of any extension: the null rotations shear every candidate partner, and the obstruction is a tautological cohomology class vanishing only for the hyperbolic plane. The ascent repairs instead the space of points: the reciprocal extends totally, continuously and uniquely to the null quadric of $\Cl(p{+}1,q{+}1)$, and the increment $(+1,+1)$ is forced twice over, by a sign argument on lifts and an embedding-independent rank computation. Central idempotents ascend by one generator, making the grade involution inner and fusing the two ideals into $\Mat_2(\Cl^0(p,q))$, minimal among the simple real algebras into which the split algebra embeds unitally; the sign of the adjoined direction is a bit the algebra forgets. Non-central idempotents demand nothing, because $eee = e$ where $vvv = 0$. Under iteration, Rule 1 fixes the mod-8 clock, fusion occurs at most once, and no unital homomorphism returns to the base. A closing section states the principle and a six-step algorithm. Every computational identity is machine-verified in exact rational arithmetic, except one flagged floating-point Moore--Penrose check.
\end{abstract}
\medskip
\noindent\textbf{MSC 2020.} Primary 15A66; Secondary 11E88, 16E50, 16U99, 51B10, 53A30.\\
\textbf{Keywords.} Clifford algebra; geometric algebra; zero divisor; null vector; idempotent; pseudo-inverse; conformal model; null rotation; Bott periodicity; transport principle.
\newpage
\noindent \textbf{Authorship and disclosure. }The founding idea of the series is the author's: that division by zero in a Clifford algebra should be read not as a prohibition but as a transport into a higher-dimensional algebra, and that the transported division need not be accompanied by an inverse operation. The direction of investigation and the critical review of the results are likewise the author's. The mathematical development and the drafting of the text were carried out in extended collaboration with Claude Fable 5 (Anthropic), an AI system.

\noindent\fbox{\parbox{\dimexpr\linewidth-2\fboxsep-2\fboxrule}{\textbf{This is not the version of record.} The official version of this paper is the one deposited on Zenodo, \href{https://doi.org/10.5281/zenodo.21743863}{doi:10.5281/zenodo.21743863}. This PDF was generated by the Transport Lab on 27 October 2026 and assembled from approved parts of its store: the results, proofs and text of the paper as the store holds them, each result with its premises and machine checks. It is regenerated when a result it contains is corrected, so it may differ from the deposited paper. Since the deposited version 3.0: 2 results changed, 14 reworded (wording, notation or citations only) and 17 cited results added. Each revised or added result is marked in the margin (\textcolor{BrickRed}{corrected}, \textcolor{BurntOrange}{changed}, \textcolor{gray}{reworded}), and the last section, \emph{Changes since the deposited version}, lists them with the reason for each.}}
\newpage

\section{Introduction}

\phantomsection\label{c:note.transport.setting-introduction}
In ordinary arithmetic, division fails at exactly one number, and the failure is handled by prohibition: division by zero is forbidden, and every calculation routes around it. This paper is about the same failure in a richer setting. A Clifford algebra is the algebra of a geometry: its elements are built from the vectors of a space, and its one rule is that a vector times itself is its squared length. In such an algebra most elements can be divided by. Every vector of nonzero length has an exact inverse, and these invertible elements are plentiful enough to build every symmetry of the underlying space \cite{HestenesSobczyk1984, DoranLasenby2003}. But with three exceptions (the real numbers, the complex numbers and the quaternions, the three Clifford algebras that are themselves division algebras), every real Clifford algebra also contains nonzero elements that cannot be inverted at all: elements that multiply with some partner to give exactly zero, the \emph{zero divisors}. Where they occur, division fails, and the standard response is arithmetic's response: exclude them and work around the exclusion.

The response examined here is the opposite one, and it is the founding idea of the programme this paper begins: read each failure of division as an instruction. A failure, precisely located, tells you what the algebra is missing. The theorems of this paper show that it also tells you the smallest thing that can be added to supply it, and that the addition is forced, not chosen.

The stance has a precedent in the foundations of mathematics, and the precedent was in fact the founding idea's model. G\"odel proved that any consistent system strong enough for arithmetic contains true statements it cannot prove. The famous part is the impossibility; the instructive part is its precision. The unprovable statement is informative: it names, exactly, the stronger system in which it becomes a theorem, and that stronger system contains new unprovable statements of its own, without end. This paper proves that division by zero in a Clifford algebra has the same structure. The impossibility is internal and exact: an inverse of a null vector $v$ would satisfy $vv^{-1} = 1$, and multiplying by $v$ reduces it to $v = 0$ (Lemma \ref{c:transport.noinverse}). The failure determines its own repair, exactly one hyperbolic plane, which is the minimality half of the type-N theorems. And the completed algebra carries a new null cone of its own, so the ascent never terminates; the resulting tower is the subject of T3 \cite{ButlerT3}. The parallel is structural, not a theorem, since provability and invertibility belong to different categories, and nothing below depends on it. What it supplies is the right way to read what follows: the limit of a system indicates the system that removes it, and mathematics has met this kind of obstruction before.

The failures come in exactly three kinds, each with a concrete description. A \emph{null vector} is a direction whose squared length is zero; in spacetime these are the directions light travels, and no algebra containing such a vector can invert it. An \emph{idempotent} is an element with $e^2 = e$, one of the algebra's projections; every projection other than $0$ and $1$ annihilates its complementary projection, so each is a zero divisor. And in certain signatures the algebra splits into two independent halves that annihilate each other; the \emph{central} idempotents that effect this split are the deepest failure of the three, because no symmetry of the algebra can even exchange the halves. We prove this list complete: every zero divisor belongs to one of the three kinds. For each kind the paper then proves the repair. The null vector cannot be repaired as an element: once the space has more than two dimensions, no assignment of pseudo-inverses (the standard weakened substitute for an inverse, recalled in the section ``Setting, and the classical facts the proofs use'') can respect the symmetries of the space, whether valued in the algebra itself, in the Clifford algebra of any extension of the space, or in a matrix algebra over one (the scope of Theorem \ref{c:transport.no-go-level-elements} and Remark \ref{c:note.transport.obstruction-parabolic}); the two-dimensional hyperbolic plane is the only exception, and an instructive one. What can be repaired is the space of points itself, by a completion that costs exactly one new plane and makes the reciprocal total; no smaller completion works. The split algebra is repaired by one added generator, which fuses its two halves into a single simple algebra, again minimally. The ordinary projection turns out to need no repair at all: it is its own pseudo-inverse. And the repairs compose cleanly, so a repair never creates a failure it cannot finish.

The plan of the paper is the plan just sketched, now with theorems attached. First classify: the three kinds have sharp existence loci, and the classification is exhaustive (the section ``A taxonomy of obstructions''). Then repair, one kind per section. \emph{Isotropic vectors} exist precisely in mixed signature; they are the most geometric of the obstructions and the only genuinely hard one, and their treatment takes the longest section of the paper (the section ``Type N: the isotropic obstruction, and Rule 1''). \emph{Central idempotents} exist precisely in the split signatures $p - q \equiv 1 \pmod 4$, where the algebra splits into two ideals that no inner symmetry can exchange, and where, on the one-dimensional extensions of the even-dimensional algebras, the repair turns out to coincide with the first rule's ascent (the section ``Type C: ideal fusion, and Rule 2''). \emph{Non-central idempotents} exist everywhere else, and demand nothing at all (the section ``Type P: the rule of no rule''). Each repair comes with the theorem that makes it forced rather than convenient, the bookkeeping of iterated repairs is settled in the section ``The mod-8 clock and the one-way door'', and the section ``The transport principle, and the algorithm'' states what it all amounts to: a general principle, and a six-step procedure whose every step is licensed by a theorem proved earlier in the paper. the table accompanying \ref{c:transport.transport-algorithm} there is the paper's contribution list in operational form.

\subsection*{Prior art}
\phantomsection\label{c:note.transport.setting-prior-art}
Two bodies of classical material are used, and neither is claimed.

The structure theory of the real Clifford algebras (the isomorphisms $\Cl(p,q) \cong \Mat_k(\mathbb{F})$ or $\Mat_k(\mathbb{F}) \oplus \Mat_k(\mathbb{F})$ with $\mathbb{F} \in \{\R, \C, \HH\}$ determined by $p - q$ mod $8$) goes back to Cartan and to the modern treatments of Atiyah, Bott and Shapiro \cite{ABS1964}, Porteous \cite{Porteous1995} and Lounesto \cite{Lounesto2001}; the quadratic-form background is in Lam \cite{Lam2005} and Chevalley \cite{Chevalley1997}. The computation of the centre and of the sign of $\omega^2$ (Lemmas \ref{c:transport.centre} and \ref{c:transport.sign}) is folklore, restated here and proved in the primer's toolkit \cite[Section 9]{ButlerT0} because the type-C theorem rests on it. Von Neumann regularity of semisimple rings, on which the exhaustiveness theorem rests, is von Neumann's \cite{vonNeumann1936}; the observation that regularity localises zero divisors at idempotents (Lemma \ref{c:transport.regularity-localises-zero-divisors-idempotents}) is elementary and surely known, though we have not found it stated in the Clifford literature in the form used here. The idempotent structure of Clifford algebras is surveyed by Lounesto \cite{Lounesto2001} and, computationally, by Ab{\l}amowicz \cite{Ablamowicz1998}.

The point-level machinery of the type-N section is the \emph{conformal model} of geometric algebra and is thoroughly established. The representation of the points of a quadratic space as null rays in a space two dimensions larger, via the embedding $X(x) = \nO + x + \tfrac12 Q(x)\nInf$, together with the realisation of conformal transformations, inversions included, as versors of $\Cl(p{+}1,q{+}1)$, originates in the program of Hestenes and coworkers \cite{Hestenes2001, HLR2001}, has antecedents in the classical projective-quadric (M\"obius and Lie sphere) constructions of the nineteenth century, was given a comprehensive computational treatment by Dorst, Fontijne and Mann \cite{DFM2007}, is standard in the physics literature through Doran and Lasenby \cite{DoranLasenby2003}, and underlies the algebraic theory of conformal groups via Clifford algebras developed by Angl\`es \cite{Angles2008}. Nothing in the section ``Setting, and the classical facts the proofs use'' and the section ``The completion: the conformal model'' is claimed as new. The reader familiar with conformal geometric algebra may skim them for notation; the reader \emph{new} to geometric algebra should begin with the companion primer \cite{ButlerT0}, which develops all of this slowly and with worked examples.

\subsection*{Contribution}
\phantomsection\label{c:note.transport.setting-contribution}
What this paper adds is not the classification and not the conformal model, but the reading they jointly permit: the organisation of the zero divisors of $\Cl(p,q)$ into three obstruction types with sharp existence loci; the no-go theorem that closes the element-level repair of the isotropic type, with its cohomological form; the minimality theorems that make each rule forced, in a representation-independent sense; and the change of reading these mandate, namely that the conformal construction is not an available convenience but the unique minimal completion demanded by the failure of division itself.

At the standard of complete proofs, with every theorem carrying its own hypotheses and every computational identity machine-verified, the paper proves: that the zero divisors of $\Cl(p,q)$ fall into exactly three types with sharp existence loci, and that the trichotomy is exhaustive because every zero divisor is annihilated by a nontrivial idempotent (the section ``A taxonomy of obstructions''); that isotropic vectors admit no equivariant element-level repair in the Clifford algebra of any extension once $n \geq 3$, the obstruction being the tautological class of the unipotent radical, whose vanishing locus is exactly the hyperbolic atom (the section ``The refuted conjecture, and the no-go'' and the section ``The obstruction as a cohomology class''); that the reciprocal nevertheless extends totally, continuously and uniquely to the null quadric of the completion (the section ``The completion: the conformal model'' and the section ``Totality of the transported reciprocal''); that the increment $(+1,+1)$ is forced twice over, by a sign argument on lifts and by a rank computation on ambient quadrics that no exotic embedding evades (the section ``Minimality among null-lift completions'' and the section ``Minimality against arbitrary embeddings''); that central idempotents fuse under a one-generator ascent into $\Mat_2(\Cl^0(p,q))$, minimal among all simple real algebras, with the sign of the adjoined direction a bit the algebra forgets, and that on every one-dimensional extension of an even-dimensional algebra exactly one sign obstructs, its ideals are copies of the base, and the opposite-sign repair is the Rule 1 ascent, so that the two rules coincide there (the section ``Type C: ideal fusion, and Rule 2''); that non-central idempotents demand nothing, because $eee = e$ where $vvv = 0$, so there is no fourth rule (the section ``Type P: the rule of no rule''); and that the bookkeeping closes, Rule 1 fixing the mod-8 clock, fusion occurring at most once, and no unital homomorphism returning to the base (the section ``The mod-8 clock and the one-way door'').

\subsection*{The destination, in advance}
\phantomsection\label{c:note.transport.setting-destination-advance}
Section \ref{c:transport.transport-algorithm} turns the theorems into a six-step procedure. It asks, in order:

   
\begin{enumerate}
  
\item Where does division fail?
\item Which of the three obstructions is it?
\item What is the minimal repair, and what does it adjoin?
\item How is the value at the obstruction computed?
\item What does descent back to the original algebra cost?
\item Does repeating the repair terminate?
\end{enumerate}
 \noindent The answers take the form of an algorithm: six steps, each with a theorem licensing it. It is set out as the table accompanying \ref{c:transport.transport-algorithm} in Section \ref{c:transport.transport-algorithm}, which names every step's instrument and its warrant, and which T5 follows in order.

Everything between here and there builds that algorithm. Each section proves the results that one or more of the steps invoke, one obstruction type at a time, and no step is stated before the theorem that licenses it, which is why the procedure itself comes last rather than first. A reader who would rather see the destination before the route can read the table accompanying \ref{c:transport.transport-algorithm} now and return here.

\emph{A reading route.} Read straight through for the mathematics in the order it was found. To enter directly: the no-go and its cohomological form are the section ``The refuted conjecture, and the no-go'' and the section ``The obstruction as a cohomology class''; the completion, its totality and its two minimality theorems are the section ``The completion: the conformal model'' to the section ``Minimality against arbitrary embeddings''; the fusion theorem is the section ``The fusion theorem''; the clock is the section ``The mod-8 clock and the one-way door''; the algorithm is Section \ref{c:transport.transport-algorithm}. the section ``Setting, and the classical facts the proofs use'' is notation and imported classical facts, and is meant to be consulted rather than read.

\section{Setting, and the classical facts the proofs use}

\phantomsection\label{c:note.transport.setting-setting-classical-facts-proofs-use}
This section fixes notation and assembles the facts the proofs use, under a strict division of labour with the companion primer \cite{ButlerT0}. The \emph{language} belongs to the primer: the geometric product, blades, involutions, versors and the groups they generate, the classification, and the structural lemmas about them, which are restated below with their proofs imported from the primer's toolkit \cite[Section 9]{ButlerT0}, each statement carrying its number there (the classification theorem alone is not reproved in the toolkit, which states it and cites the classical sources for its induction). The \emph{failure of division} belongs to this paper: every fact about zero divisors, null vectors as obstructions, pseudo-inverses and idempotents is stated and proved here, because that failure is the paper's subject; the one deliberate exception is the splitting lemma (Lemma \ref{c:transport.splitting-along-hyperbolic-plane}), a basis fact about hyperbolic planes rather than about failure, whose proof stays in the toolkit. The definitions specific to this paper (extension, equivariance) are likewise stated here and belong to this paper.

\subsection*{Quadratic spaces and Clifford algebras}
\phantomsection\label{c:note.transport.setting-quadratic-spaces-clifford-algebras}
$(\mathbb{V}, Q)$ is a finite-dimensional real quadratic space with polarisation $a \cdot b = \tfrac12(Q(a+b) - Q(a) - Q(b))$, so $a \cdot a = Q(a)$, and is \emph{nondegenerate} when $a \cdot b = 0$ for all $b$ forces $a = 0$. It then has an orthogonal basis $e_1, \dots, e_n$ with $e_i^{\,2} \mathrel{:=} Q(e_i) = +1$ for $p$ vectors and $-1$ for the remaining $q = n-p$; the \emph{signature} $(p,q)$ is an invariant (Sylvester), and is \emph{mixed} when $p, q \geq 1$.

The \emph{Clifford algebra} $\Cl(p,q) = \Cl(\mathbb{V}, Q)$ is generated by $\mathbb{V}$ subject to $v^2 = Q(v)$. Polarising gives $ab + ba = 2\,a \cdot b$, so orthogonal vectors anticommute; $ab = a \cdot b + a \wedge b$ with $a \wedge b \mathrel{:=} \tfrac12(ab-ba)$ the \emph{outer product}, pure grade two. The ordered products $e_{i_1} \cdots e_{i_k}$ over an orthogonal basis, $i_1 < \dots < i_k$, the \emph{basis blades} with $1$ as the empty product, form a linear basis, so $\dim \Cl(\mathbb{V},Q) = 2^n$; this holds for degenerate forms too \cite{Chevalley1997}, which we need once (Lemma \ref{c:transport.splitting-along-hyperbolic-plane}). An element has \emph{grade} $k$ if it combines blades of length $k$, basis-independently \cite{Porteous1995}. A vector with $Q(v) \neq 0$ is invertible, $v^{-1} = v/Q(v)$; a nonzero vector with $Q(v) = 0$ is \emph{null}, and null vectors exist exactly in mixed signature, where they form a cone \cite[Sections 2 to 4]{ButlerT0}. A word on the two notations: we write $\Cl(p,q)$ whenever only the signature matters, and $\Cl(\mathbb{V})$ or $\Cl(\mathbb{V}')$ for the algebra of a \emph{specific} quadratic space, chiefly an extension quantified over in a theorem, whose form is not assumed nondegenerate and so may carry no signature label at all.

\phantomsection\label{c:note.transport.setting-quadratic-spaces-clifford-algebras-2}
\textbf{Scope of hypotheses.} Throughout, $(\mathbb{V}, Q)$ is nondegenerate; nothing else is assumed globally, and each theorem carries its own hypotheses. The mixed-signature and dimension conditions $p, q \geq 1$ and $n \geq 3$ are local to the no-go theorem (Theorem \ref{c:transport.no-go-level-elements}), which needs them: they are exactly what guarantees that null vectors exist and that the null rotations shearing their candidate partners exist. Remarks \ref{c:note.transport.obstruction-parabolic} and \ref{c:note.transport.hyperbolic-atom} delimit the boundary: definite signatures pose no element-level question at all, and $n = 2$ is the sharp exception where the question is posed and answered affirmatively. The completion theorems, namely totality (Theorem \ref{c:transport.totality}), minimality (Theorem \ref{c:transport.minimality}) and the normal form (Proposition \ref{c:transport.swap-rescale-normal-form}), hold for \emph{every} nondegenerate signature, definite included; the section ``The Euclidean plane: inversion in a circle on the Riemann sphere'' exercises them over the base of signature $(2,0)$.

The two canonical involutions are written as in the primer \cite[Section 5]{ButlerT0}: the \emph{grade involution} $x \mapsto \hat{x}$ (the automorphism with $\hat{v} = -v$ on vectors) and the \emph{reversion} $x \mapsto \tilde{x}$ (the anti-automorphism fixing vectors). One consequence is used repeatedly enough to display, the commutation rule \cite[Lemma 9.1]{ButlerT0}: if $x$ lies in the subalgebra generated by vectors orthogonal to a vector $v$, then moving $v$ past $x$ introduces exactly a grade involution,
\begin{equation}
v\,x = \hat{x}\,v, \qquad x\,v = v\,\hat{x}.
\end{equation}

\subsection*{Null vectors, zero divisors, and hyperbolic planes}
\phantomsection\label{c:note.transport.setting-null-vectors-zero-divisors-hyperbolic}
Division fails first, and most geometrically, at the null vectors, and the facts of this subsection are the paper's own.

\begin{lemma}[A null vector has no inverse]\label{c:transport.noinverse}
\changedmark{gray}{reworded}
A null vector $v \neq 0$ has no left inverse and no right inverse in $\Cl(p,q)$, nor in any associative unital algebra containing $\Cl(p,q)$ as a unital subalgebra.
\end{lemma}
\begin{proof}
The relation $v^2 = Q(v) = 0$ (\ref{c:std-algebra.vector-squares}) persists in any algebra containing $\Cl(p,q)$ as a subalgebra, since the multiplication restricts. If $vw = 1$ then $v = v(vw) = v^2 w = 0$; if $wv = 1$ then $v = (wv)v = wv^2 = 0$. Either way a contradiction.

\end{proof}

\phantomsection\label{c:note.transport.setting-null-vectors-zero-divisors-hyperbolic-2}
Null vectors are moreover two-sided zero divisors. The obstruction they pose is packaged in the reciprocal map
\[
\rho : \mathbb{V} \setminus \{Q = 0\} \to \mathbb{V}, \qquad \rho(x) = x / Q(x),
\]
which is partial exactly on $\{Q = 0\}$, the null cone together with the origin. Both failure modes are resolved by the same completion (Theorem \ref{c:transport.totality}); only the null cone poses the element-level question of the section ``The refuted conjecture, and the no-go'', the pseudo-inverse condition being vacuous at the origin.

Null vectors come in pairs. If $v$ is null and $Q$ is nondegenerate, there is a vector $a$ with $v \cdot a \neq 0$; rescaling so that $v \cdot a = \tfrac12$ and setting $w = a - Q(a)\,v$ gives
\[
Q(w) = Q(a) - 2\,Q(a)\,(v \cdot a) = 0, \qquad v \cdot w = v \cdot a = \tfrac12 :
\]
a \emph{hyperbolic pair} $(v, w)$. The plane $H = \mathrm{span}\{v, w\}$ is a \emph{hyperbolic plane}: its Gram matrix on $(v,w)$ has determinant $-\tfrac14$, so $H$ is nondegenerate of signature $(1,1)$, and conversely every two-dimensional nondegenerate space containing a null vector is isometric to $H$ \cite{Lam2005}. Inside the Clifford algebra a hyperbolic pair satisfies
\begin{equation}
v^2 = w^2 = 0, \qquad vw + wv = 1, \qquad \text{hence} \quad vwv = (1 - wv)v = v .
\end{equation}
The last identity says $w$ is a \emph{pseudo-inverse} of $v$ in the sense of von Neumann: division by $v$ cannot be undone exactly, but $w$ undoes it up to the defect of $v$. This is the weakened form of inversion around which the section ``The refuted conjecture, and the no-go'' revolves.

\subsection*{The symmetry action, and equivariance}
\phantomsection\label{c:note.transport.setting-symmetry-action-equivariance}
Versors, the reflection formula, and the groups $\Pin(p,q)$, $\Spin(p,q)$, $\Spin^0(p,q)$ and the rotors are as in the primer \cite[Section 6]{ButlerT0}. What the statements below need is the action: a versor $g$ (a product of invertible vectors) acts on vectors by the twisted adjoint $\Ad_g(x) = \hat{g}\, x\, g^{-1}$, realising the full orthogonal group $O(p,q)$, and the action extends to the whole algebra as
\begin{equation}
\tau_g(x) \;=\;
\begin{cases}
\,g\, x\, g^{-1} & g \text{ even},\\[2pt]
\,g\, \hat{x}\, g^{-1} & g \text{ odd},
\end{cases}
\end{equation}
an algebra automorphism restricting on vectors to $\Ad_g$ \cite[Lemma 9.2]{ButlerT0}. For even elements, which are all that our proofs use, $\tau_g$ is plain conjugation.

\begin{definition}[Equivariance]\label{c:obj.transport.equivariance}
Let $\mathbb{V}' \supseteq \mathbb{V}$ be an extension of quadratic spaces (Definition \ref{c:obj.transport.extension-quadratic-spaces}) and $G$ a subgroup of $\Pin(p,q)$, acting on the null cone $\Ncone \subset \mathbb{V}$ by $\Ad$ (which preserves $\mathbb{V}$ and its null cone, being an isometry of $\mathbb{V}$) and on $\Cl(\mathbb{V}')$ by $\tau$ through the inclusion $\Cl(\mathbb{V}) \subseteq \Cl(\mathbb{V}')$. A map $\sigma : \Ncone \to \Cl(\mathbb{V}')$ is \emph{$G$-equivariant} if
\[
\sigma\bigl(\Ad_g(v)\bigr) \;=\; \tau_g\bigl(\sigma(v)\bigr) \qquad \text{for all } g \in G,\ v \in \Ncone .
\]
\end{definition}

\phantomsection\label{c:note.transport.setting-symmetry-action-equivariance-2}
Equivariance is the precise form of ``respecting the symmetries'': if two null vectors are related by a symmetry of the space, their assigned partners must be related by the same symmetry. The consequence used below is the case of a stabiliser: if $g$ fixes $v$ then $\sigma(v) = \tau_g(\sigma(v))$.

\begin{definition}[Extension of quadratic spaces]\label{c:obj.transport.extension-quadratic-spaces}
An \emph{extension} of the nondegenerate quadratic space $(\mathbb{V}, Q)$ is a real quadratic space $(\mathbb{V}', Q')$ with $\mathbb{V} \subseteq \mathbb{V}'$ and $Q'|_{\mathbb{V}} = Q$. The form $Q'$ is not assumed nondegenerate. Because $Q$ is nondegenerate on $\mathbb{V}$, the extension splits orthogonally and automatically: $\mathbb{V}' = \mathbb{V} \perp \mathbb{V}^{\perp}$, where $\mathbb{V}^{\perp} = \{ a \in \mathbb{V}' : a \cdot x = 0 \ \forall x \in \mathbb{V}\}$ \cite{Lam2005}. Extensions in this sense are exactly the ``ascents'' quantified over in Theorem \ref{c:transport.no-go-level-elements}.
\end{definition}

\begin{lemma}[Splitting along a hyperbolic plane]\label{c:transport.splitting-along-hyperbolic-plane}
\changedmark{gray}{reworded}
Let $\mathbb{V}' = H_0 \perp U'$ be an orthogonal decomposition with $H_0$ a hyperbolic plane spanned by a hyperbolic pair $(v, w_0)$; the form on $U'$ may be degenerate. Then:
\begin{enumerate}
\item $\{1, v, w_0, v w_0\}$ is a linear basis of $\Cl(H_0)$, and every element of $\Cl(\mathbb{V}')$ admits a \emph{unique} expansion
\[
\sigma \;=\; \alpha + v\beta + w_0\gamma + v w_0\,\delta, \qquad \alpha, \beta, \gamma, \delta \in \Cl(U').
\]
\item Left multiplication by $v$ is injective on $\Cl(U')$.
\end{enumerate}
\end{lemma}
\begin{proof}
Restated from \ref{c:primer.splitting-along-hyperbolic-plane} (T0), where it is proved or imported with its sources.

\end{proof}

\subsection*{Classification, the centre, and regularity}
\phantomsection\label{c:note.transport.setting-classification-centre-regularity}
The taxonomy of the section ``A taxonomy of obstructions'' rests on three classical inputs, recorded here. The first is the structure theory of the real Clifford algebras, together with the centre and pseudoscalar computations behind it, imported from the primer's toolkit because they are structure theory, the language's side of the division of labour. The second is the regularity lemma, proved here: it is a statement about zero divisors, this paper's subject, and we have not found it stated in the Clifford literature in the form used.

\begin{theorem}[Classification of the real Clifford algebras]\label{c:transport.classification-real-clifford-algebras}
\changedmark{gray}{reworded}
Let $d \equiv p - q \pmod 8$. Then $\Cl(p,q)$ is isomorphic, as a real associative algebra, to
\[
\begin{array}{llll}
d = 0:\ \Mat_{k}(\R), &
d = 1:\ \Mat_{k}(\R) \oplus \Mat_{k}(\R), &
d = 2:\ \Mat_{k}(\R), &
d = 3:\ \Mat_{k}(\C), \\[2pt]
d = 4:\ \Mat_{k}(\HH), &
d = 5:\ \Mat_{k}(\HH) \oplus \Mat_{k}(\HH), &
d = 6:\ \Mat_{k}(\HH), &
d = 7:\ \Mat_{k}(\C),
\end{array}
\]
with $k$ in each case determined by $\dim_\R \Cl(p,q) = 2^n$. In particular $\Cl(p,q)$ is a simple algebra except when $d \in \{1, 5\}$, i.e.\ except when $p - q \equiv 1 \pmod 4$, where it is the direct sum of two isomorphic simple algebras.
\end{theorem}
\begin{proof}
Restated from \ref{c:primer.classification-real-clifford-algebras} (T0), where it is proved or imported with its sources.

\end{proof}

\begin{lemma}[The centre]\label{c:transport.centre}
\changedmark{gray}{reworded}
For $n \geq 1$, the centre of $\Cl(p,q)$ is $\R \cdot 1$ if $n$ is even, and $\R \cdot 1 \oplus \R \cdot \omega$ if $n$ is odd, where $\omega = e_1 e_2 \cdots e_n$ is the pseudoscalar.
\end{lemma}
\begin{proof}
Restated from \ref{c:primer.centre} (T0), where it is proved or imported with its sources.

\end{proof}

\begin{lemma}[The sign of $\omega^2$]\label{c:transport.sign}
\changedmark{gray}{reworded}
$\displaystyle \omega^2 = (-1)^{\tfrac{n(n-1)}{2}} (-1)^q = (-1)^{\tfrac{(p-q)(p-q-1)}{2}}$. In particular $\omega^2 = +1$ if and only if $p - q \equiv 0$ or $1 \pmod 4$.
\end{lemma}
\begin{proof}
Restated from \ref{c:primer.sign} (T0), where it is proved or imported with its sources.

\end{proof}

\phantomsection\label{c:note.transport.setting-classification-centre-regularity-2}
Finally, the regularity input. A ring is \emph{von Neumann regular} if for every $a$ there is an $x$ with $axa = a$ \cite{vonNeumann1936}.

\begin{lemma}[Regularity localises zero divisors at idempotents]\label{c:transport.regularity-localises-zero-divisors-idempotents}
Every semisimple Artinian ring, in particular every real Clifford algebra of nondegenerate form, is von Neumann regular. Consequently, if $a$ is a nonzero zero divisor in such a ring $A$, and $x$ satisfies $axa = a$, then $e = xa$ and $f = ax$ are idempotents with
\[
a(1 - e) = 0, \qquad (1 - f)a = 0, \qquad e, f \notin \{0, 1\}.
\]
Moreover $ae = a$ and $fa = a$.
\end{lemma}
\begin{proof}
By Wedderburn--Artin a semisimple Artinian ring is a finite product of matrix rings over division rings, and it suffices to verify regularity in $\Mat_k(D)$: any matrix $a$ of rank $r$ can be written $a = u \begin{psmallmatrix} 1_r & 0 \\ 0 & 0\end{psmallmatrix} v$ with $u, v$ invertible, and $x = v^{-1} \begin{psmallmatrix} 1_r & 0 \\ 0 & 0 \end{psmallmatrix} u^{-1}$ gives $axa = a$. (Clifford algebras of nondegenerate forms are semisimple by Theorem \ref{c:transport.classification-real-clifford-algebras}.) For the consequence: $e^2 = xaxa = xa = e$, and $a(1-e) = a - axa = 0$; symmetrically for $f$. Also $ae = axa = a$ and $fa = axa = a$. If $e = 1$ then $xa = 1$, so $a$ is left-invertible, hence invertible in an Artinian ring, contradicting that $a$ is a zero divisor; if $e = 0$ then $a = ae = 0$, contradicting $a \neq 0$. The same for $f$.

\end{proof}

\section{A taxonomy of obstructions}

\begin{definition}[Obstruction types]\label{c:obj.transport.obstruction-types}
Let $A = \Cl(p,q)$, nondegenerate. A zero-divisor obstruction of $A$ is classified as:
\begin{itemize}
\item \textbf{Type N (isotropic vector).} A null vector $v \in \mathbb{V}$, $v \neq 0$, $v^2 = Q(v) = 0$, together with the annihilating idempotent $wv$ it acquires through a hyperbolic partner (Example \ref{c:note.transport.hyperbolic-partner-supplies-annihilating-idempotent}). Geometrically: an isotropic line, a direction without a metric partner.
\item \textbf{Type P (non-central idempotent).} An idempotent $e = e^2$ with $e \notin \{0, 1\}$ and $e$ not central. Prototype: $e = \tfrac12(1 + u)$ with $u^2 = 1$, $u \neq \pm 1$. Each such $e$ is a zero divisor via $e(1 - e) = 0$.
\item \textbf{Type C (central idempotent).} An idempotent $e \notin \{0,1\}$ lying in the centre of $A$. The two-sided ideals $eA$ and $(1-e)A$ then annihilate one another and $A$ splits as their direct sum.
\end{itemize}
\end{definition}

\phantomsection\label{c:note.transport.setting-taxonomy-obstructions}
The three types have sharp existence loci, which we now determine. First the global statement that there is anything to classify.

\begin{proposition}[The exceptional division algebras]\label{c:transport.exceptional-division-algebras}
$\Cl(p,q)$ contains no zero divisors if and only if $(p,q) \in \{(0,0), (0,1), (0,2)\}$, where it is $\R$, $\C$, $\HH$ respectively. It contains no nilpotents if and only if $(p,q) \in \{(0,0), (0,1), (0,2), (1,0), (0,3)\}$: the three division algebras together with $\Cl(1,0) \cong \R \oplus \R$ and $\Cl(0,3) \cong \HH \oplus \HH$.
\end{proposition}
\begin{proof}
By Theorem \ref{c:transport.classification-real-clifford-algebras}, $\Cl(p,q)$ is $\Mat_k(\mathbb{F})$ or $\Mat_k(\mathbb{F})^{\oplus 2}$. A matrix algebra $\Mat_k(\mathbb{F})$ with $k \geq 2$ contains the nilpotent matrix unit $E_{12} \neq 0$; a direct sum contains the zero divisors $(x, 0)$; and a division algebra contains neither. So freedom from zero divisors forces a single factor with $k = 1$: $\mathbb{F}$ itself. Dimension count fixes the signatures: $\dim = 2^n = \dim_\R \mathbb{F}$ (\ref{c:std-algebra.blade-basis}) gives $n = 0, 1, 2$ with $\mathbb{F} = \R, \C, \HH$, i.e.\ $(0,0), (0,1), (0,2)$ (the candidates $\Cl(1,0)$ and $\Cl(2,0)$, of the right dimensions, are $\R \oplus \R$ and $\Mat_2(\R)$ by the classification, not $\C$ and $\HH$). Freedom from nilpotents allows in addition direct sums of division algebras: $\Mat_k(\mathbb{F})^{\oplus 2}$ is nilpotent-free iff $k = 1$, and dimension count $2^n = 2\dim_\R \mathbb{F}$ within $d \in \{1,5\}$ gives exactly $\Cl(1,0) \cong \R \oplus \R$ ($n=1$) and $\Cl(0,3) \cong \HH \oplus \HH$ ($n=3$).

\end{proof}

\begin{proposition}[Existence locus of type N]\label{c:transport.existence-locus-type-n}
$\Cl(p,q)$ contains a null vector if and only if the signature is mixed, $pq \neq 0$. In that case the null vectors form a cone of dimension $n - 1$, and every null vector is a two-sided zero divisor without left or right inverse in any unital algebra containing $\Cl(p,q)$.
\end{proposition}
\begin{proof}
If $q = 0$ or $p = 0$ the form is definite and anisotropic: $Q(v) = 0$ forces $v = 0$. If $p, q \geq 1$, then $v = e_1 + e_n$ satisfies $v^2 = 1 - 1 + (e_1e_n + e_ne_1) = 0$ (\ref{c:std-algebra.orthogonal-anticommute}); the solution set of the homogeneous quadratic $Q = 0$ minus the origin is a cone of dimension $n-1$. Non-invertibility in every extension is Lemma \ref{c:transport.noinverse}: from $v^2 = 0$, $vw = 1$ gives $v = v^2 w = 0$.

\end{proof}

\begin{example}[The hyperbolic partner supplies the annihilating idempotent]\label{c:note.transport.hyperbolic-partner-supplies-annihilating-idempotent}
Let $v$ be null and $(v, w)$ a hyperbolic pair, $vw + wv = 1$. Then $e = wv$ is idempotent ($e^2 = w(vw)v = w(1 - wv)v = wv$, using $v^2 = w^2 = 0$), non-central (it fails to commute with $v$: $ev = wv\cdot v = wv^2 = 0$, while $ve = v(wv) = (vwv) = v \neq 0$), and
\[
v(1 - e)\;=\;v - vwv \;=\; v - v \;=\; 0 , \qquad (1-e) \ne 0.
\]
Thus the abstract annihilator idempotent of Lemma \ref{c:transport.regularity-localises-zero-divisors-idempotents} is, for a type-N obstruction, realised by the geometry itself: the pseudo-inverse $w$ of the section ``Null vectors, zero divisors, and hyperbolic planes'' produces it. Type N is in this sense a geometric refinement of type P: a null vector is a zero divisor \emph{whose annihilating idempotent comes from an isotropic line}. What Theorem \ref{c:transport.no-go-level-elements} will prove is that this refinement carries an obstruction the generic type-P case does not: the choice of $w$, hence of $e$, is sheared by the null rotations, and no equivariant choice exists in any ascent.
\end{example}

\begin{example}[Nilpotents without null vectors: the two loci differ]\label{c:note.transport.nilpotents-without-null-vectors-two-loci}
$\Cl(2,0) \cong \Mat_2(\R)$ has definite signature: no null vectors, by Proposition \ref{c:transport.existence-locus-type-n}. It nevertheless contains nilpotents, e.g.
\[
x = (1 + e_1)\,e_2, \qquad x^2 = (1+e_1)\,e_2\,(1+e_1)\,e_2 = (1+e_1)(1-e_1)\,e_2^2 = 0,
\]
using $e_2 e_1 = -e_1 e_2$ and $e_1^2 = e_2^2 = 1$. A loose reading of type N, as comprising null vectors and more generally the elements they generate, therefore needs care: the \emph{nilpotency} locus (nonzero $x$ with $x^k = 0$) is strictly larger than the locus of algebras with isotropic vectors, and it is the latter, the grade-one geometric obstruction, that defines type N and triggers Rule 1. The nilpotent $x$ above is instead type-P-generated: $x = 2\,\varepsilon\, e_2$ with $\varepsilon = \tfrac12(1 + e_1)$ a non-central idempotent, and $x$ is annihilated by the idempotent $1 - \varepsilon$ on the left. In matrix form ($e_1 \mapsto \mathrm{diag}(1,-1)$, $e_2 \mapsto \begin{psmallmatrix} 0 & 1 \\ 1 & 0\end{psmallmatrix}$), $x \mapsto 2E_{12}$: the standard nilpotent matrix unit, built from idempotents and invertible vectors with no isotropic vector anywhere in $\mathbb{V}$.
\end{example}

\begin{theorem}[Existence and classification of type C]\label{c:transport.existence-classification-type-c}
$\Cl(p,q)$ ($n \geq 1$) contains a nontrivial central idempotent if and only if
\[
p - q \equiv 1 \pmod 4 ,
\]
and in that case the nontrivial central idempotents are exactly the two elements
\[
\varepsilon_\pm = \tfrac12\,(1 \pm \omega),
\]
which are exchanged by the grade involution $\hat{\cdot}$ (an \emph{outer} automorphism of the split algebra; no inner automorphism exchanges them, centrality being conjugation-invariant), and satisfy $\varepsilon_+ \varepsilon_- = 0$ and $\varepsilon_+ + \varepsilon_- = 1$, and effect the splitting $\Cl(p,q) = \varepsilon_+ A \,\oplus\, \varepsilon_- A$ into the two simple ideals of Theorem \ref{c:transport.classification-real-clifford-algebras}.
\end{theorem}
\begin{proof}
A central idempotent lies in the centre $\Zc$, computed in Lemma \ref{c:transport.centre}. If $n$ is even, $\Zc = \R$, whose idempotents are $0$ and $1$: no type C. If $n$ is odd, $\Zc = \R[\omega]$ with $\omega^2 = \pm 1$ by Lemma \ref{c:transport.sign}. Case $\omega^2 = -1$: $\Zc \cong \C$, and an idempotent $a + b\omega$ satisfies $a^2 - b^2 = a$ and $2ab = b$; if $b \neq 0$ the second gives $a = \tfrac12$ and the first $-b^2 = \tfrac14$, impossible over $\R$; so $b = 0$ and $a \in \{0,1\}$: no type C. Case $\omega^2 = +1$: $\Zc \cong \R \oplus \R$ via $\varepsilon_\pm$, whose idempotents are $0, 1, \varepsilon_+, \varepsilon_-$; direct check $\varepsilon_\pm^2 = \tfrac14(1 \pm 2\omega + \omega^2) = \varepsilon_\pm$ and $\varepsilon_+\varepsilon_- = \tfrac14(1 - \omega^2) = 0$. By Lemma \ref{c:transport.sign}, $n$ odd and $\omega^2 = +1$ happen together precisely when $p - q$ is odd and $p - q \equiv 0, 1 \pmod 4$, i.e.\ precisely when $p - q \equiv 1 \pmod 4$. For the exchange: for $n$ odd the pseudoscalar is central, so no inner automorphism, in particular no vector conjugation, moves it; the exchanging symmetry must instead be the \emph{grade involution}: $\hat{\omega} = (-1)^n \omega = -\omega$ for $n$ odd (\ref{c:std-algebra.involution-properties}), so the grade involution, implemented on the algebra by conjugation in the larger algebra $\Cl(p,q{+}1)$ or realised abstractly as an outer automorphism, exchanges $\varepsilon_+ \leftrightarrow \varepsilon_-$. Within $A$ itself the two ideals are exchanged by no inner automorphism, centrality being conjugation-invariant, and this is exactly why type C is the deepest obstruction: no inner automorphism can act on it. The splitting into two simple ideals is Theorem \ref{c:transport.classification-real-clifford-algebras} ($d \in \{1,5\}$), each ideal $\varepsilon_\pm A \cong \Mat_k(\mathbb{F})$.

\end{proof}

\begin{remark}\label{c:note.transport.typec-depth}
The proof just given settles a point worth recording: for $n$ odd the pseudoscalar is central, so \emph{no inner automorphism} of $\Cl(p,q)$ exchanges the two ideals; the exchanging symmetry is the grade involution, which is outer on the split algebra, and its realisation by an inner automorphism is available only after the one-generator ascent of Rule 2, where a new generator anticommuting with $\omega$ exists. This sharpens the sense in which Rule 2 is a \emph{repair}: the ascent does not merely contain the fused algebra, it renders inner the very symmetry whose absence defined the obstruction. The observation will be part of the minimality argument in the section ``Type C: ideal fusion, and Rule 2''.
\end{remark}

\begin{proposition}[Existence locus of type P]\label{c:transport.existence-locus-type-p}
$\Cl(p,q)$ contains a non-central idempotent if and only if it is not a direct sum of division algebras, that is, if and only if
\[
(p,q) \;\notin\; \{(0,0),\ (0,1),\ (0,2),\ (1,0),\ (0,3)\} .
\]
In particular, for $n \geq 2$ with $(p,q) \notin \{(0,2), (0,3)\}$, type P is always present, and concretely: any invertible vector pair $a, b$ with $a \cdot b = 0$ and $a^2 b^2 = -1$ (one timelike and one spacelike factor) yields the non-central idempotent $\tfrac12(1 + u)$ with $u = ab$, since then $(ab)^2 = -a^2b^2 = +1$, and every unit blade $u$ with $u^2 = 1$, $u \neq \pm 1$, $u$ non-central yields one.
\end{proposition}
\begin{proof}
In a division algebra, and in a sum of two, every idempotent is a sum of the componentwise trivial ones, hence central: $(x,y)$ idempotent in $D_1 \oplus D_2$ forces $x \in \{0,1\}$, $y \in \{0,1\}$, all four combinations central. Conversely if $\Cl(p,q) \cong \Mat_k(\mathbb{F})$ (or a sum) with $k \geq 2$, the matrix unit $E_{11}$ is a non-central idempotent. The signature list is Proposition \ref{c:transport.exceptional-division-algebras}. For the concrete form: if $u^2 = 1$ then $\tfrac12(1+u)$ is idempotent, and it is central iff $u$ is; a product of two orthogonal invertible vectors squares to $(ab)^2 = -a^2b^2 = +1$ (\ref{c:std-algebra.orthogonal-anticommute}) iff $a^2 b^2 = -1$, i.e.\ one timelike and one spacelike factor, available whenever $pq \neq 0$; in definite signature take instead a grade-one unit $u = e_1$ with $e_1^2 = +1$ ($p \geq 1$), or in the negative-definite case a suitable blade, e.g.\ $u = e_1e_2e_3e_4$ with $u^2 = +1$ for $q \geq 4$; the classification (\ref{c:transport.classification-real-clifford-algebras}) guarantees existence in all listed cases, the explicit blade being a convenience, not the proof.

\end{proof}

\begin{theorem}[Exhaustiveness: every zero divisor is annihilated by an idempotent]\label{c:transport.exhaustiveness-every-zero-divisor-annihilated-idempotent}
Let $A = \Cl(p,q)$ be nondegenerate, $a \in A$ a nonzero zero divisor. Then:
\begin{enumerate}
\item There exist idempotents $e, f \notin \{0, 1\}$ with $a(1-e) = 0$ and $(1-f)a = 0$ (so $a$ both left-annihilates and is right-annihilated by nontrivial idempotents).
\item If $p - q \not\equiv 1 \pmod 4$, every such idempotent is non-central: the obstruction carried by $a$ is of type P, refined to type N exactly when the annihilating idempotent is the $wv$ of a hyperbolic pair, as in Example \ref{c:note.transport.hyperbolic-partner-supplies-annihilating-idempotent}.
\item If $p - q \equiv 1 \pmod 4$, the element $a$ may in addition be annihilated by a \emph{central} idempotent (this happens exactly when $\varepsilon_+ a = 0$ or $\varepsilon_- a = 0$, i.e.\ when $a$ lies entirely in one simple ideal): type C. Every zero divisor not of this kind is again of type P within its ideal(s).
\end{enumerate}
Consequently the taxonomy N/P/C of Definition \ref{c:obj.transport.obstruction-types} is exhaustive: every failure of invertibility in a real Clifford algebra of nondegenerate form is localised at an idempotent, non-central in the simple signatures, possibly central exactly in the signatures $p - q \equiv 1 \pmod 4$, and of isotropic-vector origin exactly in the mixed signatures.
\end{theorem}
\begin{proof}
(1) is Lemma \ref{c:transport.regularity-localises-zero-divisors-idempotents}. (2): for $p - q \not\equiv 1 \pmod 4$ the algebra is simple (Theorem \ref{c:transport.classification-real-clifford-algebras}), so by Theorem \ref{c:transport.existence-classification-type-c} its only central idempotents are $0$ and $1$, excluded in (1). (3): if $a = \varepsilon_- a$ (say), then $\varepsilon_+ a = \varepsilon_+ \varepsilon_- a = 0$ with $\varepsilon_+$ central and nontrivial. If neither central projector annihilates $a$, then both components $\varepsilon_\pm a$ are nonzero; apply Lemma \ref{c:transport.regularity-localises-zero-divisors-idempotents} componentwise: in each simple ideal $\varepsilon_\pm A \cong \Mat_k(\mathbb{F})$, the component $\varepsilon_\pm a$ is either invertible (within the ideal) or annihilated by a non-central idempotent of that ideal; since $a$ is a zero divisor of $A$, not both components can be invertible in their ideals, so at least one is annihilated, and the annihilating idempotent, viewed in $A$, is non-central (it is neither $0$ nor $\varepsilon_\pm$ nor central in $A$, central idempotents of $A$ being classified in Theorem \ref{c:transport.existence-classification-type-c}: any idempotent of $\varepsilon_\pm A$ central in $A$ would be a central idempotent of $A$ inside one ideal, hence $0$ or $\varepsilon_\pm$; the annihilating idempotent is nontrivial in its ideal, and if it equalled $\varepsilon_\pm$ then $\varepsilon_\pm a$ would be $0$, excluded). The classification of origins (isotropic iff mixed) is Proposition \ref{c:transport.existence-locus-type-n} with Example \ref{c:note.transport.hyperbolic-partner-supplies-annihilating-idempotent}.

\end{proof}

\begin{example}[$\Cl(1,1)$: N and P together]\label{c:note.transport.cl-1-1-n-p-together}
$v = \tfrac{1}{\sqrt2}(e_1 + e_2)$, $w = \tfrac{1}{2\sqrt2}(e_1 - e_2)$ form a hyperbolic pair; $e = wv = \tfrac12(1 + e_1e_2)$ is the annihilating idempotent of Example \ref{c:note.transport.hyperbolic-partner-supplies-annihilating-idempotent}, non-central; the algebra is $\Mat_2(\R)$ ($d = 0$), simple, and every zero divisor is N/P. This is the hyperbolic atom of Remark \ref{c:note.transport.hyperbolic-atom} below, the one mixed signature where the element-level choice of $w$ is canonical.
\end{example}

\begin{example}[$\Cl(1,0)$ and $\Cl(2,1)$: pure C, and C at matrix size]\label{c:note.transport.cl-1-0-cl-2-1}
$\Cl(1,0) \cong \R \oplus \R$ ($d = 1$): $\omega = e_1$, $\omega^2 = +1$, central; $\varepsilon_\pm = \tfrac12(1 \pm e_1)$; the only zero divisors are the multiples of $\varepsilon_\pm$: type C in its purest form, with no N (definite) and no P (division-algebra summands). $\Cl(2,1)$ ($d = 1$, dimension 8) $\cong \Mat_2(\R) \oplus \Mat_2(\R)$: the algebra of $2{+}1$-dimensional spacetime carries all three types at once: null vectors (mixed signature), matrix idempotents in each ideal, and the central pair $\tfrac12(1 \pm e_1e_2e_3)$.
\end{example}

\begin{example}[$\Cl(1,3)$: the physical case is N-only-at-grade-one, simple]\label{c:note.transport.cl-1-3-physical-case-n}
$d = -2 \equiv 6$: $\Cl(1,3) \cong \Mat_2(\HH)$ is simple: no type C, which is what the singularity paper's diagnosis requires: there the black-hole obstruction is type N and Rule 1 (the $(1,1)$ ascent to $\Cl(2,4)$) is the operative rule.
\end{example}

\section{Type N: the isotropic obstruction, and Rule 1}

\phantomsection\label{c:note.transport.setting-type-n-isotropic-obstruction-rule}
This is the hard case, and it takes the longest section of the paper. The plan: first the question, can the non-invertible element itself be repaired in the next ascent, and the theorem that closes it, with its cohomological sharpening (the section ``The refuted conjecture, and the no-go'', the section ``The obstruction as a cohomology class''); then the completion in which the repair succeeds at the level of points, and the totality theorem (the section ``The completion: the conformal model'', the section ``Totality of the transported reciprocal''); then the two minimality theorems that make the completion forced (the section ``Minimality among null-lift completions'', the section ``Minimality against arbitrary embeddings''); then the coordinate normal form and the cost of descent (the section ``Coordinate normal form, and descent''); and finally two consistency checks against classical mathematics (the section ``Consistency checks against known theory'').

\subsection*{The refuted conjecture, and the no-go}
\phantomsection\label{c:note.transport.setting-refuted-conjecture-no-go}
An exact inverse being impossible (Lemma \ref{c:transport.noinverse}), the standard weakening is a \emph{von Neumann pseudo-inverse}: an element $\sigma(v)$ with $v\,\sigma(v)\,v = v$. Pseudo-inverses of null vectors exist in abundance, since any null partner $w$ with $v \cdot w = \tfrac12$ supplies one by $v^2 = w^2 = 0, \qquad vw + wv = 1, \qquad \text{hence} \quad vwv = (1 - wv)v = v$, so the substantive question is not existence but \emph{canonicity}: can a pseudo-inverse be \emph{chosen}, one for each null vector, respecting the symmetries of the original algebra in the sense of Definition \ref{c:obj.transport.equivariance}? The natural conjecture, and the founding conjecture of this programme, asserts that it can, one ascent up:

\begin{refutation}[The original conjecture, refuted below]\label{c:transport.original-conjecture-refuted-below}
Division by zero promotes the element. For every null vector $v$ of $\Cl(p,q)$ there is, in the minimal ascent $\Cl(p{+}1,q{+}1)$, a canonically chosen partner $\sigma(v)$ satisfying $v\,\sigma(v)\,v = v$, the choice equivariant under the symmetries of the original algebra: the impossible division is repaired, element by element, one ascent up.
\end{refutation}

\phantomsection\label{c:note.transport.setting-refuted-conjecture-no-go-2}
(The phrase \emph{minimal ascent} anticipates Theorem \ref{c:transport.minimality}, which is logically independent of this section and proves the completion $\Cl(p{+}1,q{+}1)$ minimal and unique; the conjecture is quoted in the form in which it was originally posed.)

The conjecture is false, not merely in the minimal ascent but in every ascent however large, as soon as $\dim \mathbb{V} \geq 3$. The refutation is the first main result of this paper.

\begin{theorem}[No-go at the level of elements]\label{c:transport.no-go-level-elements}
Let $(\mathbb{V}, Q)$ be nondegenerate of mixed signature $(p,q)$, $p, q \geq 1$, with $n = p + q \geq 3$, and let $\Ncone \subset \mathbb{V}$ be its null cone. Let $(\mathbb{V}', Q')$ be \emph{any} extension of $(\mathbb{V}, Q)$ in the sense of Definition \ref{c:obj.transport.extension-quadratic-spaces}, possibly $\mathbb{V}$ itself. Then there is no map $\sigma : \Ncone \to \Cl(\mathbb{V}')$ satisfying $v\,\sigma(v)\,v = v$ for every null $v$ that is equivariant, in the sense of Definition \ref{c:obj.transport.equivariance}, even under the identity component $\Spin^0(p,q)$; a fortiori none equivariant under $\Spin(p,q)$ or $\Pin(p,q)$. Equivariant division by zero is impossible at the level of elements, in the Clifford algebra of every ascent, however large.
\end{theorem}
\begin{proof}
Fix a null $v \in \Ncone$ and a null partner $w_0$ with $v \cdot w_0 = \tfrac12$, by the hyperbolic-pair construction of the section ``Null vectors, zero divisors, and hyperbolic planes'' (\ref{c:std-algebra.witt-decomposition}); write $H_0 = \mathrm{span}\{v, w_0\}$ and $U = \{v, w_0\}^{\perp} \cap \mathbb{V}$, a nondegenerate subspace of $\mathbb{V}$ of dimension $n - 2 \geq 1$. By Definition \ref{c:obj.transport.extension-quadratic-spaces} the extension splits as $\mathbb{V}' = \mathbb{V} \perp \mathbb{V}^{\perp}$, so
\[
\mathbb{V}' \;=\; H_0 \perp U', \qquad U' \mathrel{:=} U \perp \mathbb{V}^{\perp},
\]
and Lemma \ref{c:transport.splitting-along-hyperbolic-plane} applies: every candidate $\sigma = \sigma(v) \in \Cl(\mathbb{V}')$ has a unique expansion $\sigma = \alpha + v\beta + w_0\gamma + v w_0 \delta$ with $\alpha, \beta, \gamma, \delta \in \Cl(U')$. Throughout, $\hat{\cdot}$ is the grade involution and we use the commutation rule $v\,x = \hat{x}\,v, \qquad x\,v = v\,\hat{x}$ (Lemma \ref{c:primer.commutation-rule}) for $v$ against $\Cl(U')$, together with the hyperbolic-pair identities $v^2 = w^2 = 0, \qquad vw + wv = 1, \qquad \text{hence} \quad vwv = (1 - wv)v = v$.

\emph{Step 1 (normal form for pseudo-inverses).} Compute $v\sigma v$ term by term:
\[
v\,\alpha\, v = v\, (\alpha v) = v\,(v \hat{\alpha}) = v^2 \hat\alpha = 0, \qquad
v\,(v\beta)\,v = v^2 \beta v = 0, \qquad
v\,(v w_0 \delta)\,v = v^2 w_0 \delta v = 0,
\]
while the third term survives:
\[
v\,(w_0 \gamma)\,v \;=\; v w_0\,(\gamma v) \;=\; v w_0\, v \hat{\gamma} \;=\; v \hat{\gamma}
\]
by $v^2 = w^2 = 0, \qquad vw + wv = 1, \qquad \text{hence} \quad vwv = (1 - wv)v = v$. Hence $v \sigma v = v \hat{\gamma}$, and since left multiplication by $v$ is injective on $\Cl(U')$ (Lemma \ref{c:transport.splitting-along-hyperbolic-plane}(2)) and $\hat{\cdot}$ is an involution, $v \sigma v = v$ holds if and only if $\gamma = 1$. Every pseudo-inverse of $v$ therefore has the form
\[
\sigma \;=\; w_0 + \alpha + v\beta + v w_0 \delta .
\]
Note that only the first von Neumann identity is used; the theorem does not even require $\sigma v \sigma = \sigma$.

\emph{Step 2 (null rotations).} For $u \in U$ set $R = 1 + vu$. Since $v \cdot u = 0$ and $v^2 = 0$,
\[
(vu)^2 = v\,(uv)\,u = -v\,(vu)\,u = -v^2 u^2 = 0, \qquad \text{and in particular} \quad vuv = -v^2u = 0,
\]
so $R = \exp(vu) = \exp(v \wedge u)$, with $R^{-1} = 1 - vu$. $R$ is even, and its reverse is $\widetilde{R} = 1 + uv = 1 - vu = R^{-1}$, so $R \widetilde{R} = 1$ and $R$ is a rotor.

Take from now on $u$ with $Q(u) \neq 0$, which is possible because $U$ is nondegenerate and nonzero; such $u$ span $U$. Then $R$ is a versor outright, not merely a rotor. Indeed the vector $u' \mathrel{:=} u + Q(u)\,v$ satisfies $Q(u') = Q(u) + 2Q(u)\,(u \cdot v) = Q(u) \neq 0$ and
\begin{equation}\label{eq:versorfactor}
u'u \;=\; u^2 + Q(u)\,vu \;=\; Q(u)\,(1 + vu), \qquad \text{so} \qquad R \;=\; \frac{u'\,u}{Q(u)},
\end{equation}
a product of two invertible vectors of equal square. Normalising both factors to unit square multiplies the product by $|Q(u)|/Q(u) = \pm 1$, and when that sign is $-1$ we have $Q(u) < 0$, so the normalised $\hat{u} = u/\sqrt{|Q(u)|}$ satisfies $\hat{u}^2 = -1$ and the sign is itself a product of unit vectors. Hence $R \in \Spin(p,q)$. The same factorisation applies with $u$ replaced by $tu$ for every $t \neq 0$, so $t \mapsto 1 + t\,vu$ is a continuous path \emph{inside the group} joining $1$ to $R$, and therefore $R \in \Spin^0(p,q)$.

Its action on vectors is computed directly. On $v$:
\[
R\,v\,R^{-1} = (1 + vu)\,v\,(1 - vu) = (v + vuv)(1 - vu) = v(1 - vu) = v - v^2u = v ,
\]
so $R$ fixes $v$. On $u'' \in U'$, using $u''v = -vu''$, $u u'' v = v u u''$ (two anticommutations) and $vuv = 0$:
\[
R\,u''\,R^{-1} = u'' - u''vu + vuu'' - vuu''vu = u'' + v(u''u + uu'') - v(vuu'')u = u'' + 2(u \cdot u'')\,v .
\]
On $w_0$, using $w_0 v = 1 - v w_0$ and $u w_0 = -w_0 u$:
\[
vuw_0 - w_0vu = -vw_0u - (1 - vw_0)u = -u, \qquad
vuw_0vu = vu(1 - vw_0)u = vu^2 - (vuv)w_0u = Q(u)\,v,
\]
whence
\[
R\,w_0\,R^{-1} = w_0 + (vuw_0 - w_0vu) - vuw_0vu = w_0 - u - Q(u)\,v .
\]
These are the \emph{null rotations}, the parabolic transvections stabilising the null direction $v$.

\emph{Step 3 (the shear).} Conjugate the normal form of Step 1 by $R$ and collect the $\Cl(U')$-component, i.e.\ the coefficient of $1 \in \Cl(H_0)$ in the unique expansion of Lemma \ref{c:transport.splitting-along-hyperbolic-plane}. Term by term, using $\alpha v = v\hat\alpha$, $\delta v = v \hat\delta$, $\beta v = v\hat\beta$ and $vuv = 0$:
\begin{align*}
R\,\alpha\,R^{-1} &= \alpha + vu\alpha - \alpha vu - vu\alpha vu
= \alpha + v\,(u\alpha - \hat{\alpha}u) && (\alpha vu = v\hat\alpha u;\ vu\alpha v = vuv\hat\alpha = 0),\\
R\,(v\beta)\,R^{-1} &= v\beta + (vuv)\beta - (v\beta v)u - (vuv)\beta vu = v\beta && (v\beta v = v^2\hat\beta = 0),\\
R\,(v w_0 \delta)\,R^{-1} &= v w_0 \delta - v w_0 \delta v u = v w_0 \delta - v\hat{\delta}u && (v w_0 \delta v u = v w_0 v \hat{\delta} u = v \hat\delta u),\\
R\,w_0\,R^{-1} &= w_0 - u - Q(u)\,v && \text{(Step 2)}.
\end{align*}
Every correction except one retains a left factor of $v$ and therefore lies in the $v$- or $vw_0$-slots of the expansion; the sole contribution to the $\Cl(U')$-component comes from $R w_0 R^{-1}$, and that component of $R\,\sigma(v)\,R^{-1}$ equals $\alpha - u$. But $R$ fixes $v$, so equivariance demands $\tau_R(\sigma(v)) = R\,\sigma(v)\,R^{-1} = \sigma(v)$; comparing $\Cl(U')$-components in the unique expansion gives $\alpha - u = \alpha$, i.e.\ $u = 0$, for \emph{every} $u \in U$ with $Q(u) \neq 0$. Since $U$ is nondegenerate of dimension $n - 2 \geq 1$, such $u$ exist, a contradiction.

\end{proof}

\begin{remark}[Locality and robustness of the obstruction]\label{c:note.transport.locality-robustness-obstruction}
The proof consumes far less than the theorem's hypotheses offer, and it is worth recording what it actually uses, since the surplus rules out the obvious weakenings.

\emph{It is local.} Only the stabiliser of a \emph{single} null vector enters. The sharp statement proved above is therefore: \emph{for one null vector $v$ there is no $\sigma \in \Cl(\mathbb{V}')$ with $v\sigma v = v$ that is invariant under the null rotations fixing $v$.} No global rule, no coherence across the cone, and no regularity of $\sigma$ (continuity, measurability, genericity) is assumed or would help; the theorem as stated follows at once.

\emph{It survives equivariance up to scale.} The natural first weakening is to ask only that $\sigma(v)$ be canonical projectively, $\tau_R(\sigma(v)) = \lambda\,\sigma(v)$ for some $\lambda \neq 0$. This fails too. In the normal form of Step 1 the $w_0$-coefficient is $\gamma = 1$, and Step 3 shows it is unchanged by conjugation; comparing $w_0$-components forces $\lambda = 1$, whereupon the $\Cl(U')$-comparison gives the same contradiction $u = 0$.
\end{remark}

\begin{remark}[The obstruction is parabolic]\label{c:note.transport.obstruction-parabolic}
The stabiliser of a null direction in an indefinite orthogonal group is not \emph{reductive}: it contains a nontrivial normal subgroup of \emph{unipotent} elements, transformations $g$ for which $g - \mathrm{id}$ is nilpotent (here $(R-1)^2 = 0$ in the algebra, while on $\mathbb{V}$ the action satisfies $(\Ad_R - \mathrm{id})^3 = 0$, with $(\Ad_R - \mathrm{id})^2 w_0 = -2Q(u)\,v$ nonzero whenever $Q(u) \neq 0$), called its unipotent radical. (``Parabolic'' parallels the M\"obius classification: like a parabolic M\"obius map, a null rotation is unipotent and fixes the null direction $[v]$; for $n = 3$ it fixes no other direction at all, while for $n \geq 4$ it fixes in addition the subspace $U \cap u^{\perp}$, which may itself contain null directions.) Step 3 shows that this radical acts on the candidate partners by translating the $\Cl(U')$-coefficient $\alpha$ by $-u$, freely and without fixed point; Proposition \ref{c:transport.pseudo-inverse-torsor} makes that action exact. Two consequences deserve emphasis. First, enlarging the algebra can never help: the shearing rotors $R = 1 + vu$ and the translation vectors $u$ lie inside $\Cl(\mathbb{V})$ and act the same way in every extension, which is why Theorem \ref{c:transport.no-go-level-elements} quantifies over all ascents at once. The quantifier ranges over the Clifford algebra of an extension quadratic space (Definition \ref{c:obj.transport.extension-quadratic-spaces}), which is what the splitting of Lemma \ref{c:transport.splitting-along-hyperbolic-plane} requires; it extends verbatim to matrix algebras $M_k\bigl(\Cl(\mathbb{V}')\bigr)$ over such an algebra, since reading $v\sigma v = v$ entrywise, with $v$ embedded as $v\,I$, makes each diagonal entry of $\sigma$ an invariant pseudo-inverse in $\Cl(\mathbb{V}')$ and so contradicts the theorem. Second, the phenomenon is absent in definite signature, where there are no null vectors and stabilisers of directions are compact and reductive: it is a genuinely mixed-signature obstruction, present exactly where the null cone is.
\end{remark}

\begin{remark}[The hyperbolic atom]\label{c:note.transport.hyperbolic-atom}
The assumption $n \geq 3$ is necessary, and the bound is sharp. For $n = 2$ the space is itself a hyperbolic plane, $U = 0$, there are no null rotations, and the assignment $\sigma(\lambda v) = \lambda^{-1} w_0$, $\sigma(\lambda w_0) = \lambda^{-1} v$ defines a genuinely $\Pin(1,1)$-equivariant pseudo-inverse section in the sense of Definition \ref{c:obj.transport.equivariance}, and no ascent is needed: the hyperbolic plane is self-complete \emph{at the level of elements}. (Its reciprocal on \emph{points} is of course still partial; Theorem \ref{c:transport.totality} applies to it like any other base, and its own transport is $\Cl(2,2)$. The two senses of completion are exactly the two halves of the transport principle, and the atom separates them.)

The section has a coordinate-free description, which is the real reason for its equivariance and simultaneously explains the failure for $n \geq 3$. For null $x$, call $y$ a \emph{partner} of $x$ if $y$ is null with $x \cdot y = \tfrac12$. Writing $y = av + bw_0 + u$ with $u \in U$, the two conditions read $b = 1$ and $a = -Q(u)$, so the partners of $v$ are exactly
\begin{equation}
y_u \;=\; w_0 - u - Q(u)\,v \qquad (u \in U),
\end{equation}
an affine space of dimension $n - 2$ on which, by Step 2 of Theorem \ref{c:transport.no-go-level-elements}, the null rotations act simply transitively. For $n = 2$ it is a single point, so ``the partner of $x$'' is a well-formed metric description and $\sigma$ is equivariant for that reason alone, with no computation; for $n \geq 3$ a torsor has no distinguished point, and Theorem \ref{c:transport.no-go-level-elements} says no enlargement of the space produces one. This is precisely why the hyperbolic plane is the unit of transport: the ascent of the section ``The completion: the conformal model'' adjoins a copy of the one quadratic space that needs no ascent of the first kind.


The same dimension count is Zeeman's. By Zeeman's theorem \cite{Zeeman1964}, for $n \geq 3$ every bijection of Minkowski space preserving the causal order is generated by Lorentz transformations, translations and dilations, while in $n = 2$ the two null coordinates can be reparametrised independently and the causal automorphism group is infinite-dimensional. The two-dimensional counterexamples to Zeeman are exactly the freedom that the null rotations remove for $n \geq 3$ by mixing each null direction with the transverse ones; the little group of a null direction has a null-rotation part of dimension $n - 2$, and both theorems hold precisely when it is present. Zeeman's theorem is imported here, not proved, and nothing below depends on it; this paragraph records that the third hypothesis of Theorem \ref{c:transport.no-go-level-elements} and Zeeman's are the same hypothesis, the existence of a null rotation, so that the hyperbolic atom is the exception to both for one reason.
\end{remark}

\phantomsection\label{c:note.transport.setting-refuted-conjecture-no-go-3}
The refutation redirects rather than terminates: what cannot be repaired is the element; what can be, as the remainder of this paper shows, is the \emph{space of points} in which the element's division is asked. The revised statement, no longer a conjecture, is proved as Theorems \ref{c:transport.totality} and \ref{c:transport.minimality} below: division by zero promotes not the element but the space the element belongs to, and does so uniquely, minimally, and totally.

\subsection*{The obstruction as a cohomology class}
\phantomsection\label{c:note.transport.setting-obstruction-cohomology-class}
Remark \ref{c:note.transport.locality-robustness-obstruction} observed that the no-go is local: only the stabiliser of a single null vector is used. That locality has an exact formulation, which we now give. It identifies the obstruction as a cohomology class, makes the enlargement-independence of Theorem \ref{c:transport.no-go-level-elements} structural rather than a quantifier, and exhibits the hyperbolic atom as the class's vanishing locus.

We keep the setting of the proof of Theorem \ref{c:transport.no-go-level-elements}, relaxed by one dimension so that the vanishing case is included rather than excluded: $(\mathbb{V}, Q)$ nondegenerate of mixed signature with $n \geq 2$; a null vector $v$ with hyperbolic partner $w_0$ ($v \cdot w_0 = \tfrac12$); $U = \{v, w_0\}^\perp \cap \mathbb{V}$, of dimension $n - 2 \geq 0$; an arbitrary extension $\mathbb{V}' \supseteq \mathbb{V}$ of quadratic spaces, split as $\mathbb{V}' = H_0 \perp U'$ with $U' = U \perp \mathbb{V}^\perp$; and the normal form of pseudo-inverses established in Step 1 there: the set
\[
S(v) \;=\; \{\, \sigma \in \Cl(\mathbb{V}') : v\sigma v = v \,\}
\;=\; \{\, w_0 + \alpha + v\beta + vw_0\delta \;:\; \alpha, \beta, \delta \in \Cl(U') \,\}.
\]

\begin{lemma}[The group of null rotations]\label{c:transport.group-null-rotations}
\changedmark{gray}{reworded}
Let $v$ be a null vector with hyperbolic partner $w$ and $U = \{v, w\}^{\perp}$ the nondegenerate complement of the plane they span. For $u \in U$ set $R_u = 1 + vu$. Then $R_u R_{u'} = R_{u + u'}$, so
\[
N_v \;=\; \{\, R_u : u \in U \,\}
\]
is an abelian subgroup of $\Spin^0(p,q)$, and $u \mapsto R_u$ is an isomorphism of groups $(U, +) \cong N_v$. Each $R_u$ fixes $v$. Moreover the assignment is equivariant for the Levi factor: every rotor $g$ in the even subalgebra of $\Cl(U)$ fixes $v$ and satisfies $g\, R_u\, g^{-1} = R_{g u g^{-1}}$, the conjugation on the right being the vector action of $\Spin(U)$ on $U$.
\end{lemma}
\begin{proof}
$R_u R_{u'} = 1 + v(u + u') + vuvu'$, and $vuvu' = (vuv)u' = 0$ since $vuv = -v^2u = 0$; so the group law holds, with $R_0 = 1$ and $R_u^{-1} = R_{-u}$. Injectivity of $u \mapsto R_u$: $vu = v \wedge u \neq 0$ for $u \neq 0$ by the blade basis (\ref{c:std-algebra.blade-basis}). Membership in $\Spin^0$ and the fixing of $v$ are Step 2 of the proof of Theorem \ref{c:transport.no-go-level-elements}. For the Levi statement: $g$ is an even product of vectors of $U$, each anticommuting with $v$ (\ref{c:std-algebra.orthogonal-anticommute}), so $g$ commutes with $v$ and $gvg^{-1} = v$; then $g R_u g^{-1} = 1 + (gvg^{-1})(gug^{-1}) = 1 + v\,(gug^{-1})$, and $gug^{-1} \in U$ is the rotor action.

\end{proof}

\begin{proposition}[The pseudo-inverse torsor]\label{c:transport.pseudo-inverse-torsor}
On $S(v)$ define $\sigma \sim \sigma'$ iff $\sigma - \sigma' \in v\,\Cl(\mathbb{V}')$, the left ideal generated by $v$. Then:
\begin{enumerate}
\item Each class contains a unique representative of the form $w_0 + \alpha$ with $\alpha \in \Cl(U')$, so the quotient $T(v) := S(v)/\!\sim$ is an affine space with translation space $\Cl(U')$, via $[w_0 + \alpha] + \xi = [w_0 + \alpha + \xi]$.
\item $N_v$ preserves $S(v)$ and the equivalence, hence acts on $T(v)$; the action is by translations along the grade-one subspace $U \subseteq \Cl(U')$:
\[
R_u \cdot [\,w_0 + \alpha\,] \;=\; [\,w_0 + \alpha - u\,].
\]
\item The $N_v$-orbit of any class is a torsor under $(U, +)$, and its translation cocycle, with respect to any base point, is the tautological homomorphism $\theta : N_v \to U$, $\theta(R_u) = u$.
\end{enumerate}
\end{proposition}
\begin{proof}
(1) In the normal form, $v\beta + vw_0\delta \in v\,\Cl(\mathbb{V}')$, and two elements $w_0 + \alpha$, $w_0 + \alpha'$ are equivalent iff $\alpha - \alpha' \in v\,\Cl(\mathbb{V}') \cap \Cl(U') = 0$: by the unique-expansion lemma (Lemma \ref{c:transport.splitting-along-hyperbolic-plane}), $v\,\Cl(\mathbb{V}') = v\,\Cl(U') \oplus vw_0\,\Cl(U')$ occupies exactly the $v$- and $vw_0$-slots, which meet the $\Cl(U')$-slot in $0$. (2) Preservation of $S(v)$: since $R_u$ fixes $v$, $v\,(R\sigma R^{-1})\,v = R\,(v\sigma v)\,R^{-1} = RvR^{-1} = v$. Preservation of the ideal: $R\,(vX)\,R^{-1} = (RvR^{-1})(RXR^{-1}) = v\,(RXR^{-1}) \in v\Cl(\mathbb{V}')$. The translation formula is Step 3 of the proof of Theorem \ref{c:transport.no-go-level-elements}: conjugating the normal form by $R_u$, every correction retains a left factor of $v$ except the contribution of $R_u w_0 R_u^{-1} = w_0 - u - Q(u)v$, whose $\Cl(U')$-component is $-u$. (3) The orbit map $u \mapsto R_u \cdot [\sigma_0] = [\sigma_0] - u$ is a bijection of $U$ onto the orbit (injectivity by (1)); the cocycle at base $[\sigma_0]$ is $c(R_u) = [\sigma_0] - R_u[\sigma_0] = u = \theta(R_u)$, independent of the base point since the coefficient action is trivial, and $c(R_uR_{u'}) = u + u' = c(R_u) + c(R_{u'})$ is the homomorphism property. The situation is the standard one: an affine action whose linear part is trivial has a fixed point exactly when its translation cocycle is a coboundary, and here the cocycle is tautological, the radical translating the candidates by exactly itself.

\end{proof}

\begin{theorem}[Cohomological form of the no-go]\label{c:transport.cohomological-form-no-go}
With trivial coefficients, $H^1(N_v;\, U) = \mathrm{Hom}(N_v, U)$, one-coboundaries vanishing; and the class of the pseudo-inverse torsor is the tautological homomorphism $\theta$, which is nonzero whenever $U \neq 0$. An $N_v$-invariant pseudo-inverse of $v$ exists, in the given extension, if and only if this class vanishes, if and only if $U = 0$, i.e.\ $n = 2$: the hyperbolic atom of Remark \ref{c:note.transport.hyperbolic-atom} is recovered as the exact vanishing locus of the class. Three consequences:
\begin{enumerate}
\item \emph{Enlargement-independence, cohomologically.} The group $N_v$, the coefficient module $U$, and the class $\theta$ are all constructed inside $\Cl(\mathbb{V})$ and act identically in every extension; the obstruction is pulled back unchanged along every ascent, which is the cohomological restatement of the quantification over all ascents in Theorem \ref{c:transport.no-go-level-elements}.
\item \emph{Levi equivariance.} By Lemma \ref{c:transport.group-null-rotations}, $\theta$ intertwines the conjugation action of $\Spin(U)$ on $N_v$ with its vector action on $U$: the class is equivariant for the Levi factor of the stabiliser.
\item \emph{The full stabiliser.} A pseudo-inverse selection equivariant under any subgroup of the stabiliser of $v$ containing $N_v$, in particular under the full parabolic stabiliser of the null direction, would restrict to an $N_v$-fixed point; nonvanishing of $\theta$ on the radical therefore obstructs every enlargement of the symmetry group at once.
\end{enumerate}
\end{theorem}
\begin{proof}
For a trivial module, cocycles $c(gg') = c(g) + c(g')$ are exactly homomorphisms and coboundaries $g \mapsto b - b$ vanish, so $H^1 = \mathrm{Hom}(N_v, U)$; under the isomorphism of Lemma \ref{c:transport.group-null-rotations} the class $\theta$ is the identity of $U$, nonzero iff $U \neq 0$. If $\sigma$ is $N_v$-fixed then so is its class, and by Proposition \ref{c:transport.pseudo-inverse-torsor}(2), $[\sigma] - u = [\sigma]$ forces $u = 0$ for every $u \in U$; conversely for $U = 0$ the group $N_v$ is trivial and the section of Remark \ref{c:note.transport.hyperbolic-atom} is invariant. (1)--(3) are read off from Lemma \ref{c:transport.group-null-rotations} and Proposition \ref{c:transport.pseudo-inverse-torsor} as indicated.

\end{proof}

\subsection*{The completion: the conformal model}
\phantomsection\label{c:note.transport.setting-completion-conformal-model}
This section is standard \cite{Hestenes2001, HLR2001, DFM2007, DoranLasenby2003}; we fix conventions and, for later use, the topology. Adjoin to $\mathbb{V}$ two orthogonal directions $e_+, e_-$ with $e_+^2 = 1$, $e_-^2 = -1$, both orthogonal to $\mathbb{V}$, giving $\Cl(p{+}1, q{+}1)$. Define the null pair
\[
\nO = \tfrac12(e_- - e_+), \qquad \nInf = e_- + e_+,
\]
so that $\nO^2 = \nInf^2 = 0$ and $\nO \cdot \nInf = -1$; the plane $H = \mathrm{span}\{\nO, \nInf\}$ they span is hyperbolic (the section ``Null vectors, zero divisors, and hyperbolic planes''). Represent each point $x \in \mathbb{V}$ by
\begin{equation}
X(x) \;=\; \nO + x + \tfrac12 Q(x)\, \nInf .
\end{equation}
Direct computation gives the two standard identities
\begin{equation}
X(x)^2 = 0, \qquad X(x)\cdot X(y) = -\tfrac12\, Q(x-y),
\end{equation}
so every point is represented by a null vector, and the representation is metrically faithful. Representatives are regarded projectively: a \emph{point} of the model is a null ray $[\,X\,] = \{\lambda X : \lambda \in \R^\times\}$ on the quadric
\[
\mathcal{Q} \;=\; \{[N] : N \in \mathbb{V} \oplus H,\ N \neq 0,\ N^2 = 0\}
\]
of the ambient space. Writing a general vector as $N = s\,\nO + x + t\,\nInf$ ($s, t \in \R$, $x \in \mathbb{V}$), one has $N \cdot \nInf = -s$; the rays with $s = 0$, an invariant condition since it reads $N \cdot \nInf = 0$, are not of the form $[X(x)]$ for any $x$, and we call them \emph{ideal points}. As is well known, the versor group of $\Cl(p{+}1,q{+}1)$ acting by the relation displayed with \ref{c:transport.noinverse} realises on $\mathcal{Q}$ the full conformal group of $(\mathbb{V}, Q)$, the group of its angle-preserving transformations, with inversions realised by reflections \cite{HLR2001, Angles2008}.

\phantomsection\label{c:note.transport.setting-completion-conformal-model-2}
\textbf{Topology.} Since the section ``Totality of the transported reciprocal'' makes claims about limits and continuity, we fix the topology once. Give $\mathbb{V} \oplus H \cong \R^{n+2}$ its standard topology, the projective space $P(\mathbb{V} \oplus H) = (\mathbb{V} \oplus H \setminus \{0\})/\R^{\times}$ the quotient topology, and $\mathcal{Q} \subset P(\mathbb{V} \oplus H)$ the subspace topology. Real projective space is compact and Hausdorff; the condition $N^2 = 0$ is well defined on rays (it scales by $\lambda^2$) and cuts out a closed subset; hence $\mathcal{Q}$ is compact and Hausdorff. The quotient map $\pi : N \mapsto [N]$ is continuous on nonzero vectors, and the lift $x \mapsto X(x)$ is polynomial, so $x \mapsto [X(x)]$ is continuous $\mathbb{V} \to \mathcal{Q}$. (The read-out map of Lemma \ref{c:transport.descent-read-out-map}, which runs in the opposite direction, is written $\dsc$ and should not be confused with $\pi$.)

\begin{lemma}[Density]\label{c:transport.density}
The lifted invertible points $\{[X(x)] : x \in \mathbb{V},\ Q(x) \neq 0\}$ are dense in $\mathcal{Q}$.
\end{lemma}
\begin{proof}
First, the lifted points exhaust the non-ideal rays: if $N = s\nO + x + t\nInf$ is null with $s \neq 0$, then rescaling by $s^{-1}$ puts it in the form $\nO + x' + t'\nInf$ with $Q(x') = 2 \cdot 1 \cdot t'$ forced by nullity (see \ref{c:transport.swap-rescale-normal-form} below), i.e.\ $N \in [X(x')]$. Second, the invertible points are dense in $\mathbb{V}$: for any $x$, pick $a$ with $Q(a) \neq 0$; then $t \mapsto Q(x + ta) = Q(x) + 2t\,(x \cdot a) + t^2 Q(a)$ is a nonconstant polynomial, hence nonzero for all small $t \neq 0$, and $x + ta \to x$. By continuity of the lift, the lifted invertible points are dense in the set of \emph{all} lifted points. It remains to show every ideal point is a limit of lifted invertible points; we exhibit explicit paths.

\emph{Ideal points with $x \neq 0$.} Let $[\,x + t\,\nInf\,]$ be ideal, so $Q(x) = 0$, $x \neq 0$, $t \in \R$. By nondegeneracy choose $u \in \mathbb{V}$ with $x \cdot u = 1$, and for $s > 0$ set
\[
y_s \;=\; s\,x + \bigl(t + \tfrac{1}{s}\bigr)\,u .
\]
Then $Q(y_s) = 2s\bigl(t + \tfrac1s\bigr) + \bigl(t + \tfrac1s\bigr)^2 Q(u) = 2ts + 2 + \bigl(t + \tfrac1s\bigr)^2 Q(u)$, which is nonzero for all large $s$: if $t \neq 0$ the term $2ts$ dominates the remaining bounded terms, and if $t = 0$ the expression tends to $2$. So $y_s$ is eventually invertible. The rescaled representative
\[
\tfrac{1}{s}\,X(y_s) \;=\; \tfrac{1}{s}\,\nO + x + \tfrac{1}{s}\bigl(t + \tfrac1s\bigr)u + \tfrac{Q(y_s)}{2s}\,\nInf
\;\longrightarrow\; x + t\,\nInf \quad (s \to \infty),
\]
since $\tfrac{Q(y_s)}{2s} = t + \tfrac1s + O(1/s) \to t$. The limit vector is nonzero, so by continuity of $\pi$, $[X(y_s)] \to [\,x + t\,\nInf\,]$.

\emph{The ideal point $[\nInf]$.} Pick any $a$ with $Q(a) \neq 0$ and set $y_s = s\,a$. Then $\tfrac{2}{s^2 Q(a)} X(y_s) = \tfrac{2}{s^2 Q(a)} \nO + \tfrac{2}{s\,Q(a)}a + \nInf \to \nInf$.

\end{proof}

\phantomsection\label{c:note.transport.setting-completion-conformal-model-3}
Note the division of labour between the two transports of the ascent. The inclusion $\Cl(p,q) \subseteq \Cl(p{+}1,q{+}1)$ transports \emph{elements}, and fixes nothing: a null $v$ remains null and, by Theorem \ref{c:transport.no-go-level-elements}, remains without equivariant repair in the enlarged algebra. The map $X$ transports \emph{points}: it is quadratic, preserves no products, and is not an algebra map; and it is on points that the repair succeeds.

\subsection*{Totality of the transported reciprocal}
\phantomsection\label{c:note.transport.setting-totality-transported-reciprocal}
Consider the single-vector versor $e_+$, acting by $N \mapsto e_+ N e_+$ (an overall sign being projectively immaterial; note $e_+^{-1} = e_+$ and $e_+ = -(\nO - \tfrac12\nInf)$). On the generators of the construction:
\begin{equation}
e_+ \nO e_+ = -\tfrac12 \nInf, \qquad e_+ \nInf e_+ = -2\nO, \qquad e_+ x e_+ = -x \quad (x \in \mathbb{V}),
\end{equation}
each line a one-step consequence of anticommutation. The reflection exchanges the origin with the point at infinity; this is the model's realisation of inversion \cite{HLR2001}.

\begin{theorem}[Totality]\label{c:transport.totality}
Let $(\mathbb{V}, Q)$ be nondegenerate, with the completion $\Cl(p{+}1,q{+}1)$ and quadric $\mathcal{Q}$ of the section ``The completion: the conformal model''. Define $\iota : \mathcal{Q} \to \mathcal{Q}$ by $\iota[N] = [\,e_+ N e_+\,]$. Then:
\begin{enumerate}
\item $\iota$ is well defined, total and continuous on $\mathcal{Q}$, and an involution.
\item For every $x$ with $Q(x) \neq 0$, \; $\iota[X(x)] = [X(x/Q(x))]$: on invertible vectors $\iota$ is the geometric reciprocal $\rho$.
\item For every $x_0$ with $Q(x_0) = 0$, that is, a null vector or the origin, $\iota[X(x_0)] = [\, x_0 + \tfrac12 \nInf \,]$, an ideal point (for $x_0 = 0$, the classical point at infinity $[\nInf]$), and this value is the limit of $\iota[X(x)]$ along every path of invertible $x \to x_0$. Consequently $\iota$ is the unique continuous map $\mathcal{Q} \to \mathcal{Q}$ agreeing with the lifted reciprocal $[X(x)] \mapsto [X(\rho(x))]$ on the lifted invertible locus.
\end{enumerate}
\end{theorem}
\begin{proof}
(1) Conjugation by the invertible element $e_+$ is an algebra automorphism of $\Cl(p{+}1,q{+}1)$ (on vectors, as $e_+^{2} = 1$, it is minus the reflection in $e_+^{\perp}$, \ref{c:std-algebra.reflection}); in particular it is injective, so it sends nonzero elements to nonzero elements. By $e_+ \nO e_+ = -\tfrac12 \nInf, \qquad e_+ \nInf e_+ = -2\nO, \qquad e_+ x e_+ = -x \quad (x \in \mathbb{V})$ and linearity it maps vectors of the ambient space to vectors, and it preserves squared lengths: $(e_+ N e_+)^2 = e_+ N^2 e_+ = N^2$ for any vector $N$, since $N^2$ is a scalar. Hence it maps nonzero null vectors to nonzero null vectors, and rays to rays ($\lambda N \mapsto \lambda\, e_+Ne_+$), so $\iota$ is well defined and total on $\mathcal{Q}$. It is continuous because it descends from an invertible linear map of $\mathbb{V} \oplus H$ through the quotient $\pi$. Finally $e_+ (e_+ N e_+) e_+ = N$ since $e_+^2 = 1$, so $\iota^2 = \mathrm{id}$.

(2) By linearity of conjugation and $e_+ \nO e_+ = -\tfrac12 \nInf, \qquad e_+ \nInf e_+ = -2\nO, \qquad e_+ x e_+ = -x \quad (x \in \mathbb{V})$,
\begin{equation}\label{eq:keydisplay}
e_+ X(x)\, e_+ \;=\; -\Bigl( Q(x)\,\nO + x + \tfrac12 \nInf \Bigr) \qquad \text{for every } x \in \mathbb{V}.
\end{equation}
For $Q(x) \neq 0$, rescaling by $-1/Q(x)$ gives $\nO + x/Q(x) + \tfrac12 (1/Q(x)) \nInf$, and since $Q(x/Q(x)) = 1/Q(x)$ this is exactly $X(x/Q(x))$ by $X(x) \;=\; \nO + x + \tfrac12 Q(x)\, \nInf$.

(3) For $Q(x_0) = 0$ the display \ref{c:transport.totality} gives $e_+ X(x_0) e_+ = -(x_0 + \tfrac12 \nInf)$, a nonzero vector (its $\nInf$-coefficient is $\tfrac12$) whose $\nO$-coefficient vanishes: an ideal point, null since $x_0 \cdot \nInf = 0$ and $x_0^2 = \nInf^2 = 0$; for $x_0 = 0$ it is the class $[\nInf]$. For the limit claim, consider the map $J : \mathbb{V} \to \mathcal{Q}$, $J(x) = [\,Q(x)\nO + x + \tfrac12\nInf\,]$. Its representative is polynomial in $x$ and never zero (the $\nInf$-coefficient is $\tfrac12$), so $J$ is continuous; and by \ref{c:transport.totality}, $J(x) = \iota[X(x)]$ for \emph{every} $x \in \mathbb{V}$, whether $Q(x)$ vanishes or not. Continuity of $J$ at $x_0$ now gives $\iota[X(x)] \to \iota[X(x_0)]$ along every path of invertible $x \to x_0$, as claimed. For uniqueness: the lifted invertible points are dense in $\mathcal{Q}$ (Lemma \ref{c:transport.density}) and $\mathcal{Q}$ is Hausdorff, so any two continuous maps $\mathcal{Q} \to \mathcal{Q}$ that agree on them agree everywhere; $\iota$ agrees with the lifted reciprocal there by (2).

\end{proof}

\begin{remark}\label{c:note.transport.destination}
The geometric content of (3), for a null vector $v$: the value of ``division by $v$'' is an ideal point at null infinity \emph{determined by} $v$: it lies over the null direction of $v$, at a position fixed by the scale of $v$. The value remembers more than the direction alone: rescaling $v$ moves it, $[\,2v + \tfrac12\nInf\,] \neq [\,v + \tfrac12\nInf\,]$ (in the coordinates of Proposition \ref{c:transport.swap-rescale-normal-form}, these are distinct points of the same ideal stratum). The answer to the forbidden division exists and is canonical; it merely does not lie in the original space. Together with Theorem \ref{c:transport.no-go-level-elements}, this is the precise sense in which the ascent performs ``division by zero without an inverse operation'': the element remains non-invertible; the space of points acquires a value for the division.
\end{remark}

\subsection*{Minimality among null-lift completions}
\phantomsection\label{c:note.transport.setting-minimality-among-null-lift-completions}
The conformal model uses \emph{two} extra dimensions of opposite square. The following shows this is not a design choice. We first fix the class of completions over which minimality is claimed.

\begin{definition}[Null-lift completion]\label{c:obj.transport.null-lift-completion}
A \emph{null-lift completion} of the nondegenerate quadratic space $(\mathbb{V}, Q)$ is a pair $(\mathbb{W}, f)$ where $\mathbb{W}$ is a nondegenerate real quadratic space adjoined orthogonally to $\mathbb{V}$ and $f : \mathbb{V} \to \mathbb{W}$ is a map with $f(0) \neq 0$ such that
\[
X(x) \;=\; f(x) + x \ \in\ \mathbb{W} \perp \mathbb{V}
\]
is null for every $x \in \mathbb{V}$. The condition $f(0) \neq 0$ says only that the origin, like every point, is represented by a nonzero vector, hence by a ray. Nothing else is assumed of $f$: not linearity, not continuity, not injectivity.
\end{definition}

\phantomsection\label{c:note.transport.setting-minimality-among-null-lift-completions-2}
Two hypotheses in this definition carry weight, and both are deliberate.

The shape $X(x) = f(x) + x$ is the normalisation that makes the minimality question well posed. Without some such condition ``minimal completion'' is not yet a definite notion: any quadratic space whatever admits \emph{some} map into \emph{some} larger space, and one can always use a dimension of the ambient space to re-encode the base rather than extend it. Requiring the $\mathbb{V}$-coordinates of a represented point to be the point itself says exactly that the representation leaves the base alone and puts new structure only into the adjoined directions; the question then asked, and answered, is how much new structure that is. (Allowing an arbitrary re-embedding of $\mathbb{V}$ alongside would merely relabel the base without changing which spaces $\mathbb{W}$ can host nullity, so nothing is lost by the normalisation, and the counting becomes meaningful.) The conformal lift $X(x) \;=\; \nO + x + \tfrac12 Q(x)\, \nInf$ is a null-lift completion with $\mathbb{W} = H$ and $f(x) = \nO + \tfrac12 Q(x)\nInf$, $f(0) = \nO \neq 0$.

Nondegeneracy of $\mathbb{W}$ is the second, and Remark \ref{c:note.transport.degenerate-adjunctions-why-they-are-excluded} records precisely where it is needed and what happens without it.

\begin{theorem}[Minimality]\label{c:transport.minimality}
Let $(\mathbb{W}, f)$ be any null-lift completion of $(\mathbb{V}, Q)$. Then $\mathbb{W}$ contains at least one direction of positive square and one of negative square; in particular $\dim \mathbb{W} \geq 2$, and enlargements of signature increment $(+k, 0)$ or $(0, +k)$ are excluded for every $k$. A single hyperbolic plane suffices, and any two-dimensional $\mathbb{W}$ admitted by the definition is a hyperbolic plane; hence the minimal completion is
\[
\Cl(p,q) \;\longrightarrow\; \Cl(p{+}1, q{+}1),
\]
with the adjoined space unique up to isometry, and the signature increment is $(+1,+1)$ in every signature.
\end{theorem}
\begin{proof}
Nullity of $X(x)$ and orthogonality of $\mathbb{W}$ to $\mathbb{V}$ give
\[
Q_{\mathbb{W}}(f(x)) = -Q(x) \qquad \text{for all } x \in \mathbb{V}.
\]
Setting $x = 0$: $Q_{\mathbb{W}}(f(0)) = 0$ with $f(0) \neq 0$, so $\mathbb{W}$ contains a nonzero null vector. A nondegenerate quadratic space containing a nonzero null vector is indefinite \cite{Lam2005} (a null vector has a hyperbolic partner, \ref{c:std-algebra.witt-decomposition}), hence contains directions of both signs and has dimension at least $2$; increments $(+k,0)$ and $(0,+k)$, whose forms are definite, are excluded for every $k$. (When $\mathbb{V}$ has mixed signature the same conclusion follows a second way, with no appeal to the origin: $-Q$ attains both signs, hence so does $Q_{\mathbb{W}}$.)

Sufficiency of one hyperbolic plane is the construction of the section ``The completion: the conformal model''--the section ``Totality of the transported reciprocal''. For uniqueness of the minimal $\mathbb{W}$: a two-dimensional nondegenerate quadratic space containing a nonzero null vector is indefinite, hence of signature $(1,1)$, and is the hyperbolic plane, unique up to isometry (the section ``Null vectors, zero divisors, and hyperbolic planes'', \cite{Lam2005}; \ref{c:std-algebra.sylvester-inertia}). Hence $\dim \mathbb{W} = 2$ with signature $(1,1)$ is minimal and unique, and the increment is $(+1,+1)$.

\end{proof}

\begin{remark}\label{c:note.transport.signargument}
Two small arguments carry the theorem, and they deserve emphasis for how little they assume. The sign argument: the adjoined directions must be capable of cancelling squared lengths of every sign, and only a mixed addition can: new null structure appears only where positive meets negative. The origin argument: merely representing the origin as a ray already requires a null vector of $\mathbb{W}$, which forces indefiniteness even over a Euclidean base. The obstruction thereby determines its own repair, including where the repair lies. Note also the consonance with Remark \ref{c:note.transport.hyperbolic-atom}: the completion adjoins precisely the hyperbolic plane, the one quadratic space that is self-complete at the level of elements, the unit of transport.
\end{remark}

\begin{remark}[Degenerate adjunctions, and why they are excluded]\label{c:note.transport.degenerate-adjunctions-why-they-are-excluded}
Nondegeneracy of $\mathbb{W}$ is a hypothesis of Definition \ref{c:obj.transport.null-lift-completion}, not a consequence, and it is worth being exact about its role, since the degenerate adjunction $\Cl(n,0,1)$ of projective geometric algebra is a familiar and useful construction.

Over a base of \emph{mixed} signature the hypothesis is idle: the parenthetical second argument in the proof shows $Q_{\mathbb{W}}$ attains both signs, so $\mathbb{W}$ has a positive and a negative direction whether or not it also carries a radical, and the increment $(+1,+1)$ is forced regardless.

Over a \emph{definite} base the hypothesis is essential to the proof. Take $\mathbb{V}$ positive definite and let $\mathbb{W} = \mathrm{span}\{r, b\}$ with $r$ spanning a radical ($Q_{\mathbb{W}}(r) = 0$, $r \cdot b = 0$) and $Q_{\mathbb{W}}(b) = -1$. Then $f(x) = r + \sqrt{Q(x)}\,b$ satisfies $Q_{\mathbb{W}}(f(x)) = -Q(x)$ and $f(0) = r \neq 0$: a null lift of increment $(0,+1)$ plus a degenerate direction. Theorem \ref{c:transport.minimality} excludes it by fiat, and rightly so for present purposes: everything the completion is asked to do here (represent points as rays on a quadric, realise the inversion as a reflection, and read out answers by descent) runs on the versor machinery of the section ``The symmetry action, and equivariance'', which needs invertible generators and Cartan--Dieudonn\'e, both unavailable along a radical direction. The minimality theorem should therefore be read as: among completions whose adjoined structure is usable as geometry, the increment is $(+1,+1)$, always.
\end{remark}

\begin{remark}[What is and is not unique]\label{c:note.transport.what-not-unique}
Uniqueness in Theorem \ref{c:transport.minimality} is uniqueness of the adjoined \emph{space} $\mathbb{W}$ (up to isometry), hence of the completed algebra up to isomorphism. The lift $f$ itself is not unique: for any $c \neq 0$, $f_c(x) = c\,\nO + \tfrac{1}{2c} Q(x)\,\nInf$ satisfies the same nullity condition. What is forced is the space in which the lift must lie, not how it is parametrised.
\end{remark}

\subsection*{Minimality against arbitrary embeddings}
\phantomsection\label{c:note.transport.setting-minimality-against-arbitrary-embeddings}
Theorem \ref{c:transport.minimality} quantifies over null-lift completions of the quadratic space, with the normalisation $X(x) = f(x) + x$ that makes the counting meaningful. One can ask more: could $\Cl(p,q)$, embedded \emph{exotically} inside some smaller Clifford algebra, find its resolution there? Exotic unital embeddings exist (Example \ref{c:note.transport.exotic-embedding-inertness}), so the question has content. The answer is that hosting a resolution is a property of the ambient \emph{quadric}, not of the embedding, and the following theorem settles it by a rank computation. It is the sharpest form of the minimality half of the trichotomy, and it discharges the scope question that the normalisation of Definition \ref{c:obj.transport.null-lift-completion} leaves open.

\begin{definition}[Normalised conformal chart]\label{c:obj.transport.normalised-conformal-chart}
Let $(\mathbb{V}'', Q'')$ be a nondegenerate real quadratic space. A \emph{normalised conformal chart} of $(\mathbb{V}, Q)$ in $\mathbb{V}''$ is a map $\Phi : \mathbb{V} \to \mathbb{V}''$ with
\[
\Phi(x) \cdot \Phi(y) \;=\; -\tfrac12\, Q(x - y) \qquad \text{for all } x, y \in \mathbb{V}
\]
(in particular every $\Phi(x)$ is null, and $\Phi$ is injective by nondegeneracy). This is exactly the metric-faithfulness identity $X(x)^2 = 0, \qquad X(x)\cdot X(y) = -\tfrac12\, Q(x-y)$ of the conformal model; the lift $X$ of the completion is a normalised conformal chart, and any projective chart with a nowhere-vanishing scale is rescaled into this form. Nothing else is assumed of $\Phi$: not linearity, not continuity.
\end{definition}

\begin{lemma}[The conformal Gram matrix]\label{c:transport.conformal-gram-matrix}
Fix an orthogonal basis $e_1, \dots, e_n$ of $\mathbb{V}$ and the $n + 2$ points
\[
p_0 = 0, \qquad p_i = e_i \ (1 \leq i \leq n), \qquad p_{n+1} = 2e_1 .
\]
Then the symmetric $(n{+}2) \times (n{+}2)$ matrix $M_{ij} = -\tfrac12 Q(p_i - p_j)$ has rank $n + 2$ and signature $(p{+}1,\, q{+}1)$.
\end{lemma}
\begin{proof}
Compute $M$ in the completion: by the model identity $X(x)^2 = 0, \qquad X(x)\cdot X(y) = -\tfrac12\, Q(x-y)$, $M$ is the Gram matrix of the lifted vectors $X(p_i) \in \mathbb{V} \oplus H$, so its rank and signature are those of the ambient form restricted to their span, and it suffices to show the lifts span the whole $(n{+}2)$-dimensional space, of signature $(p{+}1, q{+}1)$. The span contains $\nO = X(0)$; it contains $X(2e_1) - 2X(e_1) + X(0) = (2 - 2)e_1 + (2 - 1)Q(e_1)\nInf = Q(e_1)\,\nInf$ with $Q(e_1) = \pm 1 \neq 0$, hence $\nInf$; hence each $e_i = X(e_i) - X(0) - \tfrac12 Q(e_i)\nInf$. A spanning set of a nondegenerate space has Gram matrix congruent to the form (\ref{c:std-algebra.sylvester-inertia}), giving rank $n+2$ and the stated signature.

\end{proof}

\begin{theorem}[Representation-independent minimality of the completion]\label{c:transport.representation-independent-minimality-completion}
Let $(\mathbb{V}'', Q'')$ be nondegenerate of signature $(p'', q'')$ and suppose its null cone admits a normalised conformal chart of $(\mathbb{V}, Q)$. Then
\[
p'' \geq p + 1 \quad \text{and} \quad q'' \geq q + 1 ;
\]
in particular $\dim \mathbb{V}'' \geq n + 2$, and in case of equality the signature is exactly $(p{+}1, q{+}1)$ and $\Cl(\mathbb{V}'') \cong \Cl(p{+}1, q{+}1)$. Consequently the minimality of the completion is a property of the ambient quadric alone: no Clifford algebra whose quadratic space is smaller than, or differently balanced than, $(p{+}1, q{+}1)$ can host the point-level resolution, \emph{however} $\Cl(p,q)$ is embedded in it as an algebra.
\end{theorem}
\begin{proof}
Let $s_i = \Phi(p_i)$ for the points of Lemma \ref{c:transport.conformal-gram-matrix}. The chart condition gives $s_i \cdot s_j = M_{ij}$, so the pullback of $Q''|_{\mathrm{span}\{s_i\}}$ to $\R^{n+2}$ along $(c_i) \mapsto \sum c_i s_i$ is the quadratic form of $M$. Since $M$ is nondegenerate (Lemma \ref{c:transport.conformal-gram-matrix}), the map is injective, the span $S$ has dimension $n+2$, and $Q''|_S$ is nondegenerate of signature $(p{+}1, q{+}1)$. A nondegenerate subspace splits off orthogonally, $\mathbb{V}'' = S \perp S^\perp$, whence $p'' = (p{+}1) + p(S^\perp) \geq p+1$ and likewise for $q''$. Equality of dimensions forces $S = \mathbb{V}''$ and the signature claim; the algebra statement is the functoriality of $\Cl$ on isometries.

\end{proof}

\begin{example}[An exotic embedding, and its inertness]\label{c:note.transport.exotic-embedding-inertness}
$\Cl(1,3) \cong \Mat_2(\HH)$ embeds unitally in $\Cl(2,3) \cong \Mat_4(\C)$, an algebra of \emph{half} the dimension of the completion $\Cl(2,4)$: the standard unital embedding $\HH \hookrightarrow \Mat_2(\C)$ ($1, i, j, k$ mapped to $I$, $\mathrm{diag}(i, -i)$, $\begin{psmallmatrix}0&1\\-1&0\end{psmallmatrix}$, $\begin{psmallmatrix}0&i\\i&0\end{psmallmatrix}$) applied entrywise gives $\Mat_2(\HH) \hookrightarrow \Mat_4(\C)$, and $\Cl(2,3)$ has type $\Mat_4(\C)$ by Theorem \ref{c:transport.classification-real-clifford-algebras} ($d \equiv 7$). The embedding is genuine; it is also inert: the ambient quadratic space of $\Cl(2,3)$ is five-dimensional, and by Theorem \ref{c:transport.representation-independent-minimality-completion} no normalised conformal chart of four-dimensional spacetime exists on its quadric. The algebra is big enough to contain $\Cl(1,3)$ and too small to resolve it: what a resolution needs is not space for the \emph{elements} but space for the \emph{points}, and points cost two dimensions in the exact balance $(+1,+1)$.
\end{example}

\subsection*{Coordinate normal form, and descent}
\phantomsection\label{c:note.transport.setting-coordinate-normal-form-descent}
The mechanism of Theorem \ref{c:transport.totality} admits a coordinate statement that identifies it with a classical pattern. (The scalar slots are written $s, t$ to avoid collision with the $\Cl(U')$-coefficients $\alpha, \beta$ of the section ``The refuted conjecture, and the no-go''.)

\begin{proposition}[Swap-and-rescale normal form]\label{c:transport.swap-rescale-normal-form}
In the setting of Theorem \ref{c:transport.totality}, represent elements of the quadric by homogeneous triples $(s, x, t)$ with $s,t \in \R$ and $x \in \mathbb{V}$, standing for $N = s\,\nO + x + t\,\nInf$ and defined up to overall scale, subject to the incidence (nullity) condition
\begin{equation}
Q(x) = 2st .
\end{equation}
The ordinary points are the classes with $s \neq 0$, normal form $(1, x, \tfrac12 Q(x))$, i.e.\ $X(x) \;=\; \nO + x + \tfrac12 Q(x)\, \nInf$. The ideal points, the classes with $s = 0$, then satisfy $Q(x) = 0$ and fall into three disjoint strata:
\begin{enumerate}
\item[(i)] $t \neq 0$ and $x \neq 0$: unique representative $(0, x', \tfrac12)$ with $x' = x/(2t)$ null; these classes are labelled bijectively by the null \emph{vectors} of $\mathbb{V}$, and are exactly the values $\iota$ takes on lifted null points (Theorem \ref{c:transport.totality}(3));
\item[(ii)] $t = 0$, $x \neq 0$: the class retains only the null \emph{ray} $[x]$; these classes are labelled by the null directions of $\mathbb{V}$;
\item[(iii)] $x = 0$, which given $s = 0$ forces $t \neq 0$: the single classical point at infinity $[\nInf]$, which is the value $\iota[X(0)]$ by Theorem \ref{c:transport.totality}(3).
\end{enumerate}
In these coordinates the map $\iota$ of Theorem \ref{c:transport.totality} is the linear involution
\[
\iota(s, x, t) = (2t,\, x,\, \tfrac12 s),
\]
which preserves the incidence condition, is total, and recovers the reciprocal upon normalising the first slot where possible; where normalisation is impossible (first slot zero) the output is an ideal point.
\end{proposition}
\begin{proof}
The three displayed actions of $e_+ \nO e_+ = -\tfrac12 \nInf, \qquad e_+ \nInf e_+ = -2\nO, \qquad e_+ x e_+ = -x \quad (x \in \mathbb{V})$ are the three columns of $\iota$ in the basis $(\nO, \mathbb{V}, \nInf)$: $e_+ N e_+ = -\bigl(2t\,\nO + x + \tfrac12 s\,\nInf\bigr)$, and the overall factor $-1$ is projectively immaterial, as in the section ``Totality of the transported reciprocal''. Preservation of the condition is the check $2\,(2t)(\tfrac12 s) = 2st$. The stratification of the ideal points is read off from \ref{c:transport.swap-rescale-normal-form} with $s = 0$: nullity of $x$ is forced; if $x \neq 0$ and $t \neq 0$, scaling the class by $\tfrac{1}{2t}$ gives the unique representative with last slot $\tfrac12$, and Theorem \ref{c:transport.totality}(3) exhibits every such class, and no other, as $\iota[X(x')]$ for the null vector $x'$; if $x \neq 0$ and $t = 0$ the class is the ray $[x]$ and only the direction of $x$ survives the projective quotient; and if $x = 0$ then $t \neq 0$, since $s = x = t = 0$ is not a class, leaving $[\nInf]$.

\end{proof}

\begin{remark}[Where each stratum comes from]\label{c:note.transport.where-each-stratum-comes}
The three strata are distinguished by $\iota$ itself, which supplies a check on the bookkeeping. Stratum (i) is exactly the image of the null cone of $\mathbb{V}$ under $\iota \circ X$, and (iii) is the image of the origin. Stratum (ii) is the image of neither: it is fixed \emph{pointwise} by $\iota$, since $(0,x,0) \mapsto (0,x,0)$. Indeed the whole fixed locus of $\iota$ is $\{Q(x) = 1\} \cup \text{(ii)}$, as one sees by solving $(2t, x, \tfrac12 s) = \lambda(s,x,t)$ against \ref{c:transport.swap-rescale-normal-form}: $\lambda = 1$ forces $s = 2t$ and $Q(x) = 4t^2$, which after normalisation is the unit quadric when $t \neq 0$ and stratum (ii) when $t = 0$. So the null directions at infinity are the ideal fixed points of the inversion, as they must be: they are the directions along which the null cone through the origin goes to infinity, and inversion carries that cone to itself.
\end{remark}

\phantomsection\label{c:note.transport.setting-coordinate-normal-form-descent-2}
The normalisation just described is worth defining as a map in its own right, since it is the map by which answers return to the original space.

\begin{lemma}[Descent: the read-out map]\label{c:transport.descent-read-out-map}
Define $\dsc : \mathcal{Q} \setminus \{\text{ideal points}\} \to \mathbb{V}$ by
\[
\dsc\,[\, s\,\nO + x + t\,\nInf \,] \;=\; x/s .
\]
Then:
\begin{enumerate}
\item $\dsc$ is well defined and continuous, and it is inverse to the lift: $\dsc[X(x)] = x$ for every $x \in \mathbb{V}$, and $X(\dsc[N]) \in [N]$ for every non-ideal null ray $[N]$. Hence $x \mapsto [X(x)]$ and $\dsc$ are mutually inverse homeomorphisms between $\mathbb{V}$ and $\mathcal{Q} \setminus \{\text{ideal points}\}$.
\item On the invertible locus, descent returns the classical reciprocal: for every $x$ with $Q(x) \neq 0$,
\[
\dsc\bigl(\iota[X(x)]\bigr) \;=\; x/Q(x) \;=\; \rho(x).
\]
\item The composite $\dsc \circ \iota \circ [X(\cdot)]$ is defined exactly on $\{Q \neq 0\}$: for $Q(x_0) = 0$ the value $\iota[X(x_0)]$ is ideal by Theorem \ref{c:transport.totality}(3) and lies outside the domain of $\dsc$. Descent is therefore partial precisely where the original reciprocal was, and necessarily so, since the ideal points are represented by no point of $\mathbb{V}$.
\end{enumerate}
\end{lemma}
\begin{proof}
(1) Scaling a representative by $\lambda \neq 0$ scales $s$ and $x$ alike, so $x/s$ is unchanged and $\dsc$ is well defined; the map $N \mapsto x(N)/s(N)$ is continuous on the open, scale-invariant set $\{s \neq 0\}$ of the null cone, so it descends to a continuous map on the quotient. For the lift, $\dsc[X(x)] = x/1 = x$. Conversely, given a non-ideal $[N]$, normalise $s = 1$; the incidence condition \ref{c:transport.swap-rescale-normal-form} then forces $t = \tfrac12 Q(x)$, so the normalised representative is exactly $X(x)$ with $x = \dsc[N]$, whence $X(\dsc[N]) \in [N]$. Both directions are continuous, giving the homeomorphism.

(2) By Theorem \ref{c:transport.totality}(2), $\iota[X(x)] = [X(\rho(x))]$, and $\dsc[X(\rho(x))] = \rho(x)$ by (1).

(3) Immediate from Theorem \ref{c:transport.totality}(3) and the definition of the ideal points.

\end{proof}

\begin{remark}[Ascend, invert, descend]\label{c:note.transport.ascend-invert-descend}
In plain terms, the classical reciprocal is a composite of three steps. To divide by $x$, first \emph{ascend}: represent the point $x$ in the completion as the null ray $[X(x)]$. Next \emph{invert}: apply the reflection $\iota$, which is total and never fails. Finally \emph{descend}: return the answer to $\mathbb{V}$ with the read-out $\dsc$. Lemma \ref{c:transport.descent-read-out-map}(2) states that whenever $Q(x) \neq 0$ this composite gives exactly the classical answer $x/Q(x)$. When $Q(x) = 0$ the first two steps still succeed, and the composite fails only at the last: the answer $\iota[X(x)]$ exists, but it is an ideal point, and ideal points correspond to no point of the original space to which $\dsc$ could return them. The division in the completion succeeds; what fails is the return of its result to $\mathbb{V}$. The classical impossibility of dividing by a null vector is thereby relocated rather than removed: it is not a failure of the operation but a failure of the descent, and the transport explains the obstruction rather than removing it. The answer exists, and it lies permanently outside the original space.
\end{remark}

\begin{remark}\label{c:note.transport.classical}
This is the linearise-then-normalise pattern of classical mathematics: M\"obius transformations $z \mapsto (az+b)/(cz+d)$ rendered total as linear maps on homogeneous pairs $(z,1)$ with the unnormalisable pairs $(\cdot\,,0)$ serving as the point at infinity; the homogeneous coordinates of projective geometry; the deferred ``perspective divide'' of computer graphics. All nonlinearity and all division are deferred to a single optional read-out, whose refusal is information (the answer is ideal), not failure. What Theorems \ref{c:transport.no-go-level-elements}, \ref{c:transport.totality} and \ref{c:transport.minimality} add to the classical picture is a mandate: for the geometric reciprocal of a quadratic space, this pattern is not one available representation among many but the \emph{unique minimal} totalisation, realised on the null cone exactly two dimensions up, with the element-level alternative provably closed.
\end{remark}

\subsection*{Consistency checks against known theory}
\phantomsection\label{c:note.transport.setting-consistency-checks-against-known-theory}
Two specialisations connect the theorems to classical mathematics: in each, the constructions of this section specialise exactly to a textbook object. Neither is needed for the proofs; both are offered as verification, being places where the new results are obliged to reproduce old ones, and do.

\subsection*{The Euclidean plane: inversion in a circle on the Riemann sphere}
\phantomsection\label{c:note.transport.setting-euclidean-plane-inversion-circle-riemann}
Take the base $\mathbb{V} = \R^2$ with $Q$ positive definite (signature $(2,0)$), and identify $\mathbb{V}$ with $\C$ by $z = x + iy$, so that $Q(z) = |z|^2$. Definite signature has no null vectors, so the partial locus of $\rho(z) = z/Q(z)$ is the origin alone: this is the oldest division by zero there is. Note which hypotheses are at work. The no-go theorem requires $p, q \geq 1$ and $n \geq 3$ and does not apply here, since with no null vectors there is no element-level question to obstruct, while the theorems this example exercises (totality, minimality, the normal form) hold in every nondegenerate signature, as recorded in the scope paragraph of the section ``Quadratic spaces and Clifford algebras''.

The completion is $\Cl(3,1)$. Because $Q$ is definite, the incidence condition \ref{c:transport.swap-rescale-normal-form} with $s = 0$ forces $z = 0$: the ideal strata (i) and (ii) of Proposition \ref{c:transport.swap-rescale-normal-form} are empty, and the only ideal point is $[\nInf]$, stratum (iii). The quadric $\mathcal{Q}$ is therefore the plane with a single point adjoined, and the map
\[
[X(z)] \;\mapsto\; \frac{1}{|z|^2+1}\bigl(\,2x,\; 2y,\; |z|^2 - 1\,\bigr), \qquad [\nInf] \;\mapsto\; (0,0,1)
\]
identifies $\mathcal{Q}$ with the unit sphere $S^2 \subset \R^3$ by inverse stereographic projection from the north pole. This is the Riemann sphere, with $[X(0)]$ the south pole and $[\nInf]$ the point at infinity.

Under this identification, Theorem \ref{c:transport.totality} computes
\[
\iota[X(z)] = [X(z/|z|^2)] \ \ (z \neq 0), \qquad \iota[X(0)] = [\nInf], \qquad \iota[\nInf] = [X(0)],
\]
and $z \mapsto z/|z|^2 = 1/\bar{z}$ is inversion in the unit circle. On the sphere it is the reflection
\[
(X, Y, Z) \;\mapsto\; (X, Y, -Z)
\]
in the equatorial plane: the modulus inverts and the argument is unchanged, so the equator $|z| = 1$ is fixed pointwise and the two poles, $0$ and $\infty$, are exchanged. Both statements are Proposition \ref{c:transport.swap-rescale-normal-form} read in these coordinates, the second being the fixed locus $\{Q = 1\} \cup \text{(ii)}$ of Remark \ref{c:note.transport.where-each-stratum-comes} with the stratum (ii) that definite signature leaves empty. Inversion in a circle is anti-holomorphic; the holomorphic $z \mapsto 1/z$ of complex analysis is its composite with the reflection $y \mapsto -y$, which is a reflection of the base and so a versor of the model. Either way the content tested is the same, and it is the classical one: a reciprocal undefined at a single point of the plane becomes a total involution of a compact surface two dimensions up, with $1/0 = \infty$ and $1/\infty = 0$. Theorem \ref{c:transport.totality} reproduces the nineteenth-century resolution of $1/0$ by conformal compactification, as it must, and Theorem \ref{c:transport.minimality} adds that the two dimensions it uses are the fewest that will serve.

The same computation over the one-dimensional base $\mathbb{V} = \R$, $Q(x) = x^2$, has completion $\Cl(2,1)$, identifies $\mathcal{Q}$ with the real projective line $\R P^1$ in its classical homogeneous coordinates $[\,x : s\,]$, and returns for $\iota$ the classical inversion matrix $\begin{psmallmatrix} 0 & 1 \\ 1 & 0 \end{psmallmatrix}$, with $[\,0:1\,]$ and $[\,1:0\,]$ exchanged.

What no definite signature exhibits, in either dimension, is the null cone. The genuinely new content of Theorems \ref{c:transport.no-go-level-elements} and \ref{c:transport.totality}(3) lies on the ideal strata (i) and (ii), which are empty here and carry, in mixed signature, one class per null vector and one per null direction respectively. There the compactification adjoins a stratum rather than a point, and the element-level question is both posed and closed.

\subsection*{The hyperbolic atom as a matrix algebra: the Moore--Penrose pseudoinverse}
\phantomsection\label{c:note.transport.setting-hyperbolic-atom-matrix-algebra-moore}
The algebra $\Cl(1,1)$ is isomorphic to the algebra $M_2(\R)$ of real $2 \times 2$ matrices: with $e_0^2 = 1$, $e_1^2 = -1$, the assignment
\[
e_0 \mapsto E = \begin{psmallmatrix} 0 & 1 \\ 1 & 0 \end{psmallmatrix}, \qquad
e_1 \mapsto F = \begin{psmallmatrix} 0 & 1 \\ -1 & 0 \end{psmallmatrix}
\]
satisfies $E^2 = I$, $F^2 = -I$ and $EF + FE = 0$, hence extends to an algebra isomorphism (both algebras have dimension four; $M_2(\R)$ is simple). A base vector $x e_0 + y e_1$ maps to $\begin{psmallmatrix} 0 & x+y \\ x-y & 0 \end{psmallmatrix}$, of determinant $-(x^2 - y^2) = -Q$: \emph{a vector is null exactly when its matrix is singular}. In particular the hyperbolic pair $v = \tfrac12(e_0 + e_1)$, $w_0 = \tfrac12(e_0 - e_1)$ (so $v \cdot w_0 = \tfrac12$) maps to
\[
v \;\mapsto\; N = \begin{psmallmatrix} 0 & 1 \\ 0 & 0 \end{psmallmatrix}, \qquad
w_0 \;\mapsto\; N^{\mathsf T} = \begin{psmallmatrix} 0 & 0 \\ 1 & 0 \end{psmallmatrix}.
\]
Under this dictionary the objects of the section ``The refuted conjecture, and the no-go'' become textbook objects of linear algebra: the equation $v\,\sigma(v)\,v = v$ is Penrose's first equation $AXA = A$ for a \emph{generalized inverse} of a singular matrix, and the abundance of pseudo-inverses noted there is the familiar abundance of generalized inverses.

Matrix theory has a canonical selection among them: the Moore--Penrose pseudoinverse $A^{+}$, which for a rank-one matrix is $A^{+} = A^{\mathsf T}/\lVert A \rVert_F^{\,2}$. Applied to the null vectors of this base,
\[
(\lambda N)^{+} = \lambda^{-1} N^{\mathsf T}, \qquad (\lambda N^{\mathsf T})^{+} = \lambda^{-1} N,
\]
which is, under the isomorphism, \emph{exactly} the equivariant section $\sigma(\lambda v) = \lambda^{-1} w_0$, $\sigma(\lambda w_0) = \lambda^{-1} v$ of Remark \ref{c:note.transport.hyperbolic-atom}: the hyperbolic atom's self-completeness is the Moore--Penrose construction expressed in the Clifford algebra. The agreement is not a coincidence of formulas. By \ref{c:note.transport.hyperbolic-atom} the partner of a null vector in a hyperbolic plane is unique, so any rule at all that returns a null partner returns this one; the Euclidean transpose supplies it, and so does the metric description.

The agreement extends to the failure, which is the instructive part. It is classical that $A^{+}$ is canonical only relative to a choice of \emph{definite} inner product: it is built from the transpose, the adjoint of the Euclidean metric, and transforms correctly under orthogonal maps, $(UAV)^{+} = V^{\mathsf T} A^{+}\, U^{\mathsf T}$ for orthogonal $U, V$, but not under general invertible ones. Known theory therefore already signals that the obvious candidate for $\sigma$ must fail to respect an \emph{indefinite} symmetry group, which preserves no definite inner product. Theorem \ref{c:transport.no-go-level-elements} is the sharp form of that signal. For $n = 2$ the stabiliser of a null direction in $O(1,1)$ contains no null rotations, nothing shears, the partner set \ref{c:note.transport.hyperbolic-atom} is a single point, and the Moore--Penrose choice survives as an equivariant section (Remark \ref{c:note.transport.hyperbolic-atom}). For $n \geq 3$ the stabiliser acquires a unipotent radical of null rotations, the partner set becomes a torsor under it, and not only the Moore--Penrose rule but \emph{every conceivable rule}, valued in the Clifford algebra of any extension or in a matrix algebra over it, is sheared away. The no-go theorem thus marks the precise boundary of the Moore--Penrose phenomenon: canonical generalized inversion of null directions ends exactly where parabolic stabilisers begin.

\section{Type C: ideal fusion, and Rule 2}

\phantomsection\label{c:note.transport.setting-type-c-ideal-fusion-rule}
the section ``Type N: the isotropic obstruction, and Rule 1'' settled type N: no element-level repair in any ascent, a total repair at the level of points, and a minimal completion forced in two independent senses. This section settles type C, where the answer is different in form. The obstruction is not a direction without a partner but a splitting of the algebra into two ideals that no inner symmetry can exchange (Remark \ref{c:note.transport.typec-depth}); its repair costs one dimension rather than two, is delivered at the level of elements rather than points, and is unique up to isomorphism though not up to isometry.

\subsection*{Setting: one adjoined generator}
\phantomsection\label{c:note.transport.setting-setting-one-adjoined-generator}
Throughout this section $A = \Cl(p,q)$ with $p - q \equiv 1 \pmod 4$, so $n = p + q$ is odd, $\omega$ is central with $\omega^2 = +1$, and $A = \varepsilon_+ A \oplus \varepsilon_- A$ splits along the central projectors $\varepsilon_\pm = \tfrac12(1 \pm \omega)$ (Theorem \ref{c:transport.existence-classification-type-c}). We consider the Clifford algebras of one-dimensional extensions of the quadratic space: $\mathbb{V}' = \mathbb{V} \perp \langle \nu \rangle$ with $\nu$ orthogonal to $\mathbb{V}$ and $\nu^2 = \gamma \in \{+1, -1, 0\}$ (a nonzero square is rescaled to $\pm 1$; $\gamma = 0$ is the degenerate, null-generator case, which Theorem \ref{c:transport.one-generator-minimal-minimal-completions-classified} shows never fuses). Write $B_\gamma = \Cl(\mathbb{V}', Q')$; for $\gamma = +1$ this is $\Cl(p{+}1, q)$, for $\gamma = -1$ it is $\Cl(p, q{+}1)$. Since $Q$ is nondegenerate on $\mathbb{V}$, the blade basis gives a canonical unital inclusion $A \subseteq B_\gamma$ in every case, degenerate included \cite{Chevalley1997}.

\begin{lemma}[Decomposition and commutation]\label{c:transport.decomposition-commutation}
$B_\gamma = A \oplus A\nu$ as a vector space, and
\[
\nu\, a \;=\; \hat{a}\, \nu \qquad (a \in A),
\]
where $\hat{\cdot}$ is the grade involution of $A$. For $\gamma = \pm 1$, $\nu$ is invertible with $\nu^{-1} = \gamma\nu$.
\end{lemma}
\begin{proof}
The blades of $B_\gamma$ are the blades of $A$ together with the blades of $A$ times $\nu$ (ordering $\nu$ last), giving the direct sum; the count $2^{n+1} = 2^n + 2^n$ confirms it, and holds for $\gamma = 0$ as well \cite{Chevalley1997} (\ref{c:std-algebra.blade-basis}). For the commutation rule: $\nu$ anticommutes with every generator $e_i$ of $A$ (orthogonality, \ref{c:std-algebra.orthogonal-anticommute}), so for a blade $e_I$ of grade $k$, $\nu\, e_I = (-1)^k e_I\, \nu = \hat{e_I}\,\nu$, and the rule extends linearly. This is the commutation rule $v\,x = \hat{x}\,v, \qquad x\,v = v\,\hat{x}$ (\ref{c:primer.commutation-rule}) applied to the adjoined generator. Invertibility: $\nu \cdot \gamma\nu = \gamma \nu^2 = \gamma^2 = 1$ for $\gamma = \pm 1$.

\end{proof}

\begin{proposition}[The ascent renders the grade involution inner]\label{c:transport.ascent-renders-grade-involution-inner}
Let $\gamma = \pm 1$. Conjugation by $\nu$ in $B_\gamma$ restricts on $A$ to the grade involution:
\[
\nu\, a\, \nu^{-1} \;=\; \hat{a} \qquad (a \in A).
\]
In particular $\nu \omega \nu^{-1} = -\omega$ and $\nu \varepsilon_\pm \nu^{-1} = \varepsilon_\mp$: the two formerly central projectors become an inner-conjugate pair. Consequently the two-sided ideal of $B_\gamma$ generated by either projector is all of $B_\gamma$, by the one-line computation
\[
\varepsilon_+ \;+\; \nu\,\varepsilon_+\,\nu^{-1} \;=\; \varepsilon_+ + \varepsilon_- \;=\; 1 .
\]
The obstruction of Remark \ref{c:note.transport.typec-depth} is thereby resolved in the exact sense there identified: the symmetry exchanging the two ideals, outer on the split algebra, is inner in the ascent, and its availability is what fuses them.
\end{proposition}
\begin{proof}
By Lemma \ref{c:transport.decomposition-commutation}, $\nu a \nu^{-1} = \hat{a}\,\nu\nu^{-1} = \hat{a}$. Since $\omega$ has odd grade $n$, $\hat{\omega} = -\omega$ (\ref{c:std-algebra.involution-properties}), whence $\varepsilon_\pm \mapsto \varepsilon_\mp$. The displayed sum lies in the ideal generated by $\varepsilon_+$ (both summands are multiples of $\varepsilon_+$ by ring elements), and equals $1$.

\end{proof}

\subsection*{The fusion theorem}
\phantomsection\label{c:note.transport.setting-fusion-theorem}
Two lemmas prepare the sharp form of the theorem: the corner algebra cut out by $\varepsilon_+$ is the even subalgebra, and a $2 \times 2$ system of matrix units assembles $B_\gamma$ from that corner.

\begin{lemma}[The corner is the even subalgebra]\label{c:transport.corner-even-subalgebra}
The map $\phi : \Cl^0(p,q) \to \varepsilon_+ A$, $\phi(a) = \varepsilon_+ a$, is an isomorphism of unital algebras (with unit $\varepsilon_+$ on the right-hand side). Moreover $\varepsilon_+ B_\gamma\, \varepsilon_+ = \varepsilon_+ A$.
\end{lemma}
\begin{proof}
$\phi$ is an algebra map because $\varepsilon_+$ is a central idempotent of $A$: $\phi(a)\phi(b) = \varepsilon_+ a\, \varepsilon_+ b = \varepsilon_+^2 ab = \phi(ab)$, and $\phi(1) = \varepsilon_+$. Injectivity: if $a$ is even and $\varepsilon_+ a = 0$, then $a = (\varepsilon_+ + \varepsilon_-)a = \varepsilon_- a$, i.e.\ $2a = a - \omega a$, i.e.\ $\omega a = -a$. But $\omega$ has odd grade ($n$ odd) while $a$ is even, so $\omega a$ is odd and $-a$ is even; a common value is possible only if both vanish, so $a = 0$. Surjectivity follows by dimension: $\dim \Cl^0 = 2^{n-1} = \tfrac12 \dim A = \dim \varepsilon_+ A$ (the first equality is \ref{c:std-algebra.z2-grading}), the last equality because $x \mapsto (\varepsilon_+ x, \varepsilon_- x)$ splits $A$ into two summands exchanged (linearly) by the grade involution, which maps $\varepsilon_+ x \mapsto \varepsilon_- \hat{x}$ and is bijective. For the second statement: $\varepsilon_+ B_\gamma \varepsilon_+ = \varepsilon_+ A \varepsilon_+ + \varepsilon_+ A \nu \varepsilon_+$ by Lemma \ref{c:transport.decomposition-commutation}; the first term is $\varepsilon_+ A$ (centrality in $A$), and the second vanishes: $\nu \varepsilon_+ = \varepsilon_- \nu$ (Proposition \ref{c:transport.ascent-renders-grade-involution-inner}), so $\varepsilon_+ a \nu \varepsilon_+ = \varepsilon_+ a \varepsilon_- \nu = (\varepsilon_+ \varepsilon_-) a \nu = 0$.

\end{proof}

\begin{lemma}[Matrix units, and the Peirce isomorphism]\label{c:transport.matrix-units-peirce-isomorphism}
Let $\gamma = \pm 1$ and set
\[
E_{11} = \varepsilon_+, \qquad E_{22} = \varepsilon_-, \qquad E_{12} = \varepsilon_+ \nu, \qquad E_{21} = \gamma\, \varepsilon_- \nu .
\]
Then $E_{11} + E_{22} = 1$ and $E_{ij} E_{kl} = \delta_{jk} E_{il}$: a full $2 \times 2$ system of matrix units. Consequently
\[
B_\gamma \;\cong\; \Mat_2\bigl(E_{11} B_\gamma E_{11}\bigr), \qquad b \;\longmapsto\; \bigl(E_{1i}\, b\, E_{j1}\bigr)_{ij},
\]
an isomorphism of unital algebras with inverse $(c_{ij}) \mapsto \sum_{i,j} E_{i1} c_{ij} E_{1j}$.
\end{lemma}
\begin{proof}
The relations, using $\nu \varepsilon_\pm = \varepsilon_\mp \nu$ (Proposition \ref{c:transport.ascent-renders-grade-involution-inner}), $\varepsilon_\pm^2 = \varepsilon_\pm$, $\varepsilon_+\varepsilon_- = 0$, $\nu^2 = \gamma$ and $\gamma^2 = 1$:
\begin{align*}
E_{12}E_{21} &= \varepsilon_+ \nu\, \gamma \varepsilon_- \nu = \gamma\, \varepsilon_+ (\nu \varepsilon_-) \nu = \gamma\, \varepsilon_+ \varepsilon_+ \nu^2 = \gamma^2 \varepsilon_+ = E_{11},\\
E_{21}E_{12} &= \gamma\, \varepsilon_- \nu\, \varepsilon_+ \nu = \gamma\, \varepsilon_-^2 \nu^2 = E_{22},
\end{align*}
and the remaining relations follow the same way:
\[
E_{12}^2 = \varepsilon_+ \nu \varepsilon_+ \nu = \varepsilon_+ \varepsilon_- \nu^2 = 0, \qquad E_{21}^2 = 0,
\]
\[
E_{11}E_{12} = E_{12}, \qquad
E_{12}E_{11} = \varepsilon_+ \nu \varepsilon_+ = \varepsilon_+\varepsilon_-\nu = 0 .
\] For the Peirce isomorphism, write $\phi(b)_{ij} = E_{1i} b E_{j1} \in E_{11} B_\gamma E_{11}$. Multiplicativity, using $E_{k1}E_{1k} = E_{kk}$ and $E_{11} + E_{22} = 1$:
\[
\sum_k \phi(b)_{ik}\,\phi(b')_{kj}
\;=\; \sum_k E_{1i}\, b\, E_{k1}E_{1k}\, b'\, E_{j1}
\;=\; E_{1i}\, b \Bigl(\textstyle\sum_k E_{kk}\Bigr) b'\, E_{j1}
\;=\; \phi(bb')_{ij} .
\]
Unitality: $\phi(1)_{ij} = E_{1i}E_{j1} = \delta_{ij} E_{11}$. The displayed inverse is checked directly:
\[
\sum_{i,j} E_{i1}\, (E_{1i}\, b\, E_{j1})\, E_{1j} \;=\; \Bigl(\textstyle\sum_i E_{ii}\Bigr)\, b\, \Bigl(\textstyle\sum_j E_{jj}\Bigr) \;=\; b ,
\]
and conversely on matrices. This is the standard matrix-units criterion \cite{LamFC2001}, reproved for self-containment.

\end{proof}

\begin{theorem}[Rule 2: ideal fusion, sharp form]\label{c:transport.rule-2-ideal-fusion-sharp-form}
Let $p - q \equiv 1 \pmod 4$ and $\gamma = \pm 1$. Then
\[
B_\gamma \;=\; \Cl(p{+}1,q) \ \text{or}\ \Cl(p,q{+}1) \;\cong\; \Mat_2\bigl(\Cl^0(p,q)\bigr),
\]
by the matrix units of Lemma \ref{c:transport.matrix-units-peirce-isomorphism} and the corner identification of Lemma \ref{c:transport.corner-even-subalgebra}. Consequently:
\begin{enumerate}
\item $B_\gamma$ is simple, and its centre is $\R$: the fusion does not reintroduce a central splitting, and $B_\gamma$ carries no type-C obstruction.
\item The images of $\varepsilon_\pm$ in $B_\gamma$ are a conjugate pair of non-central idempotents, the diagonal matrix units of the fused algebra, exchanged by the inner automorphism $\mathrm{Ad}_\nu$, which restricts on $A$ to the grade involution (Proposition \ref{c:transport.ascent-renders-grade-involution-inner}).
\item The two ascents are isomorphic as algebras: $\Cl(p{+}1,q) \cong \Cl(p,q{+}1)$, both canonically $\Mat_2(\Cl^0(p,q))$, with an isomorphism through the matrix units that fixes the corner.
\end{enumerate}
\end{theorem}
\begin{proof}
The isomorphism $B_\gamma \cong \Mat_2(\varepsilon_+ B_\gamma \varepsilon_+) = \Mat_2(\varepsilon_+ A) \cong \Mat_2(\Cl^0(p,q))$ chains Lemmas \ref{c:transport.matrix-units-peirce-isomorphism} and \ref{c:transport.corner-even-subalgebra}; note the matrix units and the corner are the same for both signs of $\gamma$ up to the harmless factor $\gamma$ in $E_{21}$, which is why the right-hand side is $\gamma$-independent; this proves (3). For (1): $\varepsilon_+ A$ is semisimple (a direct summand of the semisimple $A$, Theorem \ref{c:transport.classification-real-clifford-algebras}) with centre $\varepsilon_+ Z(A) = \R\,\varepsilon_+$ one-dimensional (here $Z(A) = \R \oplus \R\omega = \R\varepsilon_+ \oplus \R\varepsilon_-$ by Lemma \ref{c:transport.centre}, and the centre of a direct sum is the direct sum of the centres), and a semisimple algebra whose centre is one-dimensional has exactly one Wedderburn block, i.e.\ is simple. Matrix algebras over simple algebras are simple (\ref{c:std-algebra.morita-amplification}), so $B_\gamma$ is simple; its centre is $\R$ directly by Lemma \ref{c:transport.centre}, since $n + 1$ is even. (2) is Proposition \ref{c:transport.ascent-renders-grade-involution-inner} together with the identification of $\varepsilon_\pm$ as $E_{11}, E_{22}$; a diagonal matrix unit in $\Mat_2$ of a unital algebra is never central.

\end{proof}

\begin{remark}[The $\Mat_2$ pattern: the two rules compared]\label{c:note.transport.2-pattern-two-rules-compared}
Both minimal ascents turn out to be $2 \times 2$ matrix amplifications, of the whole algebra for type N and of its even half for type C (the section ``The mod-8 clock and the one-way door'' proves the first and reads both as the transport's own bookkeeping). The comparison exposes the cost asymmetry precisely: Rule 1 adds two dimensions and preserves the Morita class, $d$ fixed; Rule 2 adds one dimension, halves the coefficient algebra to its even part, and moves $d$ by $\pm 1$.
\end{remark}

\begin{theorem}[One generator is minimal; the minimal completions classified]\label{c:transport.one-generator-minimal-minimal-completions-classified}
\changedmark{gray}{reworded}
Let $A = \Cl(p,q)$, $p - q \equiv 1 \pmod 4$, and consider the Clifford algebras $\Cl(\mathbb{V}', Q')$ of arbitrary extensions $(\mathbb{V}', Q') \supseteq (\mathbb{V}, Q)$ of the quadratic space, $Q'$ possibly degenerate, each containing $A$ canonically and unitally. Then:
\begin{enumerate}
\item $\dim \mathbb{V}' = n$ (no ascent): $A$ itself is not simple (Theorem \ref{c:transport.existence-classification-type-c}).
\item $\dim \mathbb{V}' = n + 1$, adjoined generator null ($\gamma = 0$): the fusion fails structurally. The subspace $J = A\nu$ is a nonzero proper two-sided ideal of $B_0$ with $J^2 = 0$, and $B_0 / J \cong A$: the degenerate ascent is an infinitesimal thickening of $A$, not a completion of it. (This square-zero ideal is the fibre-level germ of the degeneration sheaf of the series' sheaf-theoretic paper \cite{ButlerT4}.)
\item $\dim \mathbb{V}' = n + 1$, adjoined generator nondegenerate ($\gamma = \pm 1$ after rescaling): both ascents are simple and $\cong \Mat_2(\Cl^0(p,q))$ (Theorem \ref{c:transport.rule-2-ideal-fusion-sharp-form}).
\item $\dim \mathbb{V}' \geq n + 2$: $\dim \Cl(\mathbb{V}') \geq 2^{n+2}$.
\end{enumerate}
Hence the minimal simple Clifford completions of $A$ are exactly the two one-generator ascents $\Cl(p{+}1,q)$ and $\Cl(p,q{+}1)$, of dimension $2^{n+1}$; they are inequivalent as quadratic data (the extensions $\langle +1 \rangle$ and $\langle -1 \rangle$ are non-isometric) but isomorphic as algebras; and so the minimal completion of a type-C obstruction is \emph{unique up to isomorphism}, with the residual geometric choice, the sign of the adjoined direction's square, not detected by the completed algebra.
\end{theorem}
\begin{proof}
(1) is Theorem \ref{c:transport.existence-classification-type-c}. (2): by Lemma \ref{c:transport.decomposition-commutation} (which holds for $\gamma = 0$), $B_0 = A \oplus A\nu$ with $\nu a = \hat{a}\nu$ and $\nu^2 = 0$. Then $J = A\nu$ satisfies: $J \neq 0$ ($\nu \neq 0$ and the sum is direct); $1 \notin J$ (again directness), so $J$ is proper; two-sidedness, since $A J = A A \nu \subseteq J$, $J A = A \nu A = A \hat{A} \nu \subseteq J$, and $\nu J = \nu A \nu = \hat{A} \nu^2 = 0 \subseteq J$, $J\nu \subseteq A\nu^2 = 0$; and $J^2 = A\nu A \nu = A \hat{A} \nu^2 = 0$. A simple algebra has no such ideal. The quotient: $B_0 / J = (A \oplus A\nu)/A\nu \cong A$ as algebras, the projection sending $\nu$ to $0$. (3) is Theorem \ref{c:transport.rule-2-ideal-fusion-sharp-form}; a nonzero square rescales to $\pm 1$ by replacing $\nu$ with $\nu/\sqrt{|\nu^2|}$. (4) is the dimension formula $\dim \Cl = 2^{\dim \mathbb{V}'}$, degenerate or not (\ref{c:std-algebra.blade-basis}). The classification statement assembles (1)--(4) with Theorem \ref{c:transport.rule-2-ideal-fusion-sharp-form}(3); non-isometry of $\langle +1\rangle$ and $\langle -1 \rangle$ is Sylvester's law of inertia (\ref{c:std-algebra.sylvester-inertia}).

\end{proof}

\begin{remark}[A bit the algebra forgets]\label{c:note.transport.bit-algebra-forgets}
Theorem \ref{c:transport.one-generator-minimal-minimal-completions-classified} sharpens the disanalogy with Rule 1 into its correct form. In Rule 1 the signature increment $(+1,+1)$ is unique outright (Theorem \ref{c:transport.minimality}). In Rule 2 the \emph{algebraic} completion is likewise unique (there is one fused algebra, $\Mat_2(\Cl^0(p,q))$), but its geometric realisations are two, and no fusion-internal criterion distinguishes them: the sign $\gamma$ is a genuine geometric datum with no algebraic counterpart. The clock positions differ ($d + 1$ and $d - 1$ mod 8, i.e.\ residues $2$ and $0$, or $6$ and $4$), yet the classification table of Theorem \ref{c:transport.classification-real-clifford-algebras} is symmetric about the split residues (positions $0$ and $2$ are both real type, $4$ and $6$ both quaternionic), which is the table-level reflection of the isomorphism proved intrinsically in Theorem \ref{c:transport.rule-2-ideal-fusion-sharp-form}(3). A canonical choice of $\gamma$, if one is ever wanted, must therefore be imported from structure outside the algebra (an orientation of the ambient geometry), and we record the fact here at the level of algebra only: \emph{the fusion is forced, the sign is free}.
\end{remark}

\begin{remark}[A second refuted conjecture, recorded]\label{c:note.transport.second-refuted-conjecture-recorded}
The natural conjecture is that the two minimal completions are ``in general non-isomorphic,'' reading the distinct clock positions $d \pm 1$ as distinct algebras. The computation refutes this: the matrix-unit construction of Lemma \ref{c:transport.matrix-units-peirce-isomorphism} is independent of $\gamma$, and the two ascents are isomorphic always. The conjecture and its refutation are both recorded, in the practice of this paper's the section ``The refuted conjecture, and the no-go'', where the founding conjecture of the programme is refuted in the same way: the surviving true statement (geometric twoness, algebraic oneness) is sharper than either the conjecture or its naive negation.
\end{remark}

\begin{example}[The three smallest fusions]\label{c:note.transport.three-smallest-fusions}
$\Cl(1,0) \cong \R \oplus \R$: both ascents give $\Mat_2(\Cl^0(1,0)) = \Mat_2(\R)$, and indeed $\Cl(2,0) \cong \Cl(1,1) \cong \Mat_2(\R)$, the familiar coincidence that the Euclidean and Minkowski plane algebras agree, here encountered as forced by fusion uniqueness rather than as an accident of the classification table. $\Cl(0,3) \cong \HH \oplus \HH$: both ascents give $\Mat_2(\Cl^0(0,3)) = \Mat_2(\HH)$, i.e.\ $\Cl(1,3) \cong \Cl(0,4) \cong \Mat_2(\HH)$: the spacetime algebra is one of the two geometric realisations of the unique fusion of the double quaternions. $\Cl(2,1) \cong \Mat_2(\R) \oplus \Mat_2(\R)$ (the algebra of $2{+}1$-dimensional spacetime): both ascents give $\Mat_2(\Cl^0(2,1)) \cong \Mat_4(\R)$, i.e.\ $\Cl(3,1) \cong \Cl(2,2) \cong \Mat_4(\R)$.
\end{example}

\subsection*{Minimality against arbitrary embeddings, for the fusion}
\phantomsection\label{c:note.transport.setting-minimality-against-arbitrary-embeddings-fusion}
Theorem \ref{c:transport.one-generator-minimal-minimal-completions-classified} quantifies over Clifford algebras of extensions of the quadratic space, the same normalisation used for Rule 1 in the section ``Minimality among null-lift completions''. For Rule 1 the stronger, representation-independent statement was Theorem \ref{c:transport.representation-independent-minimality-completion}, settled on the ambient quadric. For Rule 2 the corresponding question is settled by module theory instead, and the answer is stronger still: the fused algebra is minimal among \emph{all} simple real algebras, Clifford or not, into which the split algebra embeds unitally.

\begin{theorem}[Embedding-independent minimality of the fusion]\label{c:transport.embedding-independent-minimality-fusion}
\changedmark{gray}{reworded}
Let $p - q \equiv 1 \pmod 4$, $A = \Cl(p,q) \cong \Mat_k(\mathbb{F}) \oplus \Mat_k(\mathbb{F})$ with $\mathbb{F} \in \{\R, \HH\}$, $f = \dim_\R \mathbb{F}$, so $\dim_\R A = 2k^2 f = 2^n$. Let $B$ be any \emph{simple} finite-dimensional real algebra admitting a unital algebra embedding $A \hookrightarrow B$. Then
\[
\dim_\R B \;\geq\; 2^{\,n+1},
\]
and equality holds if and only if $B \cong \Mat_{2k}(\mathbb{F}) \cong \Mat_2(\Cl^0(p,q))$: the fused algebra of Theorem \ref{c:transport.rule-2-ideal-fusion-sharp-form} is the minimal simple environment of $A$ in the absolute sense, and it is the unique one of minimal dimension up to isomorphism.
\end{theorem}
\begin{proof}
Write $B \cong \Mat_m(D)$, $D \in \{\R, \C, \HH\}$, $d = \dim_\R D$, with unique irreducible module $W = D^m$, $\dim_\R W = md$. $B$ acts faithfully on $W$, hence so does $A$; a faithful module of $A$, split as in the statement by \ref{c:transport.classification-real-clifford-algebras}, contains both irreducibles $V_\pm$ (of $\dim_\R = kf$ each), since otherwise one of the central projectors $\varepsilon_\pm \neq 0$ would act as zero. So $W \cong V_+^{\,a} \oplus V_-^{\,b}$ as an $A$-module with $a, b \geq 1$ and
\[
md \;=\; (a + b)\,k f \;\geq\; 2kf .
\]
Moreover the right $D$-action on $W$ commutes with $A$, so $D$ embeds unitally in the commutant (strictly, the right action embeds $D^{\mathrm{op}}$; harmless, as $\R$, $\C$ and $\HH$ are each isomorphic to their opposites, and the dimension bounds below are opposite-invariant) $\mathrm{End}_A(W) \cong \Mat_a(\mathbb{F}) \oplus \Mat_b(\mathbb{F})$ (using $\mathrm{End}_{\Mat_k(\mathbb{F})}(\mathbb{F}^k) \cong \mathbb{F}$ as real algebras), hence, projecting to each summand, unitally in both $\Mat_a(\mathbb{F})$ and $\Mat_b(\mathbb{F})$.

Case $\mathbb{F} = \R$ ($f = 1$, $2k^2 = 2^n$). If $D = \R$: $m \geq 2k$ and $\dim B = m^2 \geq 4k^2 = 2^{n+1}$, with equality iff $m = 2k$, $B \cong \Mat_{2k}(\R)$. If $D = \C$: $\C \hookrightarrow \Mat_a(\R)$ forces $a \geq 2$, likewise $b \geq 2$, so $2m = (a+b)k \geq 4k$, $m \geq 2k$, and $\dim B = 2m^2 \geq 8k^2 = 2^{n+2}$. If $D = \HH$: $a, b \geq 4$, so $4m \geq 8k$, $\dim B = 4m^2 \geq 16k^2$. Case $\mathbb{F} = \HH$ ($f = 4$, $8k^2 = 2^n$). If $D = \HH$: $4m \geq 8k$ gives $m \geq 2k$ and $\dim B = 4m^2 \geq 16k^2 = 2^{n+1}$, equality iff $B \cong \Mat_{2k}(\HH)$. If $D = \C$: $\C \hookrightarrow \Mat_a(\HH)$ allows $a \geq 1$, but then $2m = 4(a+b)k \geq 8k$ gives $m \geq 4k$ and $\dim B = 2m^2 \geq 32k^2$. If $D = \R$: $m = 4(a+b)k \geq 8k$ and $\dim B = m^2 \geq 64k^2$. In every case the bound holds, and the equality case is exactly $\Mat_{2k}(\mathbb{F})$. Finally $\Mat_{2k}(\mathbb{F}) \cong \Mat_2(\Mat_k(\mathbb{F})) \cong \Mat_2(\Cl^0(p,q))$ by Lemma \ref{c:transport.corner-even-subalgebra} and Theorem \ref{c:transport.rule-2-ideal-fusion-sharp-form}.

\end{proof}

\begin{remark}[The cost of transport, independently of representation]\label{c:note.transport.cost-transport-independently-representation}
Theorems \ref{c:transport.representation-independent-minimality-completion} and \ref{c:transport.embedding-independent-minimality-fusion} together make the picture uniform across the two rules, and in both cases stronger than the quadratic-space normalisations under which the rules were first proved minimal. For type N, what an ambient algebra must supply is a quadric large enough to host the \emph{points}, and the conformal Gram matrix quantifies it exactly. For type C, what it must supply is a \emph{module} large enough to carry both ideals faithfully with the commutant the fusion requires, and the module count quantifies that exactly, against arbitrary simple real algebras rather than Clifford targets alone. In both cases the cost is independent of the embedding, and in both cases the minimal environment is the one the transport rule constructs.
\end{remark}

\subsection*{One-dimensional extensions: where the obstruction lies, and the commuting square}
\phantomsection\label{c:note.transport.setting-one-dimensional-extensions-where-obstruction}
The physical algebras the series works with are even-dimensional and simple. Adjoining one direction makes them odd-dimensional, where the volume element is central and the type-C question arises. This subsection places the obstruction exactly among the one-dimensional extensions of an even-dimensional algebra, identifies its ideals, and shows that its Rule 2 repair by a generator of the opposite sign is the Rule 1 ascent of the section ``Type N: the isotropic obstruction, and Rule 1'': the two rules, compared in Remark \ref{c:note.transport.2-pattern-two-rules-compared} as amplifications of the whole algebra and of an even half, coincide on every such extension, where the even half is the base. Throughout, write $c_u(x) = u x u^{-1}$ for conjugation by an invertible $u$.

\begin{lemma}[The even subalgebra of a one-dimensional extension]\label{c:transport.even-subalgebra-one-dimensional-extension}
Let $\nu$ be a generator adjoined orthogonally to $\Cl(p,q)$. If $\nu^2 = -1$ then $e_i \mapsto e_i\nu$ extends to an isomorphism $\Cl(p,q) \cong \Cl^0(p,q{+}1)$; if $\nu^2 = +1$ then the same map gives $\Cl(q,p) \cong \Cl^0(p{+}1,q)$.
\end{lemma}
\begin{proof}
The images are even. Their squares are $(e_i\nu)^2 = e_i\nu e_i\nu = -e_i^2\nu^2$, which is $e_i^2$ for $\nu^2 = -1$ and $-e_i^2$ for $\nu^2 = +1$; for $i \neq j$, $(e_i\nu)(e_j\nu) = -e_ie_j\nu^2$ and $(e_j\nu)(e_i\nu) = -e_je_i\nu^2 = e_ie_j\nu^2$, so distinct images anticommute. By the universal property (\ref{c:std-algebra.existence-uniqueness}) the assignment extends to a unital homomorphism from $\Cl(p,q)$ (respectively $\Cl(q,p)$) into the even subalgebra. It is surjective: for $i \neq j$, $(e_i\nu)(e_j\nu) = -\nu^2 e_ie_j = \mp e_ie_j$, so every even blade of the extension not involving $\nu$ is, up to sign, a product of images, and every even blade involving $\nu$ is $e_I\nu$ with $|I|$ odd, which is $\pm(e_{i_1}\nu)(e_{i_2}\nu)(e_{i_3}\nu)\cdots$ after pairing all but the last factor. Both sides have dimension $2^n$, so the surjection is an isomorphism. The script confirms the relations, the evenness, and the dimension on $(1,3)$ and $(2,0)$ with both signs of $\nu^2$.

\end{proof}

\begin{theorem}[Where the obstruction lies, and the commuting square]\label{c:transport.where-obstruction-lies-commuting-square}
\changedmark{BurntOrange}{changed since the deposit}
Let $A = \Cl(p,q)$ with $n = p + q$ even, and let $A_s = \Cl(p,q{+}1)$ and $A_t = \Cl(p{+}1,q)$ be its spacelike and timelike one-dimensional extensions.
\begin{enumerate}
\item[(a)] Exactly one of $A_s$, $A_t$ carries the type-C obstruction: $A_s$ if $p - q \equiv 2 \pmod 4$, $A_t$ if $p - q \equiv 0 \pmod 4$. The other has complex type: its volume element is central with square $-1$, its centre is $\C$, and it is simple.
\item[(b)] Each of the two ideals of the obstructed extension is a copy of $A$: $\varepsilon_\pm A_s \cong \Cl^0(p,q{+}1) \cong A$ in the first case, $\varepsilon_\pm A_t \cong \Cl^0(p{+}1,q) \cong \Cl(q,p) \cong A$ in the second (the last isomorphism because the residues $0$ and $4$ are fixed by negation mod $8$).
\item[(c)] The Rule 2 repair of the obstructed extension by a generator of the opposite sign is $\Cl(p{+}1,q{+}1)$, the Rule 1 ascent of $A$: the two adjoined generators $e_{n+1}$ (the extension) and $\nu$ (the repair) span a hyperbolic plane, the null vectors $f = \tfrac12(u \pm w)$, $f^* = \tfrac12(u \mp w)$, with $u$ the adjoined generator of square $+1$ and $w$ that of square $-1$ (both combinations are null; the opposite signs make $ff^* + f^*f = 1$) are the ladder pair of Theorem \ref{c:transport.rule-1-morita-trivial}, and the two amplifications agree: $\Mat_2(\Cl^0(\text{extension})) = \Mat_2(A)$. The residue returns to $d$.
\item[(d)] The Rule 2 repair by a generator of the same sign is $\Cl(p,q{+}2)$ or $\Cl(p{+}2,q)$, isomorphic to $\Cl(p{+}1,q{+}1)$ as an algebra, at residue $d \mp 2$; the adjoined plane is definite and carries no null vector, so this adjunction is not the Rule 1 ascent of (c), although the algebra it produces is isomorphic to one.
\end{enumerate}
\end{theorem}
\begin{proof}
(a) By Theorem \ref{c:transport.existence-classification-type-c} with Lemmas \ref{c:transport.centre} and \ref{c:transport.sign}, an odd-dimensional algebra is type C exactly when $p' - q' \equiv 1 \pmod 4$ and complex type exactly when $p' - q' \equiv 3 \pmod 4$; these exhaust the odd residues. For $A_s$, $p' - q' = p - q - 1$, which is $\equiv 1$ iff $p - q \equiv 2$ and $\equiv 3$ iff $p - q \equiv 0 \pmod 4$. For $A_t$, $p' - q' = p - q + 1$, which is $\equiv 1$ iff $p - q \equiv 0$ and $\equiv 3$ iff $p - q \equiv 2$. Since $n$ is even, $p - q$ is even, so exactly one case obtains for each extension and the two extensions fall on opposite sides. Simplicity of the complex-type algebra: its centre $\C$ is a field, and a semisimple algebra with a field for centre has one Wedderburn block.

(b) The corner isomorphism $\varepsilon_+ A' \cong \Cl^0(A')$ is Lemma \ref{c:transport.corner-even-subalgebra}, and $\varepsilon_- A' \cong \varepsilon_+ A'$ through the grade involution, which exchanges the ideals (Proposition \ref{c:transport.ascent-renders-grade-involution-inner}). Lemma \ref{c:transport.even-subalgebra-one-dimensional-extension} identifies $\Cl^0(p,q{+}1) \cong \Cl(p,q)$ and $\Cl^0(p{+}1,q) \cong \Cl(q,p)$. In the timelike case, $p - q \equiv 0 \pmod 4$ gives $q - p \equiv -(p-q) \pmod 8$, and the classification (Theorem \ref{c:transport.classification-real-clifford-algebras}) assigns $\Cl(q,p)$ and $\Cl(p,q)$ the same type and dimension when $p - q \equiv 0 \pmod 4$ (residues $0$ and $4$ are each fixed by negation mod 8), so $\Cl(q,p) \cong \Cl(p,q)$ there (the script exercises the timelike form of Lemma \ref{c:transport.even-subalgebra-one-dimensional-extension} on $(1,3)$ and $(2,0)$).

(c) The spacelike case: the repair generator $\nu$ has $\nu^2 = +1$ and the extension generator $e_{n+1}$ has square $-1$, so the Gram matrix of the pair is $\mathrm{diag}(+1,-1)$, of signature $(1,1)$: a hyperbolic plane. Then $f = \tfrac12(\nu + e_{n+1})$ satisfies $f^2 = \tfrac14(\nu^2 + e_{n+1}^2 + \nu e_{n+1} + e_{n+1}\nu) = 0$, likewise $f^* = \tfrac12(\nu - e_{n+1})$, and $ff^* + f^*f = \tfrac14\bigl((\nu + e_{n+1})(\nu - e_{n+1}) + (\nu - e_{n+1})(\nu + e_{n+1})\bigr) = \tfrac12(\nu^2 - e_{n+1}^2) = 1$. The quadratic space is $\R^{p,q} \perp \langle e_{n+1}\rangle \perp \langle \nu \rangle = \R^{p+1,q+1}$ whichever order the two directions are adjoined in, so the algebra is $\Cl(p{+}1,q{+}1)$ on the nose, and Theorem \ref{c:transport.rule-1-morita-trivial} reads it as $\Mat_2(A)$ through the ladder pair, while Lemma \ref{c:transport.matrix-units-peirce-isomorphism} and Theorem \ref{c:transport.rule-2-ideal-fusion-sharp-form} read it as $\Mat_2(\Cl^0(A_s))$ through $\varepsilon_\pm$ and $\nu$; by (b) the two coefficient algebras are the same algebra. The residue of $\Cl(p{+}1,q{+}1)$ is $p - q$. The timelike case is the same computation with the signs of $\nu^2$ and $e_{n+1}^2$ exchanged and $f = \tfrac12(e_{n+1} + \nu)$, $f^* = \tfrac12(e_{n+1} - \nu)$.

(d) Theorem \ref{c:transport.rule-2-ideal-fusion-sharp-form}(3) gives the isomorphism of the two repairs and the residues $d' \pm 1$ of the fused algebras with $d' = d \mp 1$ the residue of the extension; the same-sign repair has residue $d \mp 2$. The adjoined plane has Gram matrix $\pm\,\mathrm{diag}(1,1)$, definite, so $(a e_{n+1} + b\nu)^2 = \pm(a^2 + b^2) \neq 0$ for $(a,b) \neq 0$: the adjoined plane carries no null vector, hence no ladder pair, so the adjunction is not a $(1,1)$ ascent; the algebra it produces is nevertheless $\Mat_2(A)$, and does contain ladder pairs, just none inside the adjoined plane.

The script verifies (a) on every even-dimensional $(p,q)$ with $2 \leq n \leq 6$ by computing the volume element's centrality and square directly, and verifies (c) and (d) on the same grid: the obstructed extension is type C, the opposite-sign repair is at residue $d$ with a hyperbolic adjoined plane carrying a ladder pair, and the same-sign repair is two residues away with a definite plane.

\end{proof}

\begin{corollary}[Which residues obstruct which way]\label{c:transport.which-residues-obstruct-which-way}
\changedmark{BurntOrange}{changed since the deposit}
The even residues whose spacelike extension is type C are $\{2, 6\}$; those whose timelike extension is type C are $\{0, 4\}$. In particular every even residue has exactly one obstructing one-dimensional extension. The physical lineage is at residue $6$ and its Rule 1 ascents never leave it (Theorem \ref{c:transport.rule-1-morita-trivial}); by Theorem \ref{c:transport.where-obstruction-lies-commuting-square} it leaves residue $6$ only through a spacelike direction, into the obstruction at residue $5$, and the opposite-sign repair returns it to $6$ as $\Mat_2$ of the algebra it started from. The same holds at residue $2$, with spacelike read as square $-1$: in T14's causal convention, where the mirror's time is its minus sign, that extension is the mirror's second time, and the mirror's own Kaluza-Klein direction gives the unobstructed $\Cl(4,1)$ (Proposition \emph{The mirror's extensions under the causal convention}).
\end{corollary}
\begin{proof}
Residues $2$ and $6$ are the even residues $\equiv 2 \pmod 4$; residues $0$ and $4$ are those $\equiv 0$. The rest is Theorem \ref{c:transport.where-obstruction-lies-commuting-square}(a) and (c) read at $d = 6$ and $d = 2$. That Rule 1 ascents fix residue $6$ is \ref{c:transport.rule-1-morita-trivial}. The script lists both sets from the grid.

\end{proof}

\phantomsection\label{c:note.transport.setting-one-dimensional-extensions-where-obstruction-2}
The residues $\{2, 6\}$ are also the oriented and causal residues of the dark-matter paper \cite[Cor.~3.11]{ButlerT14}. That the towers carrying matter are exactly the towers whose spacelike extensions obstruct is, at this level, arithmetic ($2 \bmod 4$ twice over); the coincidence is recorded and not built on.

\begin{corollary}[Odd Lorentzian dimension: exactly one real form is split]\label{c:transport.odd-lorentzian-dimension-exactly-one-real}
Let $n$ be odd and consider the two real forms of the Lorentzian algebra in dimension $n$, $\Cl(1,n{-}1)$ (time of square $+1$) and $\Cl(n{-}1,1)$ (time of square $-1$). Then $\Cl(1,n{-}1)$ is type C if and only if $n \equiv 1 \pmod 4$, $\Cl(n{-}1,1)$ is type C if and only if $n \equiv 3 \pmod 4$, and in each case the other real form has complex type. The two real forms have the same complexification, which is split for every odd $n$: over $\C$ the element $\omega' = \omega$ or $i\omega$, whichever has square $+1$, is central, and $\tfrac12(1 \pm \omega')$ are its central idempotents. Thus the obstruction is present in every odd Lorentzian dimension; the real form decides whether it is visible as a splitting or hidden as a central complex structure.
\end{corollary}
\begin{proof}
$\Cl(1,n{-}1)$ is the spacelike extension of $\Cl(1,n{-}2)$, an even-dimensional algebra with $p - q = 3 - n$, which is $\equiv 2 \pmod 4$ iff $n \equiv 1 \pmod 4$ and $\equiv 0$ iff $n \equiv 3$; Theorem \ref{c:transport.where-obstruction-lies-commuting-square}(a) gives the first claim and the complex type of the other case. $\Cl(n{-}1,1)$ is the extension of $\Cl(n{-}2,1)$, $p - q = n - 3$, by a generator of square $+1$, obstructed iff $n - 3 \equiv 0$, i.e.\ $n \equiv 3 \pmod 4$, complex type otherwise. For the complexification: over $\C$ the substitution $e_i \mapsto ie_i$ on the generators of square $-1$ turns the relations of $\Cl(p,q)$ into those of $\Cl(n,0)$, so $\Cl(p,q) \otimes \C$ depends only on $n$ (\ref{c:std-algebra.complexification}); the volume element remains central, and if $\omega^2 = -1$ then $(i\omega)^2 = +1$, so $\omega'$ as defined is a central involution and $\tfrac12(1 \pm \omega')$ are orthogonal central idempotents summing to $1$. The script checks the alternation, with its complex-type complement, in dimensions $3, 5, 7, 9, 11$.

\end{proof}

\begin{remark}[Three dimensions]\label{c:note.transport.three-dimensions}
$n = 5$: $\Cl(1,4)$ is split, $\Cl(4,1)$ complex type. $n = 3$: $\Cl(2,1)$ is split, $\Cl(1,2)$ complex type. $n = 11$: $\Cl(10,1)$ is split, $\Cl(1,10)$ complex type. In the convention of this paper (time of square $+1$) the five-dimensional extension of spacetime is the split one; the planar and eleven-dimensional cases are split in the opposite convention. Physics done with complex gamma matrices sees the splitting in every odd dimension, as the corollary says it must; what the real form governs is the reality of the spinors.
\end{remark}

\section{Type P: the rule of no rule}

\phantomsection\label{c:note.transport.setting-type-p-rule-no-rule}
Do non-central idempotents demand an ascent? The answer is that type P demands nothing, and the reason is a one-line identity whose \emph{absence} is precisely what made type N hard. We state the repair theorem in the exact form of the no-go of the section ``The refuted conjecture, and the no-go'', so that the two answers can be compared clause by clause; then we develop the structure theory that shows the triviality is a theorem about the idempotent variety, not an evasion of it.

\subsection*{The repair theorem: cost zero}
\phantomsection\label{c:note.transport.setting-repair-theorem-cost-zero}
An idempotent $e \notin \{0,1\}$ is a genuine invertibility obstruction: $e(1-e) = 0$ exhibits it as a zero divisor, and it has no inverse in any unital extension (if $ex = 1$ then $e = e(ex) = e^2 x = ex = 1$). The transport question of the section ``The refuted conjecture, and the no-go'', transposed verbatim from the null cone to the idempotent variety, is therefore well posed: does some extension of the algebra carry an equivariant assignment of von Neumann pseudo-inverses to the idempotents? For null vectors the answer was no, in every extension. For idempotents:

\begin{proposition}[Type P is self-repairing]\label{c:transport.type-p-self-repairing}
Let $A$ be any unital ring, $\mathcal{E}(A)$ its set of idempotents, and $A' \supseteq A$ any unital extension. The identity assignment
\[
\sigma : \mathcal{E}(A) \to A' , \qquad \sigma(e) = e,
\]
satisfies both von Neumann identities, $e\,\sigma(e)\,e = e$ and $\sigma(e)\,e\,\sigma(e) = \sigma(e)$, and is equivariant under \emph{every} ring automorphism of every extension: for any automorphism $\theta$ of $A'$ preserving $A$, $\sigma(\theta(e)) = \theta(\sigma(e))$ trivially. In the terminology of the taxonomy: the transport rule for type P is the identity, at signature increment $(0,0)$, and the repair is already present in the base algebra.
\end{proposition}
\begin{proof}
$e\,e\,e = e^2 \cdot e = e \cdot e = e$, twice over; equivariance of the identity map is a tautology; and $\theta(e)$ is again an idempotent, so $\sigma$ is defined wherever equivariance asks it to be.

\end{proof}

\begin{remark}[Why the shear has nothing to act on: $eee = e$ versus $vvv = 0$]\label{c:note.transport.why-shear-has-nothing-act-eee}
The proposition looks too easy to be the analogue of Theorem \ref{c:transport.no-go-level-elements}, so it is worth locating exactly where the difficulty of type N disappears. It is not in the shear, which is present on both sides: the pseudo-inverses of a null vector form a torsor under the null rotations (Proposition \ref{c:transport.pseudo-inverse-torsor}), and the idempotents sharing the image of $e$ form one under $\mathrm{Hom}$-groups that the stabiliser of $e$ shears in exactly the same unipotent manner (Remark \ref{c:note.transport.moore-penrose-boundary-idempotent-side} below). The difference is what the torsor contains. The null vector is not a member of its own repair family, since $v\,v\,v = 0 \neq v$, so equivariance must select from a set with no invariant point; the idempotent is a member of its own, since $eee = e$, so the repair need not lie in the sheared directions at all. One identity separates the two cases. The taxonomy's three transport costs, namely a hyperbolic plane for N (and even then only at the level of points), one generator for C, and nothing for P, are thus not a graded scale of difficulty but three different answers to the same equivariance question, each carrying its own theorem.
\end{remark}

\subsection*{The idempotent variety: orbits, gauge, and the P/C boundary}
\phantomsection\label{c:note.transport.setting-idempotent-variety-orbits-gauge-p}
The repair theorem shows type P poses no transport problem; the following classification shows why it could not have: a type-P idempotent carries no invariant data beyond its rank, so any two candidates of the same rank are the same candidate up to gauge.

\begin{lemma}[Units span]\label{c:transport.units-span}
In a finite-dimensional unital $\R$-algebra $A$, every element is a sum of two invertible elements.
\end{lemma}
\begin{proof}
For $a \in A$ and $\lambda \in \R$, the element $a - \lambda 1$ is invertible iff the left-multiplication operator $L_a - \lambda\, \mathrm{id}$ on the finite-dimensional space $A$ is invertible (if $L_{a - \lambda 1}$ is bijective, the preimage of $1$ is a right inverse, and finite-dimensionality upgrades one-sided inverses to two-sided). The reals $\lambda$ for which $\det(L_a - \lambda\,\mathrm{id}) = 0$ are finitely many, so choose $\lambda$ with both $a - \lambda 1$ and $\lambda 1$ invertible ($\lambda \neq 0$): then $a = (a - \lambda 1) + \lambda 1$.

\end{proof}

\begin{theorem}[Classification of idempotent orbits]\label{c:transport.classification-idempotent-orbits}
\changedmark{gray}{reworded}
Let $A = \Cl(p,q)$, nondegenerate, so that $A \cong \Mat_m(D)$ or $\Mat_m(D) \oplus \Mat_m(D)$ with $D \in \{\R, \C, \HH\}$ (Theorem \ref{c:transport.classification-real-clifford-algebras}), acting on the free right $D$-module $V = D^m$ (or on the pair of them). Define the \emph{rank} of an idempotent as $\dim_D$ of its image (a rank pair, in the split case). Then:
\begin{enumerate}
\item Two idempotents are conjugate under the inner automorphism group of $A$ if and only if they have the same rank data. In the simple case the idempotents of rank $r$, $0 < r < m$, form a single conjugacy orbit; in the split case inner automorphisms act componentwise and the orbits are indexed by rank pairs.
\item $e$ and $1 - e$ are conjugate if and only if $\mathrm{rank}\, e = m - \mathrm{rank}\, e$ (componentwise, in the split case): the labelling of a complementary pair is gauge exactly when the pair is balanced, and is canonically ordered by rank otherwise.
\item An idempotent is central if and only if its conjugation orbit is a singleton. The taxonomy's P/C boundary is therefore orbit-theoretic: type C comprises the idempotents that conjugation cannot move at all, type P those it moves through a full rank class.
\end{enumerate}
\end{theorem}
\begin{proof}
(1) Conjugation preserves rank, since $geg^{-1}$ has image $g(\mathrm{im}\, e)$. Conversely let $e, e'$ be idempotents of equal rank $r$ on $V = D^m$, on which $A \cong \Mat_m(D)$ acts (\ref{c:transport.classification-real-clifford-algebras}). Each splits $V$: $V = \mathrm{im}\,e \oplus \ker e$, both summands free (a finitely generated submodule of a free module over a division ring is free, with well-defined dimension); choose $D$-bases of $\mathrm{im}\,e$ and $\ker e$, and likewise for $e'$. The $D$-linear map $g$ carrying the first adapted basis to the second is invertible and satisfies $g e g^{-1} = e'$, both sides being the projection onto $g(\mathrm{im}\,e) = \mathrm{im}\,e'$ along $g(\ker e) = \ker e'$. In the split case, invertible elements are pairs of invertibles and conjugation acts componentwise (inner automorphisms preserve the two ideals, which are the algebra's minimal central summands), so the same argument applies in each slot and no inner automorphism exchanges slots.

(2) is (1) applied to the pair, $\mathrm{rank}(1-e) = m - \mathrm{rank}(e)$.

(3) If $e$ is central its orbit is $\{e\}$. Conversely, if $geg^{-1} = e$ for every invertible $g$, then $e$ commutes with every invertible element, hence with every sum of two of them, hence with all of $A$ by Lemma \ref{c:transport.units-span}: $e$ is central. The identification of the singleton-orbit idempotents with the taxonomy's type C, and of the moved ones with type P, is then Definition \ref{c:obj.transport.obstruction-types} read against (1).

\end{proof}

\begin{proposition}[The pair is canonical; the label is a bit]\label{c:transport.pair-canonical-label-bit}
For an idempotent $e \notin \{0,1\}$, set $u = 2e - 1$, an involution ($u^2 = 1$), with $e = \tfrac12(1+u)$ and $1 - e = \tfrac12(1-u)$ its spectral projections. The unordered pair $\{e, 1-e\}$ and the involution $\pm u$ determine one another canonically; the residual choice (which member to call $e$, equivalently the sign of $u$) is a $\mathbb{Z}_2$ labelling. By Theorem \ref{c:transport.classification-idempotent-orbits}(2) this labelling is gauge (implemented by an inner automorphism) precisely for balanced pairs, and is canonically resolved by rank order otherwise. The pattern is now visible across the whole taxonomy: each obstruction type terminates in a residual $\mathbb{Z}_2$: the labelling of $\varepsilon_\pm$ for type C, exchangeable by no inner automorphism and rendered gauge only by the ascent (Proposition \ref{c:transport.ascent-renders-grade-involution-inner}); the sign $\gamma$ of the fused ascent for Rule 2, not detected by the algebra (Remark \ref{c:note.transport.bit-algebra-forgets}); and the labelling of $\{e, 1-e\}$ for type P, gauge already in the base algebra whenever the ranks permit. The three $\mathbb{Z}_2$'s are at strictly decreasing depth: outer, forgotten, inner.
\end{proposition}
\begin{proof}
Everything except the last two sentences is contained in the displayed formulas and Theorem \ref{c:transport.classification-idempotent-orbits}(2); the comparison across types collects Proposition \ref{c:transport.ascent-renders-grade-involution-inner}, Remark \ref{c:note.transport.bit-algebra-forgets} and Theorem \ref{c:transport.classification-idempotent-orbits}(2), each proved where cited.

\end{proof}

\begin{remark}[The Moore--Penrose boundary, idempotent side]\label{c:note.transport.moore-penrose-boundary-idempotent-side}
One refinement of the type-P question does carry content, and it traces the boundary of the section ``The hyperbolic atom as a matrix algebra: the Moore--Penrose pseudoinverse'' from the other side. Within the idempotents sharing a fixed image $W$, is there a canonical one? They are in bijection with the complements of $W$, hence form a torsor under $\mathrm{Hom}(V/W, W)$ on which the stabiliser of $W$ acts by unipotent transvections, the same shear mechanism as the null rotations of Lemma \ref{c:transport.group-null-rotations}. A definite Hermitian form resolves it, the orthogonal projection $P_W$ being the unique self-adjoint member; an indefinite one fails exactly on the degenerate subspaces, where no self-adjoint member exists at all: were $e$ self-adjoint with image $W$, any nonzero $z \in W \cap W^{\perp}$ would give $\langle x, z\rangle = \langle x, ez\rangle = \langle ex, z\rangle = 0$ for every $x$, against nondegeneracy of the ambient form. So the two sides agree: canonical selection is available exactly where a definite structure is. The transport-theoretic point is the division of labour. For type P the ambiguity concerns only \emph{which} idempotent to use, and the repair theorem needs no choice at all, since $\sigma(e) = e$ is asked about a given $e$; for type N the same ambiguity affects the repair itself, which is why the null cone needs the ascent and the idempotent variety does not.
\end{remark}

\begin{corollary}[No fourth rule]\label{c:transport.no-fourth-rule}
\changedmark{gray}{reworded}
The taxonomy of Definition \ref{c:obj.transport.obstruction-types} is transport-complete with two nontrivial rules. Type N: ascend by one hyperbolic plane, with the repair delivered at the level of points (Theorems \ref{c:transport.no-go-level-elements}, \ref{c:transport.totality}, \ref{c:transport.minimality}, \ref{c:transport.representation-independent-minimality-completion}). Type C: ascend by one generator, with the fused algebra unique up to isomorphism (Theorems \ref{c:transport.rule-2-ideal-fusion-sharp-form}, \ref{c:transport.one-generator-minimal-minimal-completions-classified}, \ref{c:transport.embedding-independent-minimality-fusion}). Type P: no ascent; the identity is an equivariant repair (Proposition \ref{c:transport.type-p-self-repairing}). By the exhaustiveness theorem (Theorem \ref{c:transport.exhaustiveness-every-zero-divisor-annihilated-idempotent}), every zero divisor of a nondegenerate real Clifford algebra falls under these three rules.
\end{corollary}
\begin{proof}
The source gives no separate proof: the corollary collects the theorems its statement cites. \ref{c:transport.embedding-independent-minimality-fusion}, \ref{c:transport.exhaustiveness-every-zero-divisor-annihilated-idempotent}, \ref{c:transport.minimality}, \ref{c:transport.no-go-level-elements}, \ref{c:transport.one-generator-minimal-minimal-completions-classified}, \ref{c:transport.representation-independent-minimality-completion}, \ref{c:transport.rule-2-ideal-fusion-sharp-form}, \ref{c:transport.totality}, \ref{c:transport.type-p-self-repairing}.

\end{proof}

\section{The mod-8 clock and the one-way door}

\phantomsection\label{c:note.transport.setting-mod-8-clock-one-way}
The two rules now proved, this section settles their bookkeeping: how transport moves the algebra's Morita class (it almost never does), and what survives descent (everything transported; nothing adjoined). Both answers turn on the same construction used for Rule 2: matrix units, this time supplied by the ascent itself.

\subsection*{Rule 1 is matrix amplification: the ladder pair as matrix units}
\phantomsection\label{c:note.transport.setting-rule-1-matrix-amplification-ladder}
Fix the completion $B = \Cl(p{+}1, q{+}1) = \Cl(\mathbb{V} \perp H)$ of Rule 1, with $H$ spanned by a hyperbolic pair normalised as in the section ``Null vectors, zero divisors, and hyperbolic planes'': null vectors $f, f^*$ with $f \cdot f^* = \tfrac12$, so that
\[
f^2 = (f^*)^2 = 0, \qquad ff^* + f^*f = 1, \qquad ff^*f = f, \quad f^*ff^* = f^* .
\]
(The conformal model's pair of the section ``The completion: the conformal model'' is the same plane in a different normalisation: $\nO \cdot \nInf = -1$, so $f = \nO$ and $f^* = -\tfrac12\nInf$ is a hyperbolic pair in the present sense.) Every vector of $H$ anticommutes with every vector of $\mathbb{V}$, so (as in Lemma \ref{c:transport.decomposition-commutation}) $f x = \hat{x} f$ and $f^* x = \hat{x} f^*$ for $x \in A = \Cl(p,q)$, and the \emph{even} elements of $\Cl(H)$ (in particular $ff^*$ and $f^*f$) commute with all of $A$. The blade basis gives the Peirce-ready decomposition
\[
B \;=\; A \,\oplus\, Af \,\oplus\, Af^* \,\oplus\, Aff^* .
\]

\begin{theorem}[Rule 1 is Morita-trivial]\label{c:transport.rule-1-morita-trivial}
\changedmark{gray}{reworded}
The elements
\[
E_{11} = ff^*, \qquad E_{22} = f^*f, \qquad E_{12} = f, \qquad E_{21} = f^*
\]
form a full $2 \times 2$ system of matrix units in $B$ ($E_{11} + E_{22} = 1$, $E_{ij}E_{kl} = \delta_{jk}E_{il}$); the corner is the base, $E_{11} B E_{11} = A\,ff^* \cong A$ via $a \mapsto a\,ff^*$; and therefore
\[
\Cl(p{+}1, q{+}1) \;\cong\; \Mat_2\bigl(\Cl(p,q)\bigr),
\]
under which the ascent embedding $A \hookrightarrow B$ becomes the twisted diagonal $a \mapsto \mathrm{diag}(a, \hat{a})$. In particular Rule 1 preserves the Morita class: the Wedderburn type, and with it the simplicity or splitness and the type-C diagnosis of the section ``A taxonomy of obstructions'', are the same for $B$ as for $A$, and the clock residue $d = p - q \pmod 8$ is fixed. (The \emph{presence} of types N and P is not preserved and could not be: Proposition \ref{c:transport.no-transport-terminates} below shows every completion carries fresh obstructions of its own.) The adjoined ladder pair is, in the matrix description, exactly the off-diagonal matrix units, and the descent projectors $f^*f,\, ff^*$ are the diagonal ones.
\end{theorem}
\begin{proof}
The matrix-unit relations are the hyperbolic-pair identities $v^2 = w^2 = 0, \qquad vw + wv = 1, \qquad \text{hence} \quad vwv = (1 - wv)v = v$ read off directly: $E_{12}E_{21} = ff^* = E_{11}$, $E_{21}E_{12} = f^*f = E_{22}$, $E_{12}^2 = f^2 = 0 = E_{21}^2$, $E_{11}E_{12} = ff^*f = f = E_{12}$, $E_{12}E_{11} = f\,ff^* = 0$, and so on; $E_{11} + E_{22} = 1$ is the pair's anticommutator. For the corner: using the displayed decomposition of $B$ and $ff^* a = a\, ff^*$ (evenness), one computes $ff^* (a_0 + a_1 f + a_2 f^* + a_3 ff^*)\, ff^* = (a_0 + a_3) ff^*$, since $ff^* \cdot f \cdot ff^* = f f^* Q(f) f^* = f \cdot ff^* = f^2 f^* = 0$ and likewise the $f^*$ term vanishes, while the $1$ and $ff^*$ terms remain; so $E_{11}BE_{11} = A\,ff^*$. The map $a \mapsto a\,ff^*$ is an algebra isomorphism onto it: multiplicative since $(a\,ff^*)(b\,ff^*) = ab\,(ff^*)^2 = ab\,ff^*$, unital onto the corner unit $ff^*$, injective by the directness of the decomposition. The Peirce isomorphism of Lemma \ref{c:transport.matrix-units-peirce-isomorphism}, whose proof uses only the matrix-unit relations, then gives $B \cong \Mat_2(A)$. The embedding: for $a \in A$, $\varphi(a)_{11} = E_{11} a E_{11} = a\,ff^*$ and $\varphi(a)_{22} = E_{12}\, a\, E_{21} = f a f^* = \hat{a}\, ff^*$, while the off-diagonal entries vanish ($E_{11} a E_{21} = a\, ff^*f^* = 0$, $E_{12} a E_{11} = \hat a\,f ff^* = 0$); under the corner identification this is $\mathrm{diag}(a, \hat{a})$, which is multiplicative since the grade involution is an automorphism (\ref{c:std-algebra.involution-properties}). Morita-class preservation, and with it the invariance of the type-C diagnosis, is the statement $\Mat_2(A) \sim_{\text{Morita}} A$ (\ref{c:std-algebra.morita-amplification}) together with the fact that the Wedderburn data of $\Mat_2(\Mat_m(D))$ (or of the split case, blockwise) is that of $A$ with $m$ doubled.

\end{proof}

\begin{remark}[Prior art, and what is new here]\label{c:note.transport.prior-art-what-new-here}
The isomorphism $\Cl(p{+}1,q{+}1) \cong \Mat_2(\Cl(p,q))$ is classical: it is the algebraic form of the $(1,1)$-periodicity of Clifford algebras \cite{ABS1964, Porteous1995, Lounesto2001}, usually obtained from $\Cl(p{+}1,q{+}1) \cong \Cl(p,q)\,\hat{\otimes}\,\Cl(1,1)$ and $\Cl(1,1) \cong \Mat_2(\R)$. What we claim as expository contribution is the presentation: the matrix units \emph{are} the transported ladder pair and its occupation projectors, so the periodicity isomorphism is not auxiliary to the transport principle but is the transport principle's own bookkeeping: the ascent adjoins exactly one row and column's worth of new structure, and the relations $f^2 = 0$, $ff^* + f^*f = 1$ are the matrix-unit relations of that adjunction. The parallel with Rule 2 is now exact (Remark \ref{c:note.transport.2-pattern-two-rules-compared}): both rules amplify by $\Mat_2$, Rule 1 over the whole base, Rule 2 over its even half, and in both cases the matrix units are built from the completion's own canonical elements: the ladder pair for N, the central projectors plus the adjoined generator for C.
\end{remark}

\begin{remark}[Iteration]\label{c:note.transport.iteration}
Iterating Theorem \ref{c:transport.rule-1-morita-trivial} gives $\Cl(p{+}k, q{+}k) \cong \Mat_{2^k}(\Cl(p,q))$: the tower over any base is Morita-stationary from its first step. The structural consequences of the iterated statement (tower rigidity, the one-mode-per-rung reading, and the direct limit) are the subject of the series' tower paper \cite{ButlerT3} and are not developed here.
\end{remark}

\subsection*{The state machine becomes stationary}
\phantomsection\label{c:note.transport.setting-state-machine-becomes-stationary}
The bookkeeping of the two rules assembles into the state machine of a figure of the original, whose one theorem is the following.

\begin{proposition}[Fusion occurs at most once]\label{c:transport.fusion-occurs-most-once}
\changedmark{gray}{reworded}
Call a \emph{transport history} any finite sequence of algebras beginning at a nondegenerate $\Cl(p,q)$, each obtained from its predecessor by Rule 1 (the $(1,1)$ ascent) or, when the predecessor carries type C, by Rule 2 (ideal fusion). Then along any transport history:
\begin{enumerate}
\item the clock residue $d \bmod 8$ is constant under Rule 1 (Theorem \ref{c:transport.rule-1-morita-trivial}) and moves by $\pm 1$ under Rule 2;
\item type C is present exactly at $d \equiv 1 \pmod 4$, i.e.\ at residues $1$ and $5$ (Theorem \ref{c:transport.existence-classification-type-c}), and Rule 2 maps these to residues in $\{0, 2\}$ and $\{4, 6\}$ respectively, none of which is $\equiv 1 \pmod 4$;
\item hence Rule 2 can be applied at most once in any transport history, and after it (or from the start, if $d \not\equiv 1 \pmod 4$) every subsequent algebra in the history has the same Wedderburn type: the state machine takes at most one step on the mod-8 clock, and then remains fixed.
\end{enumerate}
In particular the physical lineage is stationary from the start: $\Cl(1,3)$ has $d \equiv 6$, is simple, never demands Rule 2, and every Rule-1 ascent $\Cl(1{+}k, 3{+}k) \cong \Mat_{2^k}(\Cl(1,3))$ remains at residue $6$, quaternionic type.
\end{proposition}
\begin{proof}
(1) is Theorem \ref{c:transport.rule-1-morita-trivial} and the count $d \mapsto d \pm 1$ for one adjoined generator. (2) is Theorem \ref{c:transport.existence-classification-type-c} with the arithmetic $1 \pm 1 = 0, 2$ and $5 \pm 1 = 4, 6$ mod $8$. (3): once the residue leaves $\{1, 5\}$ it never returns, since Rule 1 fixes it and Rule 2 is only ever invoked (indeed only ever \emph{defined}, as a transport demand) at $\{1,5\}$; Morita-stationarity after the last (hence only) fusion is Theorem \ref{c:transport.rule-1-morita-trivial} again. The closing statement on $\Cl(1,3)$ reads its residue $6$, simple and quaternionic, off \ref{c:transport.classification-real-clifford-algebras}.

\end{proof}

\subsection*{The one-way door: what descent erases is what ascent added}
\phantomsection\label{c:note.transport.setting-one-way-door-what-descent}
The asymmetry of transport (total on the way up, lossy on the way down) is now made exact in the matrix picture of Theorem \ref{c:transport.rule-1-morita-trivial}. Throughout, $P = f^*f$ and $1 - P = ff^*$ are the \emph{descent projectors}, the diagonal matrix units; they are the spectral projections $\tfrac12(1 \pm \sigma)$ of the involution $\sigma = f^*f - ff^*$, $\sigma^2 = 1$, canonically attached to the completed pair.

\begin{proposition}[The pinching, and its kernel]\label{c:transport.pinching-kernel}
Define the pinching $E : B \to B$, $E(x) = PxP + (1-P)x(1-P)$. Then:
\begin{enumerate}
\item $E$ is an idempotent linear map whose image is the commutant $C = \{x \in B : xP = Px\} = A \oplus A\,ff^*$, on which it restricts to the identity, and $E$ is a $C$-bimodule map: $E(cxc') = c\,E(x)\,c'$ for $c, c' \in C$, with $E(1) = 1$: a conditional expectation onto $C$.
\item Its kernel is exactly the adjoined ladder sector:
\[
\ker E \;=\; PB(1-P) \,\oplus\, (1-P)BP \;=\; A f^* \,\oplus\, A f .
\]
In particular $E(f) = E(f^*) = 0$: the pinching erases precisely the pair the ascent adjoined, and nothing else.
\item The round trip is exact on transported elements: for $a \in A$, $E(a) = a$, and the corner read-out $D(x) = PxP$, identified with an $A$-valued map through the natural corner isomorphism $E_{22}BE_{22} = A\,f^*f \cong A$ ($a \mapsto a\,f^*f$), satisfies $D(a) = a$: ascent-then-descent is the \emph{identity} on the base. (In the Peirce coordinates of Theorem \ref{c:transport.rule-1-morita-trivial}, which identify the $(2,2)$ block with the $(1,1)$ corner by conjugating with $E_{12} = f$, the same round trip reads as the grade involution $\hat{a}$; the two corner identifications differ by exactly the grade twist of the embedding $a \mapsto \mathrm{diag}(a,\hat a)$, and either way the map is an automorphism, faithful.) The multiplicative defect of $D$ on general elements is the exact interposition formula
\[
D(xy) - D(x)D(y) \;=\; P\,x\,(1-P)\,y\,P .
\]
\end{enumerate}
\end{proposition}
\begin{proof}
(1) $E^2 = E$ and $E(1) = P + (1-P) = 1$ are immediate from $P^2 = P$. If $xP = Px$ then the cross-blocks of $x$ vanish ($Px(1-P) = xP(1-P) = 0$ and symmetrically), so $x = PxP + (1-P)x(1-P) = E(x)$; conversely $E(x)$ commutes with $P$ for any $x$; hence $\mathrm{im}\,E = C$ and $E|_C = \mathrm{id}$. The bimodule property follows by moving $c, c'$ through the projectors they commute with. The identification $C = A \oplus A\,ff^*$: $A$ commutes with $P$ (evenness of $f^*f$ in the $H$-factor) and so does $ff^*$, giving $A \oplus A\,ff^* \subseteq C$; and the two diagonal Peirce blocks $PBP \oplus (1-P)B(1-P) = A\,f^*f \oplus A\,ff^*$ span the same $2 \cdot 2^n$-dimensional space, since $a\,f^*f = a - a\,ff^*$. (2) The off-diagonal Peirce blocks contain $Af$ and $Af^*$ ($f = ff^*\cdot f = f \cdot f^*f$ places $f \in (1-P)BP$, symmetrically $f^* \in PB(1-P)$), and the dimensions ($2^n$ each, against Peirce blocks of dimension $2^n$ each in $\Mat_2(A)$) force equality. (3) $E(a) = a$ since $a \in C$. $D(a) = P a P = a\,f^*f$ (as $a$ commutes with $P$ and $P^2 = P$), which is $a$ under the natural corner isomorphism. The Peirce-coordinate reading: the $(2,2)$ entry of $\varphi(a)$ transported to the $E_{11}$ corner is $E_{12}\,a\,E_{21} = f a f^* = \hat{a}\,ff^*$, i.e.\ $\hat{a}$, as computed in Theorem \ref{c:transport.rule-1-morita-trivial}. The defect formula is $D(xy) = Px(P + (1-P))yP = D(x)D(y) + Px(1-P)yP$.

\end{proof}

\begin{theorem}[The door is one-way]\label{c:transport.door-one-way}
Let $A = \Cl(p,q)$ be nondegenerate and let $B$ be either minimal completion of this paper: the Rule-1 ascent $\Cl(p{+}1,q{+}1) \cong \Mat_2(A)$, or (when $p - q \equiv 1 \bmod 4$) a Rule-2 ascent $\cong \Mat_2(\Cl^0(p,q))$. Then there is no unital algebra homomorphism $B \to A$ whatsoever. Every descent is therefore non-multiplicative; the conditional expectation of Proposition \ref{c:transport.pinching-kernel} is the canonical linear surrogate, and its loss, the erasure of the ladder sector, is forced by the algebra, not chosen by the formalism.
\end{theorem}
\begin{proof}
$B$ is semisimple (\ref{c:transport.classification-real-clifford-algebras}), so the kernel of a homomorphism $\theta : B \to A$ is a sum of Wedderburn blocks of $B$, and $\theta$ is injective on the complementary blocks; unitality ($\theta(1) = 1 \neq 0$) forces at least one block to survive. It suffices then that every block of $B$ has real dimension exceeding $\dim_\R A$. Rule 1, $A$ simple: $B \cong \Mat_2(A)$ (\ref{c:transport.rule-1-morita-trivial}) is a single block of dimension $4 \dim A > \dim A$. Rule 1, $A = S \oplus S$ split: $\Mat_2(A) \cong \Mat_2(S) \oplus \Mat_2(S)$, blocks of dimension $4 \dim S = 2 \dim A > \dim A$. Rule 2 ($A$ split): $B \cong \Mat_2(\Cl^0)$ is simple (Theorem \ref{c:transport.rule-2-ideal-fusion-sharp-form}) of dimension $2^{n+1} = 2\dim A > \dim A$. In every case no block embeds in $A$, contradiction.

\end{proof}

\begin{remark}[The two-sided reading]\label{c:note.transport.two-sided-reading}
Proposition \ref{c:transport.pinching-kernel} and Theorem \ref{c:transport.door-one-way} split the loose statement that descent is lossy into its two true halves. On transported elements the round trip is not lossy at all: it is the identity in the natural corner identification (the grade involution in Peirce coordinates; either way invertible), and the base algebra passes through the door unchanged in both directions. What cannot descend is the completion's \emph{new} content, the ladder pair and everything it generates off the diagonal, and this impossibility is a theorem about unital homomorphisms, not an artefact of choosing a particular descent map. The statement the computations support is sharper than the original: \emph{descent is invertible on the transported elements and impossible on the elements the ascent adjoined}. Compare the point-level statement of Lemma \ref{c:transport.descent-read-out-map}: there the descent is exact on the whole invertible locus and fails only on the ideal points, which are precisely the values the ascent added. Element-level and point-level descent lose different things, and each loses exactly what its ascent supplied.
\end{remark}

\begin{proposition}[No transport terminates]\label{c:transport.no-transport-terminates}
No algebra reachable by a nonempty transport history is free of zero divisors. Every Rule-1 completion has mixed signature and so carries type N; indeed the adjoined pair $f, f^*$ is itself null: the repair of the old null cone is effected by adjoining new members of it. Every Rule-2 completion is a matrix algebra $\Mat_2(\Cl^0(p,q))$ and so carries type P: the fusion's own matrix units supply the non-central idempotents and nilpotents. Transport never reaches a division algebra; each completion carries fresh obstructions of its own.
\end{proposition}
\begin{proof}
Rule 1 adjoins a hyperbolic plane, so the completed signature $(p{+}1, q{+}1)$ is mixed whatever the base, and Proposition \ref{c:transport.existence-locus-type-n} applies; concretely $f^2 = 0$ with $f \neq 0$. Rule 2's completion contains the matrix units of Lemma \ref{c:transport.matrix-units-peirce-isomorphism}: $E_{11}$ is a non-central idempotent (Theorem \ref{c:transport.rule-2-ideal-fusion-sharp-form}(2)) and $E_{12}$ is a nonzero nilpotent.

\end{proof}

\begin{conjecture}[Obstructions are moved, never destroyed]\label{c:transport.obstructions-are-moved-never-destroyed}
Proposition \ref{c:transport.no-transport-terminates} is the zero-dimensional instance of a principle the transport programme repeatedly meets: a completion never removes degeneracy, it relocates it. We conjecture the following family form. For any family of quadratic forms degenerating tamely along a locus (in the sense of the series' sheaf-theoretic paper \cite{ButlerT4}), any completion of the family that is nondegenerate of constant signature on one side of the locus degenerates on a displaced locus on the other side; no completion is nondegenerate globally. The one-sided completion theorem and shadow wall of that paper prove exactly this for tame corank-one degenerations; the conjecture asserts the pattern is not an artefact of the tame case, and asks for the cohomological invariant that measures it, presumably a class built from the mod-8 residues on the two sides of the wall, the state machine of Proposition \ref{c:transport.fusion-occurs-most-once} evaluated in families.
\end{conjecture}

\section{The transport principle, and the algorithm}

\begin{definition}[Transport]\label{c:obj.transport.transport}
The \emph{transport} of a nondegenerate $(\mathbb{V}, Q)$ is the pair consisting of the minimal completion $\Cl(p{+}1,q{+}1)$ of Theorem \ref{c:transport.minimality} and the null-ray representation $X(x) \;=\; \nO + x + \tfrac12 Q(x)\, \nInf$; the \emph{transported reciprocal} is the total involution $\iota$ of Theorem \ref{c:transport.totality}. For a type-C obstruction the corresponding transport is the one-generator ascent of Theorem \ref{c:transport.rule-2-ideal-fusion-sharp-form}; for type P it is the identity.
\end{definition}

\begin{principle}[Transport principle]\label{c:transport.transport-principle}
Division fails in a real Clifford algebra of nondegenerate form in exactly three ways, and each carries a forced minimal repair.
\begin{itemize}
\item \emph{At an isotropic vector} the repair is impossible for the element and total for the point. No equivariant element-level repair exists in the Clifford algebra of any extension (Theorems \ref{c:transport.no-go-level-elements}, \ref{c:transport.cohomological-form-no-go}); the reciprocal nevertheless extends totally, continuously and uniquely on the space of points, in a completion whose increment $(+1,+1)$ is forced both as a lift and as an ambient quadric (Theorems \ref{c:transport.totality}, \ref{c:transport.minimality}, \ref{c:transport.representation-independent-minimality-completion}).
\item \emph{At a central idempotent} the repair is available at the level of elements, one dimension up: the adjoined generator renders the grade involution inner, the ideals fuse, and the completion is unique up to isomorphism and minimal among all simple real algebras (Theorems \ref{c:transport.rule-2-ideal-fusion-sharp-form}, \ref{c:transport.one-generator-minimal-minimal-completions-classified}, \ref{c:transport.embedding-independent-minimality-fusion}).
\item \emph{At a non-central idempotent} no repair is needed: the element is its own (Proposition \ref{c:transport.type-p-self-repairing}).
\end{itemize}
The trichotomy is exhaustive (Theorem \ref{c:transport.exhaustiveness-every-zero-divisor-annihilated-idempotent}), so there is no fourth case (Corollary \ref{c:transport.no-fourth-rule}); in each case the cost is independent of how the algebra is represented, and is exactly what the obstruction requires. \emph{Complete the space when the obstruction is a direction, the algebra when it is a splitting, and nothing when the element is already its own repair.}
\end{principle}
\begin{proof}
Division fails in a real Clifford algebra of nondegenerate form in exactly three ways, and each carries a forced minimal repair.
\begin{itemize}
\item \emph{At an isotropic vector} the repair is impossible for the element and total for the point. No equivariant element-level repair exists in the Clifford algebra of any extension (Theorems \ref{c:transport.no-go-level-elements}, \ref{c:transport.cohomological-form-no-go}); the reciprocal nevertheless extends totally, continuously and uniquely on the space of points, in a completion whose increment $(+1,+1)$ is forced both as a lift and as an ambient quadric (Theorems \ref{c:transport.totality}, \ref{c:transport.minimality}, \ref{c:transport.representation-independent-minimality-completion}).
\item \emph{At a central idempotent} the repair is available at the level of elements, one dimension up: the adjoined generator renders the grade involution inner, the ideals fuse, and the completion is unique up to isomorphism and minimal among all simple real algebras (Theorems \ref{c:transport.rule-2-ideal-fusion-sharp-form}, \ref{c:transport.one-generator-minimal-minimal-completions-classified}, \ref{c:transport.embedding-independent-minimality-fusion}).
\item \emph{At a non-central idempotent} no repair is needed: the element is its own (Proposition \ref{c:transport.type-p-self-repairing}).
\end{itemize}
The trichotomy is exhaustive (Theorem \ref{c:transport.exhaustiveness-every-zero-divisor-annihilated-idempotent}), so there is no fourth case (Corollary \ref{c:transport.no-fourth-rule}); in each case the cost is independent of how the algebra is represented, and is exactly what the obstruction requires. \emph{Complete the space when the obstruction is a direction, the algebra when it is a splitting, and nothing when the element is already its own repair.}

\end{proof}

\phantomsection\label{c:note.transport.setting-transport-principle-algorithm}
The Principle is a summary, not an axiom: each of its clauses is one of the theorems proved above, and it adds nothing to them. What it does add is the reading. The three costs are not a graded scale of difficulty but three different answers to one equivariance question, and the question is answerable at all only because the taxonomy is exhaustive: it is the completeness of Theorem \ref{c:transport.exhaustiveness-every-zero-divisor-annihilated-idempotent}, not the ingenuity of any rule, that turns a family of repairs into a procedure.

Two closing observations. First, the shape of the type-N answer is the surprising one, and it is worth stating plainly that it is a change of level rather than a technical fix: what is repaired is not the element, which remains non-invertible in every enlargement however large, but the space in which the element's division was asked. Second, no transport removes all obstructions. Every completion carries obstructions of its own, and in the type-N case the adjoined pair is itself null (Proposition \ref{c:transport.no-transport-terminates}): the repair of the old null cone is effected by adjoining new members of it. What the bookkeeping of the section ``The mod-8 clock and the one-way door'' shows is therefore not that transport terminates but that it stabilises, the clock moving at most once and then remaining fixed. Obstruction is conserved; what the rules achieve is that it no longer changes the Wedderburn type.

The paper is intended to be independently checkable throughout: every algebraic statement can be verified from the axiom $v^2 = Q(v)$ with pencil and paper, the topological statements of the section ``Totality of the transported reciprocal'' from the elementary point-set facts recorded in the section ``The completion: the conformal model'', and every computational identity by the deposited scripts.

\begin{theorem}[The transport algorithm]\label{c:transport.transport-algorithm}
\changedmark{gray}{reworded}
The Principle says what is true; practice needs to know what to do. The theorems distil into a six-step procedure, and every step names the theorem, proved above, that licenses it. the table accompanying \ref{c:transport.transport-algorithm} is the one-page version: the theorems in operational form. That the steps form a decision procedure with no gaps is itself a theorem, Corollary \ref{c:transport.no-fourth-rule}.




\begin{enumerate}
\item[\textbf{Step 1.}] \textbf{Locate the failure.} Identify the element at which invertibility fails, or, if the construction is a family, the locus at which it fails. Nothing later depends on how this is done; it is a computation in the case at hand. What the algebra guarantees is that the failure cannot be argued away: a null vector has no inverse in \emph{any} unital algebra containing the base (Lemma \ref{c:transport.noinverse}), and neither has an idempotent.

\item[\textbf{Step 2.}] \textbf{Diagnose the type.} Ask, in order:
  \begin{itemize}
  \item Is the failing element a vector $v$ with $Q(v) = 0$? Then the obstruction is \textbf{type N}. These exist exactly in mixed signature, $pq \neq 0$ (Proposition \ref{c:transport.existence-locus-type-n}).
  \item Is it a \emph{central} idempotent? Then \textbf{type C}. These exist exactly when $p - q \equiv 1 \pmod 4$, and are then exactly the two projectors $\tfrac12(1 \pm \omega)$ (Theorem \ref{c:transport.existence-classification-type-c}); equivalently, exactly when the algebra is not simple (Theorem \ref{c:transport.classification-real-clifford-algebras}).
  \item Otherwise \textbf{type P}: a non-central idempotent, or an element annihilated by one.
  \end{itemize}
  The trichotomy is exhaustive: every zero divisor of a nondegenerate real Clifford algebra is annihilated by a nontrivial idempotent, possibly central in the split signatures, and then exactly when the element lies in one ideal (Theorem \ref{c:transport.exhaustiveness-every-zero-divisor-annihilated-idempotent}). Two cautions. The nilpotency locus is strictly larger than the type-N locus, and nilpotence alone does not diagnose N (Example \ref{c:note.transport.nilpotents-without-null-vectors-two-loci}); and a single algebra may carry all three types at once (Example \ref{c:note.transport.cl-1-0-cl-2-1}).

\item[\textbf{Step 3.}] \textbf{Apply the rule.}
  \begin{itemize}
  \item \textbf{Type N, Rule 1.} Adjoin one hyperbolic plane: $\Cl(p,q) \to \Cl(p{+}1,q{+}1)$. Do not expect the element to be repaired: no equivariant assignment of a partner to $v$ exists, in this ascent or any other (Theorem \ref{c:transport.no-go-level-elements}), and the obstruction is a cohomology class that every enlargement pulls back unchanged (Theorem \ref{c:transport.cohomological-form-no-go}). What is repaired is the space of points (Step 4). The increment is forced: $(+1,+1)$, never $(+2,0)$ or $(0,+2)$, among null-lift completions (Theorem \ref{c:transport.minimality}), and no algebra with a smaller or differently balanced quadric can host the resolution however the base is embedded in it (Theorem \ref{c:transport.representation-independent-minimality-completion}).
  \item \textbf{Type C, Rule 2.} Adjoin one generator $\nu$ orthogonal to $\mathbb{V}$, of either nonzero square. The grade involution becomes inner, the two ideals fuse, and the completion is $\Mat_2(\Cl^0(p,q))$, simple (Theorem \ref{c:transport.rule-2-ideal-fusion-sharp-form}). The null-generator ascent is not an option: it produces a square-zero ideal and never fuses (Theorem \ref{c:transport.one-generator-minimal-minimal-completions-classified}(2)). The sign of $\nu^2$ is free and the algebra forgets it (Remark \ref{c:note.transport.bit-algebra-forgets}); if it matters for the application, it must be fixed by structure outside the algebra.
  \item \textbf{Type P.} Do nothing. The identity assignment $\sigma(e) = e$ is a total, fully equivariant von Neumann repair, at signature increment $(0,0)$ (Proposition \ref{c:transport.type-p-self-repairing}).
  \end{itemize}

\item[\textbf{Step 4.}] \textbf{Compute in the completion.} For type N this is the substantive step, and it is where the classical formulas apply. Represent points, not elements: lift $x \in \mathbb{V}$ to the null ray of
  \[
  X(x) \;=\; \nO + x + \tfrac12 Q(x)\,\nInf ,
  \]
  which satisfies the incidence condition $Q(x) = 2st$ in the homogeneous coordinates $(s,x,t)$ of Proposition \ref{c:transport.swap-rescale-normal-form}. The reciprocal is then the linear involution
  \[
  \iota(s,\,x,\,t) \;=\; (2t,\ x,\ \tfrac12 s),
  \]
  total on the quadric and never failing (Theorem \ref{c:transport.totality}). All nonlinearity and all division have been deferred to a single read-out. This is the linearise-then-normalise pattern of homogeneous coordinates and M\"obius matrices, encountered here as forced rather than chosen (Remark \ref{c:note.transport.classical}).

\item[\textbf{Step 5.}] \textbf{Descend, and state what descent costs.} Return the answer to $\mathbb{V}$ with the read-out $\dsc[\,s\nO + x + t\nInf\,] = x/s$. On the invertible locus the composite gives exactly the classical answer (Lemma \ref{c:transport.descent-read-out-map}(2)); where the first slot vanishes the answer is an ideal point and corresponds to no point of $\mathbb{V}$ (Lemma \ref{c:transport.descent-read-out-map}(3)). \emph{The refusal of the read-out is information, not failure}: it says the answer exists and lies outside the original space. At the level of elements the corresponding statement is stronger and worth knowing before relying on any descent: the pinching onto the base is a conditional expectation whose kernel is exactly the adjoined ladder sector (Proposition \ref{c:transport.pinching-kernel}), and there is no unital algebra homomorphism from either completion back to its base at all (Theorem \ref{c:transport.door-one-way}). Descent is invertible on the transported elements and impossible on the elements the ascent adjoined.

\item[\textbf{Step 6.}] \textbf{Iterate, or stop.} Ask whether the completion carries an obstruction of its own. It does: no transport terminates in a division algebra (Proposition \ref{c:transport.no-transport-terminates}). But the bookkeeping is finite in the following sense: Rule 1 fixes the mod-8 residue, so the taxonomy of the completion is that of the base and iterating changes nothing structural (Theorem \ref{c:transport.rule-1-morita-trivial}), while Rule 2 moves the residue off the split classes and can therefore occur at most once in any transport history (Proposition \ref{c:transport.fusion-occurs-most-once}). A transport history takes at most one step on the clock, and then remains fixed.
\end{enumerate}
\end{theorem}
\begin{proof}
Each step names the result that licenses it; the derivation is those results.

The Principle says what is true; practice needs to know what to do. The theorems distil into a six-step procedure, and every step names the theorem, proved above, that licenses it. the table accompanying \ref{c:transport.transport-algorithm} is the one-page version: the theorems in operational form. That the steps form a decision procedure with no gaps is itself a theorem, Corollary \ref{c:transport.no-fourth-rule}.




\begin{enumerate}
\item[\textbf{Step 1.}] \textbf{Locate the failure.} Identify the element at which invertibility fails, or, if the construction is a family, the locus at which it fails. Nothing later depends on how this is done; it is a computation in the case at hand. What the algebra guarantees is that the failure cannot be argued away: a null vector has no inverse in \emph{any} unital algebra containing the base (Lemma \ref{c:transport.noinverse}), and neither has an idempotent.

\item[\textbf{Step 2.}] \textbf{Diagnose the type.} Ask, in order:
  \begin{itemize}
  \item Is the failing element a vector $v$ with $Q(v) = 0$? Then the obstruction is \textbf{type N}. These exist exactly in mixed signature, $pq \neq 0$ (Proposition \ref{c:transport.existence-locus-type-n}).
  \item Is it a \emph{central} idempotent? Then \textbf{type C}. These exist exactly when $p - q \equiv 1 \pmod 4$, and are then exactly the two projectors $\tfrac12(1 \pm \omega)$ (Theorem \ref{c:transport.existence-classification-type-c}); equivalently, exactly when the algebra is not simple (Theorem \ref{c:transport.classification-real-clifford-algebras}).
  \item Otherwise \textbf{type P}: a non-central idempotent, or an element annihilated by one.
  \end{itemize}
  The trichotomy is exhaustive: every zero divisor of a nondegenerate real Clifford algebra is annihilated by a nontrivial idempotent, possibly central in the split signatures, and then exactly when the element lies in one ideal (Theorem \ref{c:transport.exhaustiveness-every-zero-divisor-annihilated-idempotent}). Two cautions. The nilpotency locus is strictly larger than the type-N locus, and nilpotence alone does not diagnose N (Example \ref{c:note.transport.nilpotents-without-null-vectors-two-loci}); and a single algebra may carry all three types at once (Example \ref{c:note.transport.cl-1-0-cl-2-1}).

\item[\textbf{Step 3.}] \textbf{Apply the rule.}
  \begin{itemize}
  \item \textbf{Type N, Rule 1.} Adjoin one hyperbolic plane: $\Cl(p,q) \to \Cl(p{+}1,q{+}1)$. Do not expect the element to be repaired: no equivariant assignment of a partner to $v$ exists, in this ascent or any other (Theorem \ref{c:transport.no-go-level-elements}), and the obstruction is a cohomology class that every enlargement pulls back unchanged (Theorem \ref{c:transport.cohomological-form-no-go}). What is repaired is the space of points (Step 4). The increment is forced: $(+1,+1)$, never $(+2,0)$ or $(0,+2)$, among null-lift completions (Theorem \ref{c:transport.minimality}), and no algebra with a smaller or differently balanced quadric can host the resolution however the base is embedded in it (Theorem \ref{c:transport.representation-independent-minimality-completion}).
  \item \textbf{Type C, Rule 2.} Adjoin one generator $\nu$ orthogonal to $\mathbb{V}$, of either nonzero square. The grade involution becomes inner, the two ideals fuse, and the completion is $\Mat_2(\Cl^0(p,q))$, simple (Theorem \ref{c:transport.rule-2-ideal-fusion-sharp-form}). The null-generator ascent is not an option: it produces a square-zero ideal and never fuses (Theorem \ref{c:transport.one-generator-minimal-minimal-completions-classified}(2)). The sign of $\nu^2$ is free and the algebra forgets it (Remark \ref{c:note.transport.bit-algebra-forgets}); if it matters for the application, it must be fixed by structure outside the algebra.
  \item \textbf{Type P.} Do nothing. The identity assignment $\sigma(e) = e$ is a total, fully equivariant von Neumann repair, at signature increment $(0,0)$ (Proposition \ref{c:transport.type-p-self-repairing}).
  \end{itemize}

\item[\textbf{Step 4.}] \textbf{Compute in the completion.} For type N this is the substantive step, and it is where the classical formulas apply. Represent points, not elements: lift $x \in \mathbb{V}$ to the null ray of
  \[
  X(x) \;=\; \nO + x + \tfrac12 Q(x)\,\nInf ,
  \]
  which satisfies the incidence condition $Q(x) = 2st$ in the homogeneous coordinates $(s,x,t)$ of Proposition \ref{c:transport.swap-rescale-normal-form}. The reciprocal is then the linear involution
  \[
  \iota(s,\,x,\,t) \;=\; (2t,\ x,\ \tfrac12 s),
  \]
  total on the quadric and never failing (Theorem \ref{c:transport.totality}). All nonlinearity and all division have been deferred to a single read-out. This is the linearise-then-normalise pattern of homogeneous coordinates and M\"obius matrices, encountered here as forced rather than chosen (Remark \ref{c:note.transport.classical}).

\item[\textbf{Step 5.}] \textbf{Descend, and state what descent costs.} Return the answer to $\mathbb{V}$ with the read-out $\dsc[\,s\nO + x + t\nInf\,] = x/s$. On the invertible locus the composite gives exactly the classical answer (Lemma \ref{c:transport.descent-read-out-map}(2)); where the first slot vanishes the answer is an ideal point and corresponds to no point of $\mathbb{V}$ (Lemma \ref{c:transport.descent-read-out-map}(3)). \emph{The refusal of the read-out is information, not failure}: it says the answer exists and lies outside the original space. At the level of elements the corresponding statement is stronger and worth knowing before relying on any descent: the pinching onto the base is a conditional expectation whose kernel is exactly the adjoined ladder sector (Proposition \ref{c:transport.pinching-kernel}), and there is no unital algebra homomorphism from either completion back to its base at all (Theorem \ref{c:transport.door-one-way}). Descent is invertible on the transported elements and impossible on the elements the ascent adjoined.

\item[\textbf{Step 6.}] \textbf{Iterate, or stop.} Ask whether the completion carries an obstruction of its own. It does: no transport terminates in a division algebra (Proposition \ref{c:transport.no-transport-terminates}). But the bookkeeping is finite in the following sense: Rule 1 fixes the mod-8 residue, so the taxonomy of the completion is that of the base and iterating changes nothing structural (Theorem \ref{c:transport.rule-1-morita-trivial}), while Rule 2 moves the residue off the split classes and can therefore occur at most once in any transport history (Proposition \ref{c:transport.fusion-occurs-most-once}). A transport history takes at most one step on the clock, and then remains fixed.
\end{enumerate}

\end{proof}

\subsection*{The transport algorithm}
\phantomsection\label{c:note.transport.setting-transport-algorithm}
\begin{tabular}{c>{ \arraybackslash}p{3.5cm}>{ \arraybackslash}p{4.1cm}>{ \arraybackslash}p{4.3cm}}
\toprule
Step & Question answered & Instrument & Warrant \\
\midrule
1 & where does division fail? & the case at hand & Lemma \ref{c:transport.noinverse} \\
2 & which of the three obstructions is it? & $pq \neq 0$; $p - q \equiv 1 \bmod 4$ & Prop.\ \ref{c:transport.existence-locus-type-n}; Thms \ref{c:transport.existence-classification-type-c}, \ref{c:transport.exhaustiveness-every-zero-divisor-annihilated-idempotent} \\
3 & what is the minimal repair, and what does it adjoin? & $(+1,+1)$; one generator; nothing & Thms \ref{c:transport.minimality}, \ref{c:transport.representation-independent-minimality-completion}, \ref{c:transport.one-generator-minimal-minimal-completions-classified}, \ref{c:transport.embedding-independent-minimality-fusion}; Prop.\ \ref{c:transport.type-p-self-repairing} \\
4 & how is the value at the obstruction computed? & lift, then $\iota(s,x,t) = (2t,x,\tfrac12 s)$ & Thm \ref{c:transport.totality}, Prop.\ \ref{c:transport.swap-rescale-normal-form} \\
5 & what does descent back to the original algebra cost? & read-out $\dsc$; pinching $E$ & Lem.\ \ref{c:transport.descent-read-out-map}, Prop.\ \ref{c:transport.pinching-kernel}, Thm \ref{c:transport.door-one-way} \\
6 & does repeating the repair terminate? & the mod-8 clock & Thm \ref{c:transport.rule-1-morita-trivial}, Prop.\ \ref{c:transport.fusion-occurs-most-once} \\
\bottomrule
\end{tabular}

\emph{The transport algorithm. Each row is a question a practitioner meets in order, the instrument that answers it, and the theorem that licenses the instrument. Step 3 is the only step where the three obstruction types diverge; Steps 4 and 5 are stated for type N, the only type whose repair is delivered at the level of points rather than elements.}

\begin{remark}[A worked instance]\label{c:note.transport.worked-instance}
The procedure is run end to end, on the case that motivated the programme, in the series' gravitational paper \cite{ButlerT5}, whose chain of steps mirrors the table accompanying \ref{c:transport.transport-algorithm} row for row. It is worth carrying the outline through the paper as a running example. In it: the failure is located at the centre of the Schwarzschild solution, where the light cones fold onto a single direction (Step 1); the degenerating direction is null and the algebra $\Cl(1,3)$ is simple, with $d \equiv 6 \bmod 8$, so the diagnosis is type N and type C is unavailable (Step 2, by Theorem \ref{c:transport.existence-classification-type-c} and Example \ref{c:note.transport.cl-1-3-physical-case-n}); Rule 1 gives $\Cl(1,3) \to \Cl(2,4)$, the conformal algebra of spacetime, forced (Step 3); the transported division is total there (Step 4); and the descent is lossy in exactly the way Step 5 states, which is where that paper's physical content begins. The example is worth keeping in view for one reason in particular: the answer to the forbidden division is not a number in spacetime, and the algorithm says so in advance rather than after the fact.
\end{remark}

\begin{remark}[What the procedure does not do]\label{c:note.transport.what-procedure-does-not-do}
Three boundaries, each a theorem rather than a caution.

It never returns an inverse. For type N no element-level repair exists at all (Theorem \ref{c:transport.no-go-level-elements}); what the ascent supplies is a total operation on points together with a canonical destination for the values the base cannot hold. A reader who wants $1/v$ will not find it here, and Theorem \ref{c:transport.no-go-level-elements} is the statement that nobody will.

It assumes the form is nondegenerate. Everything above is stated for nondegenerate $(\mathbb{V},Q)$. A degenerate form is not a harder case of the same problem but a different one: its radical generates a square-zero ideal rather than a resolvable obstruction (Theorem \ref{c:transport.one-generator-minimal-minimal-completions-classified}(2) is the fibre-level instance), and completing it is the subject of the series' sheaf-theoretic paper \cite{ButlerT4}.

It is a procedure for one algebra. When division fails not at a point of an algebra but along a locus of a family, the rank of the relevant bundle cannot jump, and the correct object is a pair of sheaves rather than a single ascent: the fibrewise version of Step 3, with the completion performed over the degenerate locus and only there, is \cite[Thm.\ 4.4, Cor.\ 4.5]{ButlerT4}.
\end{remark}

\phantomsection\label{c:note.transport.setting-transport-algorithm-2}
The six steps are the whole paper in operational form: the section ``Setting, and the classical facts the proofs use'' and the section ``A taxonomy of obstructions'' supplied Steps 1 and 2; the section ``Type N: the isotropic obstruction, and Rule 1'' supplied Steps 3 to 5 for the hard type, which took half the paper; the section ``Type C: ideal fusion, and Rule 2'' and the section ``Type P: the rule of no rule'' supplied Step 3 for the other two types; and the section ``The mod-8 clock and the one-way door'' supplied Step 6.

\section{Position in the series}

\phantomsection\label{c:note.transport.setting-position-series}
This paper is T1 of the series.

\emph{This paper relies on:} the companion primer T0 \cite{ButlerT0}, and on nothing else in the series. The division of labour is strict: the language belongs to the primer (the structural lemmas of the section ``Setting, and the classical facts the proofs use'' are restated with their toolkit numbers, and proved in the primer's toolkit \cite[Section 9]{ButlerT0}, save the classification theorem, cited there to its classical sources), while every fact about the failure of division, zero divisors, pseudo-inverses and idempotents included, is stated and proved in this paper. Beyond the toolkit the primer is a convenience for readers new to geometric algebra. The paper's other input is the founding idea of the programme itself, whose informal taxonomy it rigorises.

\emph{Papers that rely on this one:} every other research paper of the series. In brief: T2 \cite{ButlerT2} takes the element-level lemma that a square-zero element is invertible in no unital extension (Lemma \ref{c:transport.noinverse}), the type-N reading of a radical direction with the taxonomy frame, and the scope of the sharpenings of the section ``Minimality against arbitrary embeddings'', and proves its own completion theorem for degenerate algebras against that frame; T3 \cite{ButlerT3} iterates Rule 1 into the tower, using the Morita theorem and the pinching; T4 \cite{ButlerT4} localises the taxonomy and Rule 1 over degenerate loci, using the splitting lemma and the null-generator clause of Theorem \ref{c:transport.one-generator-minimal-minimal-completions-classified}; T5 \cite{ButlerT5} runs the whole algorithm on the black-hole singularity, importing the no-go, totality, minimality, the type-C exclusion for $\Cl(1,3)$, the pinching, and the one-way door; and the papers after T5 \cite{ButlerT5} import these results through T3 \cite{ButlerT3}, T4 \cite{ButlerT4} and T5 \cite{ButlerT5} rather than directly, with each paper's Fact environments carrying the exact provenance. The capstone T17 \cite{ButlerT17} maps the series and records the status of every claim.

One subsection points forward without relying on anything: the sentence after Corollary \ref{c:transport.which-residues-obstruct-which-way} names the dark-matter paper T14 \cite{ButlerT14} for the residues $\{2, 6\}$. It is a pointer, and the theorems of the section ``One-dimensional extensions: where the obstruction lies, and the commuting square'' use only this paper.

\section{Machine verification}

\phantomsection\label{c:note.transport.setting-machine-verification}
Every computational identity in this paper has been checked by computer in exact rational arithmetic. The consolidated script \texttt{verify\_T01.py} \cite{ButlerVerifyT01}, deposited alongside the series' shared exact-arithmetic library \texttt{GA.py}, runs one hundred and fifteen checks and reports each; no floating point is used anywhere except the Moore--Penrose comparisons of the section ``The hyperbolic atom as a matrix algebra: the Moore--Penrose pseudoinverse'', which are float and compared at machine precision, and the script exits nonzero on any failure.

The twenty-two checks inherited from the published deposit's script cover: the hyperbolic-pair and versor identities; Lemma \ref{c:transport.splitting-along-hyperbolic-plane}, including over an extension carrying a degenerate direction, where the dimension count is the essential step; all three steps of the proof of Theorem \ref{c:transport.no-go-level-elements} in $\Cl(2,1)$ and in $\Cl(3,2)$, the latter exercising the shear for \emph{every} $u$ in a three-dimensional $U$; the scale-robustness of Remark \ref{c:note.transport.locality-robustness-obstruction}, the nilpotency data of Remark \ref{c:note.transport.obstruction-parabolic}, and the partner parametrisation \ref{c:note.transport.hyperbolic-atom} with the simple transitivity on it; both branches of Theorem \ref{c:transport.totality} and the limit paths of Lemma \ref{c:transport.density} over bases of signature $(1,1)$ and $(2,1)$; Proposition \ref{c:transport.swap-rescale-normal-form} with the fixed loci of Remark \ref{c:note.transport.where-each-stratum-comes}, and Lemma \ref{c:transport.descent-read-out-map}; and the two consistency checks of the section ``Consistency checks against known theory''.

The eight checks of the section ``One-dimensional extensions: where the obstruction lies, and the commuting square'' cover: Theorem \ref{c:transport.where-obstruction-lies-commuting-square}(a) on every even-dimensional $(p,q)$ with $2 \leq n \leq 6$, with Corollary \ref{c:transport.which-residues-obstruct-which-way}'s two residue sets read off the grid, Corollary \ref{c:transport.odd-lorentzian-dimension-exactly-one-real} in dimensions $3, 5, 7, 9, 11$, and Remark \ref{c:note.transport.three-dimensions}'s three cases; Lemma \ref{c:transport.even-subalgebra-one-dimensional-extension} on $(1,3)$ and $(2,0)$ with both signs of $\nu^2$; and Theorem \ref{c:transport.where-obstruction-lies-commuting-square}(c), (d) on the whole grid, the repairs built from the signature tuple of the extension found to obstruct, with Proposition \ref{c:transport.fusion-occurs-most-once} along the square (the centre of $\Cl(p{+}1,q{+}1)$ computed blade-wise and found trivial, so fusion occurs at most once).

The other eighty-five checks (the taxonomy, the type-N sharpenings, type C, type P, and the clock) cover: Lemmas \ref{c:transport.centre} and \ref{c:transport.sign} exhaustively over \emph{all} signatures with $n \leq 8$, blade-wise, following the proofs, with the type-C criterion where it applies; the taxonomy's examples and Lemma \ref{c:transport.regularity-localises-zero-divisors-idempotents}; the full Rule-2 package for the three bases $\Cl(1,0)$, $\Cl(0,3)$, $\Cl(2,1)$ and \emph{both} ascents of each, the $\gamma$-independence of every check being the computational form of Theorem \ref{c:transport.rule-2-ideal-fusion-sharp-form}(3); the full Rule-1 package of Theorem \ref{c:transport.rule-1-morita-trivial} and Proposition \ref{c:transport.pinching-kernel} over $\Cl(1,1)$ and, for its physical relevance, $\Cl(1,3) \to \Cl(2,4)$ in dimension $64$; the type-P orbit facts with both clauses of Remark \ref{c:note.transport.moore-penrose-boundary-idempotent-side}; Theorem \ref{c:transport.cohomological-form-no-go} instantiated in $\Cl(2,2)$, with the torsor translation of Proposition \ref{c:transport.pseudo-inverse-torsor} tested modulo the left ideal $v\,\Cl(\mathbb{V}')$ by exact membership; and Lemma \ref{c:transport.conformal-gram-matrix} over five signatures, confirmed by exact congruence diagonalisation against the model lift, with Example \ref{c:note.transport.exotic-embedding-inertness} realised concretely over the Gaussian rationals.

A second layer of machine checking, of a different kind, is deposited alongside: \texttt{T01\_kernel.lean} \cite{ButlerKernelT01}, seventy-six theorems in Lean 4 checked against Mathlib, with no \texttt{sorry} and no axioms beyond Lean's own. The two layers prove different things and neither replaces the other. The script checks named identities in named signatures, exactly and exhaustively within the range it enumerates; the Lean file proves the quantified statements, for every associative ring, every element, every transport history, and every $(p,q)$. Where both apply, the second subsumes the first: Lemma \ref{c:transport.sign}'s two expressions for $\omega^2$ are checked blade-wise for $n \leq 8$ by the script and proved for all $p,q$ in Lean, the bridge between them being the polynomial identity $n(n-1) + 2q = (p-q)(p-q-1) + 4pq$; and Proposition \ref{c:transport.fusion-occurs-most-once} is confirmed on the square by the script and proved for transport histories of every length in Lean. The Lean file is not a formalisation of this paper: it contains no Clifford algebras, no quadratic forms and no signatures, and it takes as hypotheses what this paper takes from the literature, namely semisimplicity, the injectivity of a unital homomorphism out of a simple ring, and Theorem \ref{c:transport.rule-1-morita-trivial}. Its header carries a concordance from each statement label used here to the theorems that carry it, and a list of what it does not cover.

\section{Results used from elsewhere}
The proofs above use the following results of the series, stated briefly; each is proved in the paper named.

\begin{fact}[The centre]\label{c:primer.centre}
For $n \geq 1$ the centre of $\Cl(p,q)$ is $\R \cdot 1$ when $n$ is even and $\R \cdot 1 \oplus \R \cdot \omega$ when $n$ is odd.
\emph{Source:} T0, Lem 9.4 (\href{https://doi.org/10.5281/zenodo.22714501}{doi:10.5281/zenodo.22714501}).
\end{fact}

\begin{fact}[Classification of the real Clifford algebras]\label{c:primer.classification-real-clifford-algebras}
$\Cl(p,q)$ is a matrix algebra over $\R$, $\C$ or $\HH$, or a sum of two such, according to $p - q \bmod 8$, and it is simple except when $p - q \equiv 1 \pmod 4$.
\emph{Source:} T0, Thm 9.3 (\href{https://doi.org/10.5281/zenodo.22714501}{doi:10.5281/zenodo.22714501}).
\end{fact}

\begin{fact}[The commutation rule]\label{c:primer.commutation-rule}
If $x$ lies in the subalgebra generated by vectors orthogonal to $v$, then $v x = \hat{x} v$ and $x v = v \hat{x}$.
\emph{Source:} T0, Lem 9.1 (\href{https://doi.org/10.5281/zenodo.22714501}{doi:10.5281/zenodo.22714501}).
\end{fact}

\begin{fact}[The sign of $\omega^2$]\label{c:primer.sign}
The pseudoscalar of $\Cl(p,q)$ satisfies $\omega^{2} = (-1)^{(p-q)(p-q-1)/2}$, so $\omega^{2} = +1$ exactly when $p - q \equiv 0$ or $1 \pmod 4$.
\emph{Source:} T0, Lem 9.5 (\href{https://doi.org/10.5281/zenodo.22714501}{doi:10.5281/zenodo.22714501}).
\end{fact}

\begin{fact}[Splitting along a hyperbolic plane]\label{c:primer.splitting-along-hyperbolic-plane}
Along a hyperbolic plane spanned by a null pair $(v, w_0)$, every element of $\Cl(H_0 \perp U')$ expands uniquely as $\alpha + v\beta + w_0\gamma + v w_0\delta$ with coefficients in $\Cl(U')$, and left multiplication by $v$ is injective on $\Cl(U')$.
\emph{Source:} T0, Lem 9.6 (\href{https://doi.org/10.5281/zenodo.22714501}{doi:10.5281/zenodo.22714501}).
\end{fact}

\begin{fact}[The blade basis, degenerate forms included]\label{c:std-algebra.blade-basis}
For any basis of $V$ the $2^{n}$ increasing products $e_I$ form a basis of $\Cl(V, Q)$, degenerate forms included; so $V \hookrightarrow \Cl(V, Q)$.
\emph{Source:} the algebra core of the Transport Lab, proved there with an exact check; standard, see C. Chevalley, The Algebraic Theory of Spinors, Columbia University Press (1954), Ch. II (the Clifford algebra of a quadratic space).
\end{fact}

\begin{fact}[Complexification, and the odd-dimensional split]\label{c:std-algebra.complexification}
$\Cl(V) \otimes \C \cong \Cl_n(\C)$, a complex matrix algebra (doubled for odd $n$); for odd $n$, $\Cl \cong \Cl^{+} \oplus \Cl^{+}$ if $\omega^{2} = 1$ and $\Cl \cong \Cl^{+} \otimes \C$ if $\omega^{2} = -1$.
\emph{Source:} the algebra core of the Transport Lab, proved there with an exact check; standard, see P. Lounesto, Clifford Algebras and Spinors, 2nd ed., London Mathematical Society Lecture Note Series 286, Cambridge University Press (2001), Ch. 16 (matrix representations and periodicity of 8).
\end{fact}

\begin{fact}[The Clifford algebra exists and is unique]\label{c:std-algebra.existence-uniqueness}
$\Cl(V, Q) = T(V)/\langle v \otimes v - Q(v) \rangle$ exists and is unique up to unique isomorphism; the wedge-plus-contraction operators on $\Lambda V$ give a concrete representation.
\emph{Source:} the algebra core of the Transport Lab, proved there with an exact check; standard, see C. Chevalley, The Algebraic Theory of Spinors, Columbia University Press (1954), Ch. II (the Clifford algebra of a quadratic space).
\end{fact}

\begin{fact}[The three involutions on each grade]\label{c:std-algebra.involution-properties}
On $\Cl^{k}$: $\hat{x} = (-1)^{k}x$, $\tilde{x} = (-1)^{k(k-1)/2}x$, $\bar{x} = (-1)^{k(k+1)/2}x$; the first is an automorphism, the others anti-automorphisms.
\emph{Source:} the algebra core of the Transport Lab, proved there with an exact check; standard, see P. Lounesto, Clifford Algebras and Spinors, 2nd ed., London Mathematical Society Lecture Note Series 286, Cambridge University Press (2001), Ch. 14 (definitions of the Clifford algebra).
\end{fact}

\begin{fact}[An algebra and its matrix amplification]\label{c:std-algebra.morita-amplification}
$\Zc(\Mat_m(A)) = \Zc(A)$, the ideals of $\Mat_m(A)$ are $\Mat_m(I)$, and $A$ and $\Mat_m(A)$ have equivalent module categories.
\emph{Source:} the algebra core of the Transport Lab, proved there with an exact check; standard, see T. Y. Lam, Lectures on Modules and Rings, Graduate Texts in Mathematics 189, Springer, New York (1999), matrix rings, idempotents and Morita theory.
\end{fact}

\begin{fact}[Orthogonal vectors anticommute]\label{c:std-algebra.orthogonal-anticommute}
$uv = -vu$ iff $B(u,v) = 0$; in an orthogonal basis $e_I e_J = \pm \prod_{I \cap J} Q(e_i)\; e_{I \triangle J}$ with the sorting sign.
\emph{Source:} the algebra core of the Transport Lab, proved there with an exact check; standard, see C. Doran and A. Lasenby, Geometric Algebra for Physicists, Cambridge University Press (2003), Ch. 4 (foundations of geometric algebra).
\end{fact}

\begin{fact}[Reflection in a non-null vector]\label{c:std-algebra.reflection}
For $Q(u) \ne 0$, $x \mapsto -uxu^{-1} = x - 2\frac{B(x,u)}{Q(u)}u$ is the reflection in $u^{\perp}$: an involutive isometry of determinant $-1$ fixing $\operatorname{rad} V$.
\emph{Source:} the algebra core of the Transport Lab, proved there with an exact check; standard, see C. Doran and A. Lasenby, Geometric Algebra for Physicists, Cambridge University Press (2003), Ch. 4 (foundations of geometric algebra: rotors, versors).
\end{fact}

\begin{fact}[Orthogonal bases and the law of inertia]\label{c:std-algebra.sylvester-inertia}
Every real quadratic space has an orthogonal basis squaring to $1, -1, 0$, and the counts $(p, q, r)$ are invariants that classify it up to isometry.
\emph{Source:} T. Y. Lam, Introduction to Quadratic Forms over Fields, AMS (2005)., the diagonalisation of forms and Sylvester's law of inertia (real case).
\end{fact}

\begin{fact}[The product of two vectors]\label{c:std-algebra.vector-squares}
$v^{2} = Q(v)$ and $uv + vu = 2B(u, v)$, so $uv = B(u, v) + \tfrac12(uv - vu)$.
\emph{Source:} the algebra core of the Transport Lab, proved there with an exact check; standard, see C. Doran and A. Lasenby, Geometric Algebra for Physicists, Cambridge University Press (2003), Ch. 4 (foundations of geometric algebra).
\end{fact}

\begin{fact}[Hyperbolic pairs and the Witt decomposition]\label{c:std-algebra.witt-decomposition}
A null vector in a nondegenerate plane has a unique hyperbolic partner, and $V = H_1 \perp \dots \perp H_{\min(p,q)} \perp A \perp \operatorname{rad} V$ with $A$ definite; maximal totally null subspaces have dimension $\min(p,q) + r$.
\emph{Source:} the algebra core of the Transport Lab, proved there with an exact check; standard, see T. Y. Lam, Introduction to Quadratic Forms over Fields, Graduate Studies in Mathematics 67, American Mathematical Society (2005), Ch. I (hyperbolic planes, the Witt decomposition, Witt's theorems).
\end{fact}

\begin{fact}[The even subalgebra and the Z2 grading]\label{c:std-algebra.z2-grading}
$\Cl = \Cl^{+} \oplus \Cl^{-}$ is a $\mathbb{Z}_2$-grading, $\dim \Cl^{+} = 2^{n-1}$, and $\Cl^{j}\Cl^{k}$ has grades $|j-k|, |j-k|+2, \dots, j+k$.
\emph{Source:} the algebra core of the Transport Lab, proved there with an exact check; standard, see P. Lounesto, Clifford Algebras and Spinors, 2nd ed., London Mathematical Society Lecture Note Series 286, Cambridge University Press (2001), Ch. 14 (definitions of the Clifford algebra).
\end{fact}

\begin{refutation}[The two one-generator ascents differ as algebras]\label{c:refuted.the-two-one-generator-ascents-differ-as-algebras}
Refuted: the two one-generator ascents are always isomorphic as algebras, the sign of the adjoined generator being a geometric bit the algebra forgets.
\end{refutation}

\section*{Changes since the deposited version}
\addcontentsline{toc}{section}{Changes since the deposited version}
Since the deposited version 3.0: 2 results changed, 14 reworded (wording, notation or citations only) and 17 cited results added. The comparison is with the source of version 3.0 on Zenodo, \href{https://doi.org/10.5281/zenodo.21743863}{doi:10.5281/zenodo.21743863}. A result counts as revised when the words or mathematics of its statement or proof differ from that source; a citation added in a parenthesis does not count. A result the paper cites from elsewhere is compared with the deposit of the paper that proves it. Each revision is classed as corrected, changed or reworded; the classes were assigned by a review of each difference and the store's record of why it was made, and a revision not yet reviewed is counted as changed.
\subsection*{Changed}
The content changed by a later decision, not because the deposited version was wrong.
\begin{itemize}\setlength{\itemsep}{2pt}
\item Theorem~\ref{c:transport.where-obstruction-lies-commuting-square} (\emph{Where the obstruction lies, and the commuting square}), Thm 5.14 in the deposit: statement. The ladder pair is now written in the two adjoined generators u (square +1) and w (square -1), with their squares stated.
\item Corollary~\ref{c:transport.which-residues-obstruct-which-way} (\emph{Which residues obstruct which way}), Cor 5.15 in the deposit: statement and proof. New clause: at residue 2, in T14's causal convention, the extension is the mirror's second time, with unobstructed Cl(4,1).
\end{itemize}
\subsection*{Reworded}
Wording, notation, terminology or citations only; the mathematics and the claim are as deposited.
\begin{itemize}\setlength{\itemsep}{2pt}
\item Lemma~\ref{c:transport.noinverse} (\emph{A null vector has no inverse}), Lem 2.1 in the deposit: proof. Notation and citation only: v\^{}2 written as Q(v) and the vector-squares fact cited.
\item Lemma~\ref{c:transport.splitting-along-hyperbolic-plane} (\emph{Splitting along a hyperbolic plane}), Lem 2.4 in the deposit: proof. Attribution only: the proof now points to the primer, where the lemma is proved.
\item Theorem~\ref{c:transport.classification-real-clifford-algebras} (\emph{Classification of the real Clifford algebras}), Thm 2.5 in the deposit: proof. Attribution only: the proof now points to the primer's classification theorem.
\item Lemma~\ref{c:transport.centre} (\emph{The centre}), Lem 2.6 in the deposit: proof. Attribution only: the proof now points to the primer, where the centre is computed.
\item Lemma~\ref{c:transport.sign} (\emph{The sign of $\omega^2$}), Lem 2.7 in the deposit: proof. Attribution only: the proof now points to the primer, where the sign of the pseudoscalar square is computed.
\item Lemma~\ref{c:transport.group-null-rotations} (\emph{The group of null rotations}), Lem 4.6 in the deposit: statement. The setting (null v, hyperbolic partner w, complement U) is stated explicitly and supporting facts cited; same result.
\item Theorem~\ref{c:transport.one-generator-minimal-minimal-completions-classified} (\emph{One generator is minimal; the minimal completions classified}), Thm 5.7 in the deposit: proof. Citations only: numbered references replaced by named results, and Sylvester's law of inertia cited.
\item Theorem~\ref{c:transport.embedding-independent-minimality-fusion} (\emph{Embedding-independent minimality of the fusion}), Thm 5.11 in the deposit: proof. Citations only: references updated and the classification cited for the module splitting.
\item Theorem~\ref{c:transport.classification-idempotent-orbits} (\emph{Classification of idempotent orbits}), Thm 6.4 in the deposit: proof. Citations only: references updated and the classification cited where the matrix form is used.
\item Corollary~\ref{c:transport.no-fourth-rule} (\emph{No fourth rule}), Cor 6.7 in the deposit: proof. Presentation only: a note that the corollary collects the cited theorems, with references updated.
\item Theorem~\ref{c:transport.rule-1-morita-trivial} (\emph{Rule 1 is Morita-trivial}), Thm 7.1 in the deposit: proof. Notation and citation only: f.f written as Q(f), with standard facts cited.
\item Proposition~\ref{c:transport.fusion-occurs-most-once} (\emph{Fusion occurs at most once}), Prop 7.4 in the deposit: proof. Citations only: references updated and the classification cited for the residue of Cl(1,3).
\item Theorem~\ref{c:transport.transport-algorithm} (\emph{The transport algorithm}), Sec 8.1 in the deposit: proof. Citations only: numbered theorem references replaced by named results.
\end{itemize}
\noindent In addition, 17 results from elsewhere in the series or the algebra core are stated here that the deposited paper cited or used without stating; 4 results gained citations only. These are not marked.
\begin{thebibliography}{99}
\bibitem{HestenesSobczyk1984} D.~Hestenes and G.~Sobczyk,
\emph{Clifford Algebra to Geometric Calculus},
Reidel, Dordrecht (1984).
\bibitem{DoranLasenby2003} C.~Doran and A.~Lasenby,
\emph{Geometric Algebra for Physicists},
Cambridge University Press (2003).
\bibitem{ButlerT3} L.~Butler,
\emph{The Tower and the CAR Limit: Iterated Transport, the Frozen Clock, and Second Quantisation},
Transport Papers T3 (2026), doi:\href{https://doi.org/10.5281/zenodo.22721556}{10.5281/zenodo.22721556}.
\bibitem{ABS1964} M.~F.~Atiyah, R.~Bott and A.~Shapiro,
Clifford modules,
\emph{Topology} \textbf{3} (1964), suppl.~1, 3--38.
\bibitem{Porteous1995} I.~R.~Porteous,
\emph{Clifford Algebras and the Classical Groups},
Cambridge University Press (1995).
\bibitem{Lounesto2001} P.~Lounesto,
\emph{Clifford Algebras and Spinors}, 2nd ed.,
London Mathematical Society Lecture Note Series 286, Cambridge University Press (2001).
\bibitem{Lam2005} T.~Y.~Lam,
\emph{Introduction to Quadratic Forms over Fields},
Graduate Studies in Mathematics 67, American Mathematical Society (2005).
\bibitem{Chevalley1997} C.~Chevalley,
\emph{The Algebraic Theory of Spinors and Clifford Algebras},
Collected Works, Vol.~2, Springer, Berlin (1997).
\bibitem{ButlerT0} L.~Butler,
\emph{An Introduction to Geometric Algebra: A Primer for the Transport Papers},
Transport Papers T0 (2026).
\bibitem{vonNeumann1936} J.~von Neumann,
On regular rings,
\emph{Proc.\ Natl.\ Acad.\ Sci.\ USA} \textbf{22} (1936), 707--713.
\bibitem{Ablamowicz1998} R.~Ab{\l}amowicz,
Spinor representations of Clifford algebras: a symbolic approach,
\emph{Computer Physics Communications} \textbf{115} (1998), 510--535.
\bibitem{Hestenes2001} D.~Hestenes,
Old wine in new bottles: a new algebraic framework for computational geometry,
in \emph{Geometric Algebra with Applications in Science and Engineering},
E.~Bayro-Corrochano and G.~Sobczyk (eds.), Birkh\"auser, Boston (2001).
\bibitem{HLR2001} H.~Li, D.~Hestenes and A.~Rockwood,
Generalized homogeneous coordinates for computational geometry,
in \emph{Geometric Computing with Clifford Algebras},
G.~Sommer (ed.), Springer (2001), 27--59.
\bibitem{DFM2007} L.~Dorst, D.~Fontijne and S.~Mann,
\emph{Geometric Algebra for Computer Science: An Object-Oriented Approach to Geometry},
Morgan Kaufmann (2007).
\bibitem{Angles2008} P.~Angl\`es,
\emph{Conformal Groups in Geometry and Spin Structures},
Progress in Mathematical Physics 50, Birkh\"auser (2008).
\bibitem{Zeeman1964} E.~C.~Zeeman,
Causality implies the Lorentz group,
\emph{J. Math. Phys.} \textbf{5} (1964) 490--493.
\bibitem{LamFC2001} T.~Y.~Lam,
\emph{A First Course in Noncommutative Rings}, 2nd ed.,
Graduate Texts in Mathematics 131, Springer (2001).
\bibitem{ButlerT4} L.~Butler,
\emph{Clifford Sheaves over Degenerate Metrics: Nilpotent Ideals, One-Sided Completion, and the Exceptional Algebra},
Transport Papers T4 (2026), doi:\href{https://doi.org/10.5281/zenodo.22721669}{10.5281/zenodo.22721669}.
\bibitem{ButlerT14} L.~Butler,
\emph{Dark Matter and the Arrow of Time: The Five Dark Towers},
Transport Papers T14 (2026).
\bibitem{ButlerT5} L.~Butler,
\emph{Division by Zero at the Gravitational Singularity: The Transport Principle Applied to the Black Hole, and the Information Paradox as a Category Error},
Transport Papers T5 (2026), doi:\href{https://doi.org/10.5281/zenodo.22721790}{10.5281/zenodo.22721790}.
\bibitem{ButlerT2} L.~Butler,
\emph{Projective Completes to Conformal: The Minimal Nondegenerate Completion of a Degenerate Clifford Algebra, and Where PGA's Missing Division Goes},
Transport Papers T2 (2026), doi:\href{https://doi.org/10.5281/zenodo.22714886}{10.5281/zenodo.22714886}.
\bibitem{ButlerT17} L.~Butler,
\emph{The Transport Programme: Method, Status, and the Falsifiable Ledger},
Transport Papers T17 (2026).
\bibitem{ButlerVerifyT01} L.~Butler,
\emph{verify\_T01.py: exact-arithmetic verification of the computational identities of ``The Transport Principle''},
Python script, deposited with this paper alongside the shared library GA.py (2026).
\bibitem{ButlerKernelT01} L.~Butler,
\emph{T01\_kernel.lean: a Lean 4 formalisation of the algebraic kernel of ``The Transport Principle''},
Lean 4 source checked against Mathlib, deposited with this paper (2026).
\end{thebibliography}
\end{document}